% !TEX program = lualatex
% ---------------------------------------------------------------------------------------------
%  Readback note 16: Partner data, vanishing on standard data, and the exchange identity at the top of the tower
%
%  Build:  python3 tools/rb.py note <slug>
%
%  One unit of a Lean artifact, rewritten as ordinary mathematics. The prose is hand-written. Lean names,
%  line ranges and counts come from data.tex; quoted Lean statements are copied from the artifact by the build.
%
%  Rules of the text: standard terminology; every symbol is explained where it first appears and listed again
%  in the notation at the end; no letter has two meanings; a proof is complete or says what it takes for granted;
%  Lean names appear only in the LEAN lines and in the concordance.
%
%  The signature rule (the build enforces it): a function, relation or predicate that the artifact defines is
%  written only through its notation macro, which the build makes from the spec (`notation`) with one slot for
%  each explicit argument of the Lean definition. No argument is fixed and dropped, also not "for brevity" and
%  not for one section. Every argument and every hypothesis of a stated declaration stands inside its
%  statement, or in a Setting that the statement cites by number; only implicit and instance arguments (the
%  ambient objects) may stand in the conventions.
% ---------------------------------------------------------------------------------------------
\documentclass[11pt,a4paper]{article}
\ifdefined\pdfvariable\pdfvariable suppressoptionalinfo \numexpr 1+2+4+8\relax\fi
\usepackage[a4paper,hmargin=35mm,top=22mm,bottom=27mm,footskip=12mm,nohead]{geometry}

\usepackage{fontspec}
\directlua{
  luaotfload.add_fallback("gfallback", {"STIXTwoMath-Regular.otf:mode=node;", "DejaVuSans.ttf:mode=node;"})
  luaotfload.add_fallback("gmonofallback", {"DejaVuSansMono.ttf:mode=node;", "STIXTwoMath-Regular.otf:mode=node;"})
}
\setmainfont{NewCM10}[Extension=.otf,UprightFont=*-Book,ItalicFont=*-BookItalic,
  BoldFont=*-Bold,BoldItalicFont=*-BoldItalic,SmallCapsFont=*-Book,SmallCapsFeatures={Letters=SmallCaps},
  RawFeature={fallback=gfallback}]
\setsansfont{NewCMSans10}[Extension=.otf,UprightFont=*-Book,ItalicFont=*-BookOblique,
  BoldFont=*-Bold,BoldItalicFont=*-BoldOblique,RawFeature={fallback=gfallback}]
\setmonofont{JetBrainsMono}[Extension=.otf,UprightFont=*-Regular,ItalicFont=*-Italic,
  BoldFont=*-SemiBold,BoldItalicFont=*-SemiBoldItalic,Scale=MatchLowercase,
  RawFeature={-calt;-liga;fallback=gmonofallback}]
\usepackage{amsmath}
\usepackage{amsthm}
\usepackage[math-style=TeX]{unicode-math}
\setmathfont{NewCMMath-Book.otf}
\emergencystretch=2.5em\relax

\usepackage{xcolor}
\usepackage{tikz}
\usetikzlibrary{calc,arrows.meta,positioning,backgrounds}
\usepackage[protrusion=true,expansion=true]{microtype}
\usepackage{enumitem}
\usepackage{needspace}
\usepackage{booktabs}
\usepackage{array}
\usepackage{longtable}
\usepackage{fancyhdr}
\usepackage{lastpage}
\usepackage{multicol}
\usepackage[hidelinks,pdfusetitle]{hyperref}

\definecolor{ink}{HTML}{0F172A}
\definecolor{inksoft}{HTML}{475569}
\definecolor{inkfaint}{HTML}{94A3B8}
\definecolor{rulegrey}{HTML}{CBD5E1}
\definecolor{arith}{HTML}{C2410C}
\definecolor{est}{HTML}{7C3AED}
\definecolor{ana}{HTML}{0891B2}
\definecolor{ok}{HTML}{047857}
\definecolor{leancol}{HTML}{3730A3}

\title{Partner data, vanishing on standard data, and the exchange identity at the top of the tower}
\author{grh-compact}
\date{}

% ------------------------------------------------------------------------------------- macros
\newcommand{\leanfont}{\ttfamily\fontsize{8.6}{10.5}\selectfont\color{leancol}}
\newcommand{\lean}[1]{\csname LN@#1\endcsname}
\newcommand{\fillin}[1]{{\color{arith}\itshape[#1]}}  % a place of the template that is not written yet; the build reports every one that is left
\newcommand{\leanquote}[1]{\csname LQ@#1\endcsname}
% generated by rb.py note; do not edit
\newcommand{\ConcDecls}{47}
\newcommand{\Pin}{72fead10}
\newcommand{\PinText}{that of commit {\leanfont 72fead10}}
\newcommand{\ArtifactName}{grh-compact}
\expandafter\newcommand\csname LN@existscommoncosetreps\endcsname{\href{https://github.com/lsernau/grh-lean/blob/72fead102f663d5c7a407964a13dc0cf9d977207/grh-compact/src/New/Exchange/BlockPartners.lean\#L20-L29}{{\leanfont exists\_common\_coset\_reps}}}
\expandafter\newcommand\csname LN@localSection\endcsname{\href{https://github.com/lsernau/grh-lean/blob/72fead102f663d5c7a407964a13dc0cf9d977207/grh-compact/src/New/Exchange/BlockPartners.lean\#L33-L34}{{\leanfont localSection}}}
\expandafter\newcommand\csname LN@matchingPairs\endcsname{\href{https://github.com/lsernau/grh-lean/blob/72fead102f663d5c7a407964a13dc0cf9d977207/grh-compact/src/New/Exchange/BlockPartners.lean\#L39-L41}{{\leanfont matchingPairs}}}
\expandafter\newcommand\csname LN@mulfirstCols\endcsname{\href{https://github.com/lsernau/grh-lean/blob/72fead102f663d5c7a407964a13dc0cf9d977207/grh-compact/src/New/Exchange/BlockPartners.lean\#L44-L44}{{\leanfont mul\_firstCols}}}
\expandafter\newcommand\csname LN@translatemem\endcsname{\href{https://github.com/lsernau/grh-lean/blob/72fead102f663d5c7a407964a13dc0cf9d977207/grh-compact/src/New/Exchange/BlockPartners.lean\#L46-L51}{{\leanfont matchingPairs.\allowbreak{}translate\_mem}}}
\expandafter\newcommand\csname LN@localSectionofprimitive\endcsname{\href{https://github.com/lsernau/grh-lean/blob/72fead102f663d5c7a407964a13dc0cf9d977207/grh-compact/src/New/Exchange/BlockPartners.lean\#L55-L58}{{\leanfont localSection\_of\_primitive}}}
\expandafter\newcommand\csname LN@isClopenprimitive\endcsname{\href{https://github.com/lsernau/grh-lean/blob/72fead102f663d5c7a407964a13dc0cf9d977207/grh-compact/src/New/Exchange/BlockPartners.lean\#L60-L65}{{\leanfont TameLocalChar.\allowbreak{}isClopen\_primitive}}}
\expandafter\newcommand\csname LN@existslinearfactor\endcsname{\href{https://github.com/lsernau/grh-lean/blob/72fead102f663d5c7a407964a13dc0cf9d977207/grh-compact/src/New/Exchange/BlockPartners.lean\#L70-L89}{{\leanfont matchingPairs.\allowbreak{}exists\_linear\_factor}}}
\expandafter\newcommand\csname LN@partnerblock\endcsname{\href{https://github.com/lsernau/grh-lean/blob/72fead102f663d5c7a407964a13dc0cf9d977207/grh-compact/src/New/Exchange/BlockPartners.lean\#L94-L108}{{\leanfont partner\_block}}}
\expandafter\newcommand\csname LN@adjointthetaeqformLift\endcsname{\href{https://github.com/lsernau/grh-lean/blob/72fead102f663d5c7a407964a13dc0cf9d977207/grh-compact/src/New/Exchange/CharacterLiftNorm.lean\#L36-L44}{{\leanfont AuxPair.\allowbreak{}adjoint\_theta\_eq\_formLift}}}
\expandafter\newcommand\csname LN@liftauxDatumeqliftmodDatum\endcsname{\href{https://github.com/lsernau/grh-lean/blob/72fead102f663d5c7a407964a13dc0cf9d977207/grh-compact/src/New/Exchange/CharacterLiftNorm.lean\#L49-L56}{{\leanfont AuxPair.\allowbreak{}lift\_auxDatum\_eq\_lift\_modDatum}}}
\expandafter\newcommand\csname LN@levelIndicatorsupportandnormSq\endcsname{\href{https://github.com/lsernau/grh-lean/blob/72fead102f663d5c7a407964a13dc0cf9d977207/grh-compact/src/New/Exchange/CharacterLiftNorm.lean\#L62-L75}{{\leanfont AuxPair.\allowbreak{}levelIndicator\_support\_and\_normSq}}}
\expandafter\newcommand\csname LN@gramFactor\endcsname{\href{https://github.com/lsernau/grh-lean/blob/72fead102f663d5c7a407964a13dc0cf9d977207/grh-compact/src/New/Exchange/CharacterLiftNorm.lean\#L79-L80}{{\leanfont AuxPair.\allowbreak{}gramFactor}}}
\expandafter\newcommand\csname LN@rallisTerm\endcsname{\href{https://github.com/lsernau/grh-lean/blob/72fead102f663d5c7a407964a13dc0cf9d977207/grh-compact/src/New/Exchange/CharacterLiftNorm.lean\#L84-L86}{{\leanfont Data.\allowbreak{}rallisTerm}}}
\expandafter\newcommand\csname LN@existsisSphericalWithiota\endcsname{\href{https://github.com/lsernau/grh-lean/blob/72fead102f663d5c7a407964a13dc0cf9d977207/grh-compact/src/New/Exchange/CharacterLiftNorm.lean\#L91-L101}{{\leanfont Data.\allowbreak{}exists\_isSphericalWith\_iota}}}
\expandafter\newcommand\csname LN@innerauxLift\endcsname{\href{https://github.com/lsernau/grh-lean/blob/72fead102f663d5c7a407964a13dc0cf9d977207/grh-compact/src/New/Exchange/CharacterLiftNorm.lean\#L105-L106}{{\leanfont Data.\allowbreak{}inner\_auxLift}}}
\expandafter\newcommand\csname LN@auxLiftinjectiveandhasSuminner\endcsname{\href{https://github.com/lsernau/grh-lean/blob/72fead102f663d5c7a407964a13dc0cf9d977207/grh-compact/src/New/Exchange/CharacterLiftNorm.lean\#L113-L132}{{\leanfont Data.\allowbreak{}auxLift\_injective\_and\_hasSum\_inner}}}
\expandafter\newcommand\csname LQ@auxLiftinjectiveandhasSuminner\endcsname{\leanquoteline{0}{lemma Data.auxLift\_injective\_and\_hasSum\_inner : ∃ A ≠ 0, ∀ ϑ, Injective (D.auxLift ϑ j) ∧ ∀ x y,}\leanquoteline{2}{HasSum (A * (D.auxPair ϑ j).gramFactor * D.rallisTerm ϑ j \_ x y ·)}\leanquoteline{2}{(inner ℂ (D.auxLift ϑ j x) (D.auxLift ϑ j y))}}
\expandafter\newcommand\csname LN@liftsauxLift\endcsname{\href{https://github.com/lsernau/grh-lean/blob/72fead102f663d5c7a407964a13dc0cf9d977207/grh-compact/src/New/Exchange/CharacterLiftNorm.lean\#L137-L143}{{\leanfont Data.\allowbreak{}lifts\_auxLift}}}
\expandafter\newcommand\csname LN@auxLiftinjective\endcsname{\href{https://github.com/lsernau/grh-lean/blob/72fead102f663d5c7a407964a13dc0cf9d977207/grh-compact/src/New/Exchange/CharacterLiftNorm.lean\#L147-L148}{{\leanfont Data.\allowbreak{}auxLift\_injective}}}
\expandafter\newcommand\csname LN@auxLiftprincipalRep\endcsname{\href{https://github.com/lsernau/grh-lean/blob/72fead102f663d5c7a407964a13dc0cf9d977207/grh-compact/src/New/Exchange/CharacterLiftNorm.lean\#L152-L158}{{\leanfont Data.\allowbreak{}auxLift\_principalRep}}}
\expandafter\newcommand\csname LN@auxPairliftstrivialLift\endcsname{\href{https://github.com/lsernau/grh-lean/blob/72fead102f663d5c7a407964a13dc0cf9d977207/grh-compact/src/New/Exchange/CharacterLiftNorm.lean\#L174-L184}{{\leanfont LValueData.\allowbreak{}auxPair\_lifts\_trivialLift}}}
\expandafter\newcommand\csname LN@innertrivialLiftselfeqL\endcsname{\href{https://github.com/lsernau/grh-lean/blob/72fead102f663d5c7a407964a13dc0cf9d977207/grh-compact/src/New/Exchange/CharacterLiftNorm.lean\#L191-L252}{{\leanfont LValueData.\allowbreak{}inner\_trivialLift\_self\_eq\_L}}}
\expandafter\newcommand\csname LN@posCoordMatrix\endcsname{\href{https://github.com/lsernau/grh-lean/blob/72fead102f663d5c7a407964a13dc0cf9d977207/grh-compact/src/New/Exchange/DeterminantVectors.lean\#L27-L28}{{\leanfont posCoordMatrix}}}
\expandafter\newcommand\csname LN@detVectorFn\endcsname{\href{https://github.com/lsernau/grh-lean/blob/72fead102f663d5c7a407964a13dc0cf9d977207/grh-compact/src/New/Exchange/DeterminantVectors.lean\#L32-L33}{{\leanfont detVectorFn}}}
\expandafter\newcommand\csname LN@detVectors\endcsname{\href{https://github.com/lsernau/grh-lean/blob/72fead102f663d5c7a407964a13dc0cf9d977207/grh-compact/src/New/Exchange/DeterminantVectors.lean\#L41-L42}{{\leanfont detVectors}}}
\expandafter\newcommand\csname LN@frameAtrealunitary\endcsname{\href{https://github.com/lsernau/grh-lean/blob/72fead102f663d5c7a407964a13dc0cf9d977207/grh-compact/src/New/Exchange/DeterminantVectors.lean\#L45-L50}{{\leanfont frameAt\_real\_unitary}}}
\expandafter\newcommand\csname LN@posProjeq\endcsname{\href{https://github.com/lsernau/grh-lean/blob/72fead102f663d5c7a407964a13dc0cf9d977207/grh-compact/src/New/Exchange/DeterminantVectors.lean\#L53-L54}{{\leanfont posProj\_eq}}}
\expandafter\newcommand\csname LN@existssliceeq\endcsname{\href{https://github.com/lsernau/grh-lean/blob/72fead102f663d5c7a407964a13dc0cf9d977207/grh-compact/src/New/Exchange/DeterminantVectors.lean\#L59-L83}{{\leanfont PositivePart.\allowbreak{}exists\_slice\_eq}}}
\expandafter\newcommand\csname LN@surjOnconeintegraldetVectors\endcsname{\href{https://github.com/lsernau/grh-lean/blob/72fead102f663d5c7a407964a13dc0cf9d977207/grh-compact/src/New/Exchange/DeterminantVectors.lean\#L88-L144}{{\leanfont surjOn\_cone\_integral\_detVectors}}}
\expandafter\newcommand\csname LN@thetaSmall\endcsname{\href{https://github.com/lsernau/grh-lean/blob/72fead102f663d5c7a407964a13dc0cf9d977207/grh-compact/src/New/Exchange/ExchangeDifference.lean\#L17-L18}{{\leanfont thetaSmall}}}
\expandafter\newcommand\csname LN@existsisEquivariantLift\endcsname{\href{https://github.com/lsernau/grh-lean/blob/72fead102f663d5c7a407964a13dc0cf9d977207/grh-compact/src/New/Exchange/ExchangeDifference.lean\#L26-L60}{{\leanfont exists\_isEquivariantLift}}}
\expandafter\newcommand\csname LN@thetaOnRealData\endcsname{\href{https://github.com/lsernau/grh-lean/blob/72fead102f663d5c7a407964a13dc0cf9d977207/grh-compact/src/New/Exchange/LowRankIdentity.lean\#L21-L22}{{\leanfont thetaOnRealData}}}
\expandafter\newcommand\csname LN@thetaTermstdDatumspec\endcsname{\href{https://github.com/lsernau/grh-lean/blob/72fead102f663d5c7a407964a13dc0cf9d977207/grh-compact/src/New/Exchange/LowRankIdentity.lean\#L27-L61}{{\leanfont thetaTerm\_stdDatum\_spec}}}
\expandafter\newcommand\csname LN@thetaOnRealDataboxProdeqmul\endcsname{\href{https://github.com/lsernau/grh-lean/blob/72fead102f663d5c7a407964a13dc0cf9d977207/grh-compact/src/New/Exchange/LowRankIdentity.lean\#L65-L71}{{\leanfont thetaOnRealData\_boxProd\_eq\_mul}}}
\expandafter\newcommand\csname LN@thetaLargeeqapplyoneofevenrank\endcsname{\href{https://github.com/lsernau/grh-lean/blob/72fead102f663d5c7a407964a13dc0cf9d977207/grh-compact/src/New/Exchange/LowRankIdentity.lean\#L75-L86}{{\leanfont thetaLarge\_eq\_apply\_one\_of\_even\_rank}}}
\expandafter\newcommand\csname LN@equalRankThetaRatio\endcsname{\href{https://github.com/lsernau/grh-lean/blob/72fead102f663d5c7a407964a13dc0cf9d977207/grh-compact/src/New/Exchange/LowRankIdentity.lean\#L90-L91}{{\leanfont equalRankThetaRatio}}}
\expandafter\newcommand\csname LN@StdExchange\endcsname{\href{https://github.com/lsernau/grh-lean/blob/72fead102f663d5c7a407964a13dc0cf9d977207/grh-compact/src/New/Exchange/LowRankIdentity.lean\#L96-L99}{{\leanfont StdExchange}}}
\expandafter\newcommand\csname LQ@StdExchange\endcsname{\leanquoteline{0}{def StdExchange \{t' r\} (X : Space E m r) (V : Space E n s) (V' : Space E b t') μ σ v y (ℓ :}\leanquoteline{2}{\_ → ℝ) := IsRapidDatum v ∧ IsRapidDatum y ∧ ∃ C, ∀ 𝔠 : SplitTameChar E μ, ∃ c : ℂ, C * ‖c‖ \textasciicircum{} 2 =}\leanquoteline{2}{lValueProdInt 𝔠 (card n - card m / 2) * ℓ 𝔠 ∧ ∀ z, IsIntMatrix O z →}\leanquoteline{2}{thetaLarge X V 𝔠 v (stdRowDatum σ 𝔠.χ z) = c • thetaSmall X V' 𝔠 y (stdColDatum 𝔠.χ z)}}
\expandafter\newcommand\csname LN@stdExchangeofperiods\endcsname{\href{https://github.com/lsernau/grh-lean/blob/72fead102f663d5c7a407964a13dc0cf9d977207/grh-compact/src/New/Exchange/LowRankIdentity.lean\#L111-L205}{{\leanfont stdExchange\_of\_periods}}}
\expandafter\newcommand\csname LN@centreCoeffeqzeroofnotnegative\endcsname{\href{https://github.com/lsernau/grh-lean/blob/72fead102f663d5c7a407964a13dc0cf9d977207/grh-compact/src/New/Exchange/NonnegativeVanishing.lean\#L21-L103}{{\leanfont centreCoeff\_eq\_zero\_of\_not\_negative}}}
\expandafter\newcommand\csname LN@memmatchingPairs\endcsname{\href{https://github.com/lsernau/grh-lean/blob/72fead102f663d5c7a407964a13dc0cf9d977207/grh-compact/src/New/Exchange/PartnersMatch.lean\#L14-L26}{{\leanfont IsPartner.\allowbreak{}mem\_matchingPairs}}}
\expandafter\newcommand\csname LN@existsrapidwithrealSlice\endcsname{\href{https://github.com/lsernau/grh-lean/blob/72fead102f663d5c7a407964a13dc0cf9d977207/grh-compact/src/New/Exchange/RealPair.lean\#L19-L35}{{\leanfont exists\_rapid\_with\_realSlice}}}
\expandafter\newcommand\csname LN@existscontDiffwitt\endcsname{\href{https://github.com/lsernau/grh-lean/blob/72fead102f663d5c7a407964a13dc0cf9d977207/grh-compact/src/New/Exchange/RealPair.lean\#L40-L54}{{\leanfont exists\_contDiff\_witt}}}
\expandafter\newcommand\csname LN@existsstdExchangetop\endcsname{\href{https://github.com/lsernau/grh-lean/blob/72fead102f663d5c7a407964a13dc0cf9d977207/grh-compact/src/New/Exchange/RealPair.lean\#L59-L132}{{\leanfont exists\_stdExchange\_top}}}
\expandafter\newcommand\csname LQ@existsstdExchangetop\endcsname{\leanquoteline{0}{lemma exists\_stdExchange\_top \{t t'\} \{W : Space E a t\} (D₁ :}\leanquoteline{2}{W₁.Splitting W Unit) [IsBaseAdapted D₁] (V' : Space E b t') σ' (hn : 2 ∣ Fintype.card n) μ \{v v₁ :}\leanquoteline{2}{PairFock ..\} (hv : IsRapidDatum v) (r₁ : IsRapidDatum v₁) : ∃ x x', realSlice D.symm V' 0 x' = v ∧}\leanquoteline{2}{StdExchange X (V'.sum W₁) V' μ σ' x x' (equalRankThetaRatio W₁ · v₁)}}
\expandafter\newcommand\csname LN@existsstdRowDatumeqstdDatum\endcsname{\href{https://github.com/lsernau/grh-lean/blob/72fead102f663d5c7a407964a13dc0cf9d977207/grh-compact/src/New/Exchange/StableSpaces.lean\#L15-L21}{{\leanfont IntegralFrame.\allowbreak{}exists\_stdRowDatum\_eq\_stdDatum}}}
\expandafter\newcommand\csname LN@partnermemofstdColDatum\endcsname{\href{https://github.com/lsernau/grh-lean/blob/72fead102f663d5c7a407964a13dc0cf9d977207/grh-compact/src/New/Exchange/StableSpaces.lean\#L29-L51}{{\leanfont partner\_mem\_of\_stdColDatum}}}
\expandafter\newcommand\csname LN@partnermemofunimodularvanishing\endcsname{\href{https://github.com/lsernau/grh-lean/blob/72fead102f663d5c7a407964a13dc0cf9d977207/grh-compact/src/New/Exchange/StableSpaces.lean\#L57-L88}{{\leanfont IntegralFrame.\allowbreak{}partner\_mem\_of\_unimodular\_vanishing}}}
\expandafter\newcommand\csname LN@applystdColDatumeqzero\endcsname{\href{https://github.com/lsernau/grh-lean/blob/72fead102f663d5c7a407964a13dc0cf9d977207/grh-compact/src/New/Exchange/StandardCoefficients.lean\#L18-L78}{{\leanfont IsEquivariantLift.\allowbreak{}apply\_stdColDatum\_eq\_zero}}}
\newcommand{\ConcordanceRows}{\crow{existscommoncosetreps}{\href{https://github.com/lsernau/grh-lean/blob/72fead102f663d5c7a407964a13dc0cf9d977207/grh-compact/src/New/Exchange/BlockPartners.lean\#L20-L29}{{\leanfont exists\_common\_coset\_reps}}}{BlockPartners.lean}{20--29}
\crow{localSection}{\href{https://github.com/lsernau/grh-lean/blob/72fead102f663d5c7a407964a13dc0cf9d977207/grh-compact/src/New/Exchange/BlockPartners.lean\#L33-L34}{{\leanfont localSection}}}{BlockPartners.lean}{33--34}
\crow{matchingPairs}{\href{https://github.com/lsernau/grh-lean/blob/72fead102f663d5c7a407964a13dc0cf9d977207/grh-compact/src/New/Exchange/BlockPartners.lean\#L39-L41}{{\leanfont matchingPairs}}}{BlockPartners.lean}{39--41}
\crow{mulfirstCols}{\href{https://github.com/lsernau/grh-lean/blob/72fead102f663d5c7a407964a13dc0cf9d977207/grh-compact/src/New/Exchange/BlockPartners.lean\#L44-L44}{{\leanfont mul\_firstCols}}}{BlockPartners.lean}{44--44}
\crow{translatemem}{\href{https://github.com/lsernau/grh-lean/blob/72fead102f663d5c7a407964a13dc0cf9d977207/grh-compact/src/New/Exchange/BlockPartners.lean\#L46-L51}{{\leanfont matchingPairs.\allowbreak{}translate\_mem}}}{BlockPartners.lean}{46--51}
\crow{localSectionofprimitive}{\href{https://github.com/lsernau/grh-lean/blob/72fead102f663d5c7a407964a13dc0cf9d977207/grh-compact/src/New/Exchange/BlockPartners.lean\#L55-L58}{{\leanfont localSection\_of\_primitive}}}{BlockPartners.lean}{55--58}
\crow{isClopenprimitive}{\href{https://github.com/lsernau/grh-lean/blob/72fead102f663d5c7a407964a13dc0cf9d977207/grh-compact/src/New/Exchange/BlockPartners.lean\#L60-L65}{{\leanfont TameLocalChar.\allowbreak{}isClopen\_primitive}}}{BlockPartners.lean}{60--65}
\crow{existslinearfactor}{\href{https://github.com/lsernau/grh-lean/blob/72fead102f663d5c7a407964a13dc0cf9d977207/grh-compact/src/New/Exchange/BlockPartners.lean\#L70-L89}{{\leanfont matchingPairs.\allowbreak{}exists\_linear\_factor}}}{BlockPartners.lean}{70--89}
\crow{partnerblock}{\href{https://github.com/lsernau/grh-lean/blob/72fead102f663d5c7a407964a13dc0cf9d977207/grh-compact/src/New/Exchange/BlockPartners.lean\#L94-L108}{{\leanfont partner\_block}}}{BlockPartners.lean}{94--108}
\crow{adjointthetaeqformLift}{\href{https://github.com/lsernau/grh-lean/blob/72fead102f663d5c7a407964a13dc0cf9d977207/grh-compact/src/New/Exchange/CharacterLiftNorm.lean\#L36-L44}{{\leanfont AuxPair.\allowbreak{}adjoint\_theta\_eq\_formLift}}}{CharacterLiftNorm.lean}{36--44}
\crow{liftauxDatumeqliftmodDatum}{\href{https://github.com/lsernau/grh-lean/blob/72fead102f663d5c7a407964a13dc0cf9d977207/grh-compact/src/New/Exchange/CharacterLiftNorm.lean\#L49-L56}{{\leanfont AuxPair.\allowbreak{}lift\_auxDatum\_eq\_lift\_modDatum}}}{CharacterLiftNorm.lean}{49--56}
\crow{levelIndicatorsupportandnormSq}{\href{https://github.com/lsernau/grh-lean/blob/72fead102f663d5c7a407964a13dc0cf9d977207/grh-compact/src/New/Exchange/CharacterLiftNorm.lean\#L62-L75}{{\leanfont AuxPair.\allowbreak{}levelIndicator\_support\_and\_normSq}}}{CharacterLiftNorm.lean}{62--75}
\crow{gramFactor}{\href{https://github.com/lsernau/grh-lean/blob/72fead102f663d5c7a407964a13dc0cf9d977207/grh-compact/src/New/Exchange/CharacterLiftNorm.lean\#L79-L80}{{\leanfont AuxPair.\allowbreak{}gramFactor}}}{CharacterLiftNorm.lean}{79--80}
\crow{rallisTerm}{\href{https://github.com/lsernau/grh-lean/blob/72fead102f663d5c7a407964a13dc0cf9d977207/grh-compact/src/New/Exchange/CharacterLiftNorm.lean\#L84-L86}{{\leanfont Data.\allowbreak{}rallisTerm}}}{CharacterLiftNorm.lean}{84--86}
\crow{existsisSphericalWithiota}{\href{https://github.com/lsernau/grh-lean/blob/72fead102f663d5c7a407964a13dc0cf9d977207/grh-compact/src/New/Exchange/CharacterLiftNorm.lean\#L91-L101}{{\leanfont Data.\allowbreak{}exists\_isSphericalWith\_iota}}}{CharacterLiftNorm.lean}{91--101}
\crow{innerauxLift}{\href{https://github.com/lsernau/grh-lean/blob/72fead102f663d5c7a407964a13dc0cf9d977207/grh-compact/src/New/Exchange/CharacterLiftNorm.lean\#L105-L106}{{\leanfont Data.\allowbreak{}inner\_auxLift}}}{CharacterLiftNorm.lean}{105--106}
\crow{auxLiftinjectiveandhasSuminner}{\href{https://github.com/lsernau/grh-lean/blob/72fead102f663d5c7a407964a13dc0cf9d977207/grh-compact/src/New/Exchange/CharacterLiftNorm.lean\#L113-L132}{{\leanfont Data.\allowbreak{}auxLift\_injective\_and\_hasSum\_inner}}}{CharacterLiftNorm.lean}{113--132}
\crow{liftsauxLift}{\href{https://github.com/lsernau/grh-lean/blob/72fead102f663d5c7a407964a13dc0cf9d977207/grh-compact/src/New/Exchange/CharacterLiftNorm.lean\#L137-L143}{{\leanfont Data.\allowbreak{}lifts\_auxLift}}}{CharacterLiftNorm.lean}{137--143}
\crow{auxLiftinjective}{\href{https://github.com/lsernau/grh-lean/blob/72fead102f663d5c7a407964a13dc0cf9d977207/grh-compact/src/New/Exchange/CharacterLiftNorm.lean\#L147-L148}{{\leanfont Data.\allowbreak{}auxLift\_injective}}}{CharacterLiftNorm.lean}{147--148}
\crow{auxLiftprincipalRep}{\href{https://github.com/lsernau/grh-lean/blob/72fead102f663d5c7a407964a13dc0cf9d977207/grh-compact/src/New/Exchange/CharacterLiftNorm.lean\#L152-L158}{{\leanfont Data.\allowbreak{}auxLift\_principalRep}}}{CharacterLiftNorm.lean}{152--158}
\crow{auxPairliftstrivialLift}{\href{https://github.com/lsernau/grh-lean/blob/72fead102f663d5c7a407964a13dc0cf9d977207/grh-compact/src/New/Exchange/CharacterLiftNorm.lean\#L174-L184}{{\leanfont LValueData.\allowbreak{}auxPair\_lifts\_trivialLift}}}{CharacterLiftNorm.lean}{174--184}
\crow{innertrivialLiftselfeqL}{\href{https://github.com/lsernau/grh-lean/blob/72fead102f663d5c7a407964a13dc0cf9d977207/grh-compact/src/New/Exchange/CharacterLiftNorm.lean\#L191-L252}{{\leanfont LValueData.\allowbreak{}inner\_trivialLift\_self\_eq\_L}}}{CharacterLiftNorm.lean}{191--252}
\crow{posCoordMatrix}{\href{https://github.com/lsernau/grh-lean/blob/72fead102f663d5c7a407964a13dc0cf9d977207/grh-compact/src/New/Exchange/DeterminantVectors.lean\#L27-L28}{{\leanfont posCoordMatrix}}}{DeterminantVectors.lean}{27--28}
\crow{detVectorFn}{\href{https://github.com/lsernau/grh-lean/blob/72fead102f663d5c7a407964a13dc0cf9d977207/grh-compact/src/New/Exchange/DeterminantVectors.lean\#L32-L33}{{\leanfont detVectorFn}}}{DeterminantVectors.lean}{32--33}
\crow{detVectors}{\href{https://github.com/lsernau/grh-lean/blob/72fead102f663d5c7a407964a13dc0cf9d977207/grh-compact/src/New/Exchange/DeterminantVectors.lean\#L41-L42}{{\leanfont detVectors}}}{DeterminantVectors.lean}{41--42}
\crow{frameAtrealunitary}{\href{https://github.com/lsernau/grh-lean/blob/72fead102f663d5c7a407964a13dc0cf9d977207/grh-compact/src/New/Exchange/DeterminantVectors.lean\#L45-L50}{{\leanfont frameAt\_real\_unitary}}}{DeterminantVectors.lean}{45--50}
\crow{posProjeq}{\href{https://github.com/lsernau/grh-lean/blob/72fead102f663d5c7a407964a13dc0cf9d977207/grh-compact/src/New/Exchange/DeterminantVectors.lean\#L53-L54}{{\leanfont posProj\_eq}}}{DeterminantVectors.lean}{53--54}
\crow{existssliceeq}{\href{https://github.com/lsernau/grh-lean/blob/72fead102f663d5c7a407964a13dc0cf9d977207/grh-compact/src/New/Exchange/DeterminantVectors.lean\#L59-L83}{{\leanfont PositivePart.\allowbreak{}exists\_slice\_eq}}}{DeterminantVectors.lean}{59--83}
\crow{surjOnconeintegraldetVectors}{\href{https://github.com/lsernau/grh-lean/blob/72fead102f663d5c7a407964a13dc0cf9d977207/grh-compact/src/New/Exchange/DeterminantVectors.lean\#L88-L144}{{\leanfont surjOn\_cone\_integral\_detVectors}}}{DeterminantVectors.lean}{88--144}
\crow{thetaSmall}{\href{https://github.com/lsernau/grh-lean/blob/72fead102f663d5c7a407964a13dc0cf9d977207/grh-compact/src/New/Exchange/ExchangeDifference.lean\#L17-L18}{{\leanfont thetaSmall}}}{ExchangeDifference.lean}{17--18}
\crow{existsisEquivariantLift}{\href{https://github.com/lsernau/grh-lean/blob/72fead102f663d5c7a407964a13dc0cf9d977207/grh-compact/src/New/Exchange/ExchangeDifference.lean\#L26-L60}{{\leanfont exists\_isEquivariantLift}}}{ExchangeDifference.lean}{26--60}
\crow{thetaOnRealData}{\href{https://github.com/lsernau/grh-lean/blob/72fead102f663d5c7a407964a13dc0cf9d977207/grh-compact/src/New/Exchange/LowRankIdentity.lean\#L21-L22}{{\leanfont thetaOnRealData}}}{LowRankIdentity.lean}{21--22}
\crow{thetaTermstdDatumspec}{\href{https://github.com/lsernau/grh-lean/blob/72fead102f663d5c7a407964a13dc0cf9d977207/grh-compact/src/New/Exchange/LowRankIdentity.lean\#L27-L61}{{\leanfont thetaTerm\_stdDatum\_spec}}}{LowRankIdentity.lean}{27--61}
\crow{thetaOnRealDataboxProdeqmul}{\href{https://github.com/lsernau/grh-lean/blob/72fead102f663d5c7a407964a13dc0cf9d977207/grh-compact/src/New/Exchange/LowRankIdentity.lean\#L65-L71}{{\leanfont thetaOnRealData\_boxProd\_eq\_mul}}}{LowRankIdentity.lean}{65--71}
\crow{thetaLargeeqapplyoneofevenrank}{\href{https://github.com/lsernau/grh-lean/blob/72fead102f663d5c7a407964a13dc0cf9d977207/grh-compact/src/New/Exchange/LowRankIdentity.lean\#L75-L86}{{\leanfont thetaLarge\_eq\_apply\_one\_of\_even\_rank}}}{LowRankIdentity.lean}{75--86}
\crow{equalRankThetaRatio}{\href{https://github.com/lsernau/grh-lean/blob/72fead102f663d5c7a407964a13dc0cf9d977207/grh-compact/src/New/Exchange/LowRankIdentity.lean\#L90-L91}{{\leanfont equalRankThetaRatio}}}{LowRankIdentity.lean}{90--91}
\crow{StdExchange}{\href{https://github.com/lsernau/grh-lean/blob/72fead102f663d5c7a407964a13dc0cf9d977207/grh-compact/src/New/Exchange/LowRankIdentity.lean\#L96-L99}{{\leanfont StdExchange}}}{LowRankIdentity.lean}{96--99}
\crow{stdExchangeofperiods}{\href{https://github.com/lsernau/grh-lean/blob/72fead102f663d5c7a407964a13dc0cf9d977207/grh-compact/src/New/Exchange/LowRankIdentity.lean\#L111-L205}{{\leanfont stdExchange\_of\_periods}}}{LowRankIdentity.lean}{111--205}
\crow{centreCoeffeqzeroofnotnegative}{\href{https://github.com/lsernau/grh-lean/blob/72fead102f663d5c7a407964a13dc0cf9d977207/grh-compact/src/New/Exchange/NonnegativeVanishing.lean\#L21-L103}{{\leanfont centreCoeff\_eq\_zero\_of\_not\_negative}}}{NonnegativeVanishing.lean}{21--103}
\crow{memmatchingPairs}{\href{https://github.com/lsernau/grh-lean/blob/72fead102f663d5c7a407964a13dc0cf9d977207/grh-compact/src/New/Exchange/PartnersMatch.lean\#L14-L26}{{\leanfont IsPartner.\allowbreak{}mem\_matchingPairs}}}{PartnersMatch.lean}{14--26}
\crow{existsrapidwithrealSlice}{\href{https://github.com/lsernau/grh-lean/blob/72fead102f663d5c7a407964a13dc0cf9d977207/grh-compact/src/New/Exchange/RealPair.lean\#L19-L35}{{\leanfont exists\_rapid\_with\_realSlice}}}{RealPair.lean}{19--35}
\crow{existscontDiffwitt}{\href{https://github.com/lsernau/grh-lean/blob/72fead102f663d5c7a407964a13dc0cf9d977207/grh-compact/src/New/Exchange/RealPair.lean\#L40-L54}{{\leanfont exists\_contDiff\_witt}}}{RealPair.lean}{40--54}
\crow{existsstdExchangetop}{\href{https://github.com/lsernau/grh-lean/blob/72fead102f663d5c7a407964a13dc0cf9d977207/grh-compact/src/New/Exchange/RealPair.lean\#L59-L132}{{\leanfont exists\_stdExchange\_top}}}{RealPair.lean}{59--132}
\crow{existsstdRowDatumeqstdDatum}{\href{https://github.com/lsernau/grh-lean/blob/72fead102f663d5c7a407964a13dc0cf9d977207/grh-compact/src/New/Exchange/StableSpaces.lean\#L15-L21}{{\leanfont IntegralFrame.\allowbreak{}exists\_stdRowDatum\_eq\_stdDatum}}}{StableSpaces.lean}{15--21}
\crow{partnermemofstdColDatum}{\href{https://github.com/lsernau/grh-lean/blob/72fead102f663d5c7a407964a13dc0cf9d977207/grh-compact/src/New/Exchange/StableSpaces.lean\#L29-L51}{{\leanfont partner\_mem\_of\_stdColDatum}}}{StableSpaces.lean}{29--51}
\crow{partnermemofunimodularvanishing}{\href{https://github.com/lsernau/grh-lean/blob/72fead102f663d5c7a407964a13dc0cf9d977207/grh-compact/src/New/Exchange/StableSpaces.lean\#L57-L88}{{\leanfont IntegralFrame.\allowbreak{}partner\_mem\_of\_unimodular\_vanishing}}}{StableSpaces.lean}{57--88}
\crow{applystdColDatumeqzero}{\href{https://github.com/lsernau/grh-lean/blob/72fead102f663d5c7a407964a13dc0cf9d977207/grh-compact/src/New/Exchange/StandardCoefficients.lean\#L18-L78}{{\leanfont IsEquivariantLift.\allowbreak{}apply\_stdColDatum\_eq\_zero}}}{StandardCoefficients.lean}{18--78}}
% the notations of the defined functions: one slot for each explicit argument of the Lean declaration
\newcommand{\localSection}[2]{\operatorname{sec}_{#1}(#2)}
\newcommand{\matchingPairs}[2]{\operatorname{Match}(#1,#2)}
\newcommand{\gramFactor}[1]{\operatorname{gram}(#1)}
\newcommand{\rallisTerm}[7]{\operatorname{ral}_{#1,#2,#3}(#4;#5,#6;#7)}
\newcommand{\posCoordMatrix}[4]{\operatorname{pos}_{#1,#2}(#3,#4)}
\newcommand{\detVectorFn}[7]{\operatorname{dv}_{#1,#2}^{#3,#4}(#5,#6;#7)}
\newcommand{\detVectors}[2]{\operatorname{Det}(#1,#2)}
\newcommand{\thetaSmall}[5]{\theta^{\mathrm{sm}}_{#1,#2}(#3;#4,#5)}
\newcommand{\thetaOnRealData}[7]{\theta^{\infty}_{#1,#2}(#3,#4;#5,#6,#7)}
\newcommand{\equalRankThetaRatio}[3]{\varrho(#1;#2,#3)}
\newcommand{\StdExchange}[8]{\operatorname{Exch}(#1,#2,#3;#4,#5;#6,#7;#8)}
\newcommand{\UnitMembers}{305}
\newcommand{\UnitInNote}{47}
\newcommand{\UnitProofs}{122}
\newcommand{\UnitSteps}{879}
\newcommand{\UnitAllSteps}{933}


\linespread{1.08}                       % 11 pt on 14.7 pt
\setlength{\parindent}{0pt}
\setlength{\parskip}{5.5pt plus 1.5pt minus 0.5pt}
\setlength{\abovedisplayskip}{7.5pt plus 2pt minus 3pt}
\setlength{\belowdisplayskip}{7.5pt plus 2pt minus 3pt}
\setlength{\abovedisplayshortskip}{3pt plus 2pt}
\setlength{\belowdisplayshortskip}{6pt plus 2pt minus 2pt}
\emergencystretch=1.5em
\clubpenalty=10000 \widowpenalty=10000 \displaywidowpenalty=10000
\relpenalty=7000 \binoppenalty=9000
\raggedbottom
\pagestyle{fancy}\fancyhf{}\renewcommand{\headrulewidth}{0pt}
\fancyfoot[C]{\sffamily\fontsize{7.6}{9}\selectfont\color{inkfaint}\thepage\,/\,\pageref*{LastPage}}
\setlist{topsep=2pt,partopsep=0pt,parsep=0pt,itemsep=4pt,leftmargin=2.1em,labelsep=0.55em}

\newcommand{\Q}{\mathbb{Q}}
\newcommand{\Z}{\mathbb{Z}}
\newcommand{\F}{\mathbb{F}}
\newcommand{\len}{\operatorname{len}}
\newcommand{\cok}{\operatorname{cok}}
\newcommand{\End}{\operatorname{End}}
\newcommand{\Hom}{\operatorname{Hom}}
\newcommand{\Gal}{\operatorname{Gal}}
\newcommand{\diag}{\operatorname{diag}}
\newcommand{\kw}[1]{\emph{#1}}
\newcommand{\m}[1]{\mbox{$#1$}}        % a short formula that must not be broken across lines
\newcommand{\head}[1]{\par\addvspace{17pt plus 4pt minus 2pt}\Needspace{6\baselineskip}%
  \noindent{\fontsize{13.2}{16}\selectfont\bfseries\color{ink}#1}\par\nopagebreak\vspace{-1pt}}
\newcommand{\sect}[2]{\refstepcounter{section}\label{#2}\par\addvspace{19pt plus 4pt minus 2pt}\Needspace{7\baselineskip}%
  \noindent\llap{{\fontsize{13.2}{16}\selectfont\bfseries\color{inkfaint}\thesection}\hspace{9pt}}%
  {\fontsize{13.2}{16}\selectfont\bfseries\color{ink}#1}\par\nopagebreak\vspace{-1pt}}
\newcommand{\subhead}[1]{\par\addvspace{9pt plus 2pt}\Needspace{4\baselineskip}\noindent\textbf{#1}\enspace\ignorespaces}
% where a statement lives in the Lean archive
\newcommand{\leantag}{{\sffamily\bfseries\fontsize{6.3}{8}\selectfont\color{leancol}\addfontfeatures{LetterSpace=7}LEAN}}
\newcommand{\sep}{\unskip\nobreak\hspace{5pt}{\color{inkfaint}\textperiodcentered}\hspace{5pt plus 2pt}\ignorespaces}
\newcommand{\leanline}[1]{\par\nopagebreak\vspace{-2.5pt}{\fontsize{8.9}{12.4}\selectfont\color{inksoft}\raggedright
  \noindent\llap{\leantag\hspace{9pt}}#1\par}\vspace{1pt}}
% a displayed statement in a tinted box: #1 colour, #2 text
\newcommand{\statement}[2]{\par\addvspace{7pt}\begingroup\setlength{\fboxsep}{8.5pt}\setlength{\fboxrule}{0.6pt}%
  \noindent\fcolorbox{#1!30!white}{#1!3!white}{\parbox{\dimexpr\linewidth-2\fboxsep-2\fboxrule\relax}{#2}}%
  \endgroup\par\addvspace{7pt}}
\newcommand{\nt}[2]{\par\noindent\hangindent=1.1em\hangafter=1 #1\enspace{\color{inksoft}#2}\par\vspace{1pt}}
\newcommand{\legdot}[1]{\tikz[baseline=-0.62ex]\fill[#1] (0,0) circle (2.5pt);\hspace{3pt}}
\newcommand{\crow}[4]{\csname CS@#1\endcsname & #2 & \textcolor{inksoft}{#3} & \textcolor{inksoft}{#4}\\}
\newcommand{\leanquoteline}[2]{\par\noindent\hangindent=\dimexpr#1\fontdimen2\font+2.2em\relax\hangafter=1
  \hspace*{#1\fontdimen2\font}#2}
\newcommand{\quoted}[1]{\par\vspace{3pt}{\ttfamily\fontsize{7.5}{10.4}\selectfont\color{leancol}\raggedright
  \setlength{\parskip}{0pt}\leanquote{#1}\par}\vspace{4pt}}
\newcommand{\bx}[3]{{\sffamily\bfseries\fontsize{7.3}{9.2}\selectfont\color{#1!82!black}#2}\\[1.5pt]#3}
\newcommand{\cs}[2]{\expandafter\newcommand\csname CS@#1\endcsname{#2}}

% theorems
\newtheoremstyle{plainx}{7pt plus 2pt}{7pt plus 2pt}{\itshape}{}{\bfseries}{.}{0.5em}{}
\newtheoremstyle{uprightx}{7pt plus 2pt}{7pt plus 2pt}{}{}{\bfseries}{.}{0.5em}{}
\theoremstyle{plainx}
\newtheorem{theorem}{Theorem}[section]
\newtheorem{lemma}[theorem]{Lemma}
\newtheorem{proposition}[theorem]{Proposition}
\newtheorem{corollary}[theorem]{Corollary}
\theoremstyle{uprightx}
\newtheorem{definition}[theorem]{Definition}
\newtheorem{setting}[theorem]{Setting}          % hypotheses shared by several statements; each of them cites it by number
\newtheorem{remark}[theorem]{Remark}
\newtheorem{example}[theorem]{Example}
\renewcommand{\qedsymbol}{{\setlength{\fboxsep}{0pt}\setlength{\fboxrule}{0.5pt}\fbox{\rule{0pt}{0.46em}\hspace{0.46em}}}}
% what holds throughout the note: the ambient objects only (implicit and instance arguments), never a hypothesis
\newenvironment{conventions}{\par\addvspace{5pt}\noindent\textbf{Conventions.}\enspace\ignorespaces}{\par\addvspace{3pt}}
\makeatletter
\renewenvironment{proof}[1][Proof]{\par\pushQED{\qed}\normalfont\topsep0pt\partopsep0pt
  \trivlist\item[\hskip\labelsep\itshape #1.]\ignorespaces}{\popQED\endtrivlist\@endpefalse}
\makeatother

% where each declaration of the concordance is treated: \cs{key}{place}, the place naming a \label or a lettered item
\cs{existsstdExchangetop}{Theorem 1}
\cs{applystdColDatumeqzero}{Theorem 2}
\cs{stdExchangeofperiods}{Theorem 3}
\cs{auxLiftinjectiveandhasSuminner}{Theorem 4}
\cs{innertrivialLiftselfeqL}{Theorem 5}
\cs{existscommoncosetreps}{Lemma~\ref{l:reps}}
\cs{localSection}{Definition~\ref{d:locsec}}
\cs{matchingPairs}{Definition~\ref{d:match}}
\cs{mulfirstCols}{Lemma~\ref{l:translate}}
\cs{translatemem}{Lemma~\ref{l:translate}}
\cs{localSectionofprimitive}{Lemma~\ref{l:section}}
\cs{isClopenprimitive}{Lemma~\ref{l:clopen}}
\cs{existslinearfactor}{Proposition~\ref{p:factor}}
\cs{partnerblock}{Lemma~\ref{l:block}}
\cs{memmatchingPairs}{Lemma~\ref{l:matched}}
\cs{existsstdRowDatumeqstdDatum}{Lemma~\ref{l:rowstd}}
\cs{partnermemofstdColDatum}{Lemma~\ref{l:stable}}
\cs{partnermemofunimodularvanishing}{Proposition~\ref{p:unimod}}
\cs{thetaSmall}{Definition~\ref{d:thsm}}
\cs{existsisEquivariantLift}{Proposition~\ref{p:lift}}
\cs{thetaOnRealData}{Definition~\ref{d:threal}}
\cs{equalRankThetaRatio}{Definition~\ref{d:exch}}
\cs{StdExchange}{Definition~\ref{d:exch}}
\cs{centreCoeffeqzeroofnotnegative}{Proposition~\ref{p:nonneg}}
\cs{thetaTermstdDatumspec}{Proposition~\ref{p:terms}}
\cs{thetaOnRealDataboxProdeqmul}{Lemma~\ref{l:box}}
\cs{thetaLargeeqapplyoneofevenrank}{Proposition~\ref{p:const}}
\cs{existsrapidwithrealSlice}{Lemma~\ref{l:slice}}
\cs{existscontDiffwitt}{Lemma~\ref{l:witt}}
\cs{posCoordMatrix}{Definition~\ref{d:detvec}}
\cs{detVectorFn}{Definition~\ref{d:detvec}}
\cs{detVectors}{Definition~\ref{d:detvec}}
\cs{frameAtrealunitary}{Lemma~\ref{l:frame}}
\cs{posProjeq}{Lemma~\ref{l:frame}}
\cs{existssliceeq}{Lemma~\ref{l:posslice}}
\cs{surjOnconeintegraldetVectors}{Proposition~\ref{p:cone}}
\cs{adjointthetaeqformLift}{Lemma~\ref{l:adjoint}}
\cs{liftauxDatumeqliftmodDatum}{Lemma~\ref{l:moddatum}}
\cs{levelIndicatorsupportandnormSq}{Lemma~\ref{l:indicator}}
\cs{gramFactor}{Definition~\ref{d:gram}}
\cs{rallisTerm}{Definition~\ref{d:gram}}
\cs{existsisSphericalWithiota}{Lemma~\ref{l:spherical}}
\cs{innerauxLift}{Lemma~\ref{l:inner}}
\cs{liftsauxLift}{Lemma~\ref{l:lifts}}
\cs{auxLiftinjective}{Corollary~\ref{c:inj}}
\cs{auxLiftprincipalRep}{Lemma~\ref{l:rep}}
\cs{auxPairliftstrivialLift}{Lemma~\ref{l:trivlifts}}
% letters that the coordinator may change in one place
\newcommand{\dL}{\check{\Psi}}         % a function on the large side (Lean Φ; the letter Φ is taken by Φ_{d,M})
\newcommand{\dS}{\Psi}                 % a datum on the small side
\newcommand{\sF}{\mathsf{F}}
\newcommand{\sO}{\mathsf{O}}
\newcommand{\sR}{\mathsf{R}}
\newcommand{\ba}{\mathbf{a}}
\newcommand{\bb}{\mathbf{b}}
\newcommand{\bm}{\mathbf{m}}
\newcommand{\bn}{\mathbf{n}}
\newcommand{\br}{\mathbf{r}}
\newcommand{\sW}{\mathsf{W}}
\newcommand{\sV}{\mathsf{V}}
\newcommand{\sX}{\mathsf{X}}
\newcommand{\sY}{\mathsf{Y}}
\newcommand{\sZ}{\mathsf{Z}}
\newcommand{\sI}{\mathsf{I}}
\newcommand{\se}{\mathsf{e}}
\newcommand{\sg}{\mathsf{g}}
\newcommand{\sh}{\mathsf{h}}
\newcommand{\sk}{\mathsf{k}}
\newcommand{\sx}{\mathsf{x}}
\newcommand{\sy}{\mathsf{y}}
\newcommand{\sz}{\mathsf{z}}
\newcommand{\ssx}{\mathsf{s}}
\newcommand{\sq}{\mathsf{q}}
\newcommand{\fc}{\mathfrak{c}}
\newcommand{\fp}{\mathfrak{p}}
\newcommand{\fv}{\mathfrak{v}}
\newcommand{\fx}{\mathfrak{x}}
\newcommand{\fS}{\mathfrak{S}}
\newcommand{\fA}{\mathfrak{A}}
\newcommand{\fa}{\mathfrak{a}}
\newcommand{\fl}{\mathfrak{l}}
\newcommand{\rc}{\mathrm{c}}
\newcommand{\rC}{\mathrm{C}}
\newcommand{\A}{\mathbb{A}}
\newcommand{\CC}{\mathbb{C}}
\newcommand{\RR}{\mathbb{R}}
\newcommand{\NN}{\mathbb{N}}
\newcommand{\GL}{\operatorname{GL}}
\newcommand{\U}{\operatorname{U}}
\newcommand{\tr}{\operatorname{tr}}
\newcommand{\cS}{\mathcal{S}}
\newcommand{\cM}{\mathcal{M}}
\newcommand{\cD}{\mathcal{D}}
\newcommand{\cL}{\mathcal{L}}
\newcommand{\cN}{\mathcal{N}}
\newcommand{\cP}{\mathcal{Q}}
\newcommand{\Mat}[2]{\operatorname{M}_{#1\times#2}}
\newcommand{\rowk}[2]{\operatorname{row}_{#1}(#2)}
\newcommand{\colk}[2]{\operatorname{col}_{#1}(#2)}
\newcommand{\stc}[2]{\dS^{\mathrm{st}}_{#1,#2}}
\newcommand{\str}[3]{\dL^{\mathrm{st}}_{#1,#2,#3}}
\newcommand{\inu}[1]{#1^{\flat}}
\newcommand{\thL}[5]{\theta^{\mathrm{lg}}_{#1,#2}(#3;#4,#5)}
\newcommand{\thI}[5]{\operatorname{I}_{#1,#2}(#3,#4;#5)}
\newcommand{\LP}[2]{\Pi(#1,#2)}
\newcommand{\cf}[3]{\operatorname{coef}_{#1}(#2;#3)}
\newcommand{\loc}[2]{#1\otimes#2}
\newcommand{\negc}{\operatorname{neg}}
\newcommand{\Np}{\operatorname{N}}
\newcommand{\pet}[2]{\langle #1,#2\rangle_{\mathrm{Pet}}}

\begin{document}
\color{ink}

{\raggedright
{\sffamily\bfseries\fontsize{7.2}{9}\selectfont\color{leancol}\addfontfeatures{LetterSpace=8}READBACK NOTE 16\par}
\vspace{5pt}
{\fontsize{21.5}{25.5}\selectfont\bfseries Partner data, vanishing on standard data, and the exchange identity at
the top of the tower\par}
\vspace{5pt}
{\fontsize{12.4}{16}\selectfont\itshape\color{inksoft} One unit of the Lean artifact \ArtifactName, read back into
mathematics\par}}
\vspace{8pt}
{\color{rulegrey}\hrule height 0.6pt}
\vspace{7pt}
\thispagestyle{empty}

\textbf{Summary.} Let $\sX$ be a hermitian space over a CM field and let $\sV$, $\sV'$ be hermitian spaces whose
ranks add up to the rank of $\sX$. The \emph{exchange identity} says that the theta integrals of the pairs
$(\sX,\sV)$ and $(\sX,\sV')$, formed with data that are \emph{partners} at a set of split places, are proportional,
and that the square of the absolute value of the constant is a fixed number times a product of Hecke $L$-values.
This note proves the identity at the top of a tower of spaces, for standard data (Theorem~1). The difference
of the two integrals is an equivariant functional of the partner datum alone; such a functional vanishes on all
standard data once certain Fourier coefficients of one of its values vanish (Theorem~2), and these coefficients
are found by unfolding (Theorem~3). The note also proves two norm formulas for theta lifts to definite groups
(Theorems 4 and~5).

\statement{leancol}{\fontsize{9.6}{13.2}\selectfont\textbf{What this note is.} The artifact \ArtifactName{} is a
body of Lean code: the compact edition of a formal proof, attributed by its authors to L.~Alpöge with the AI system
Claude, whose final theorems are the Riemann hypothesis and its generalisation to Dirichlet $L$-functions. That
headline is a claim. The authors report that every file compiles and that Lean's kernel accepted the proofs; the
handout's record of checks says who checked what. No mathematician has read the argument through, and nothing here
says that the Riemann hypothesis is proved. This note is the first of three on one unit of the artifact, the unit of
the exchange identity, which belongs to the proof of Theorem~C of the handout (integral forms are long). It treats
\UnitInNote{} of the unit's \UnitMembers{} declarations; the notes \emph{The exchange theorem, one square of lifts,
and the triangle of lifts and descents} and \emph{The grid of lifts and descents, and the crude rows} treat the
others. Every statement of the note is proved here from the structured proofs of the unit's dossier, which a
program compared with the Lean text; what the proofs take from other units is recalled in §\ref{s:setting} and
listed as Inputs, and it is not proved here. The five theorems, the definitions and the statements of
§\ref{s:local} and §\ref{s:lift} were compared with their Lean statements by the writer; §\ref{s:notes} says what
was not. A second reader compared all statements of the unit and a passage of every structured proof with the Lean text after the dossier was made; where her reading differs from the dossier's, this note follows hers (§\ref{s:notes}). The Lean text read is \PinText; Appendix~A is the concordance of \ConcDecls{} declarations.}

\sect{The statements}{s:intro}
The notions used in the five theorems are recalled or defined in §\ref{s:setting}: hermitian spaces and their
splittings, real data, split tame characters $\fc$ with their local characters $\chi_{\fc}$, the theta integral
on the large side, the standard data $\str{\se}{\chi}{\sz}$ and $\stc{\chi}{\sz}$, the products $\LP{\fc}{e}$ of
$L$-values, and equivariant lifts. Four notions are defined later: the theta integral on the small side
(Definition~\ref{d:thsm}), the theta functional on real data (Definition~\ref{d:threal}), and the equal-rank
ratio and the standard exchange identity (Definition~\ref{d:exch}).

\statement{arith}{\phantomsection\label{t:1}\textbf{Theorem 1} (the exchange identity at the top of the
tower).\enspace \emph{Let $E$ be a CM field with a CM type and Weil data, and let $\sX$, $\sW_1$, $\sW$, $\sV'$ be
hermitian spaces over $E$ with index sets $\bm$, $\bn$, $\ba$, $\bb$. Let $\fS$ be a base-adapted splitting of
$\sX$ over $\sW_1$ along $\bb$ and $\fS_1$ a base-adapted splitting of $\sW_1$ over $\sW$ along one line. Suppose
that $|\bn|$ is even. Let $\se'$ be an identification $\bb\sqcup(\bb\sqcup\bn)\simeq\bm$, let $\mu$ be weights,
and let $\fv$ and $\fv_1$ be rapid real data for the pairs $(\sW_1,\sV')$ and $(\sW_1,\sW_1)$. Then there are real
data $\fx$ for $(\sX,\sV'\oplus\sW_1)$ and $\fx'$ for $(\sX,\sV')$ such that the slice of $\fx'$ at $0$ for the
splitting opposite to $\fS$ is $\fv$, and such that}
\[\StdExchange{\sX}{\sV'\oplus\sW_1}{\sV'}{\mu}{\se'}{\fx}{\fx'}{\fc\mapsto\equalRankThetaRatio{\sW_1}{\fc}{\fv_1}}\]
\emph{holds: $\fx$ and $\fx'$ are rapid, and there is a real number $\rC$ such that every split tame character
$\fc$ of weight $\mu$ has a complex number $\rc$ with
$\rC\,|\rc|^2=\LP{\fc}{|\bn|/2}\cdot\equalRankThetaRatio{\sW_1}{\fc}{\fv_1}$ and}
\[\thL{\sX}{\sV'\oplus\sW_1}{\fc}{\fx}{\str{\se'}{\chi_{\fc}}{\sz}}
=\rc\cdot\thetaSmall{\sX}{\sV'}{\fc}{\fx'}{\stc{\chi_{\fc}}{\sz}}\]
\emph{for every integral $\bb\times\bm$ matrix $\sz$ over the product of the completions of $E$ at the places of
$\fc$.}}
\leanline{\lean{existsstdExchangetop}}

\statement{arith}{\phantomsection\label{t:2}\textbf{Theorem 2} (vanishing on standard data).\enspace
\emph{Let $E$ be a CM field with a CM type and Weil data, with the element $\jmath$; let $\sX$, $\sW_1$, $\sW$,
$\sV'$ be hermitian spaces over $E$ with index sets $\bm$, $\bn$, $\ba$, $\bb$; let $\fc$ be a split tame character
and $\chi'$ a multiplicative character of the product of the completions of $E$ at the places of $\fc$. Let
$(\cM,\cD)$ be an equivariant lift for $\fc$, $\sX$ and $\chi'$, let $\fS$ be a base-adapted splitting of $\sX$
over $\sW_1$ along $\bb$ and $\fS_1$ a base-adapted splitting of $\sW_1$ over $\sW$ along one line. Write
$\dS_0=\stc{\chi_{\fc}}{\fA_3'}$ for the standard datum of the block $\fA_3'$ of the frame of $\fS$ adapted to the places
of $\fc$. Suppose that $\dS_0$ lies in $\cM$, and that the centre coefficient}
\[\cf{\fS}{\jmath T}{\cD(\dS_0)}\bigl(\operatorname{rad}_{\fS}(\mathsf{c},\mathsf{w})\,\operatorname{lev}_{\fS}(\sq,\sg)\bigr)\]
\emph{vanishes for all $T$, $\mathsf{c}$, $\mathsf{w}$, $\sq$, $\sg$, where $T$ is a $\bb\times\bb$ matrix over
$E$ and $\operatorname{rad}_{\fS}(\mathsf{c},\mathsf{w})$, $\operatorname{lev}_{\fS}(\sq,\sg)$ are a radical and a
Levi element of the parabolic subgroup of $\fS$, such that one of the following holds:}
\begin{enumerate}[label=\textup{(\roman*)},leftmargin=2.3em,itemsep=1pt]
\item \emph{$T$ is the Gram matrix $\mathrm{J}_{\sV'}$ of $\sV'$, and there are an adelic $\bb\times\bb$ matrix $Y$ and
$\sk\in\GL_{\bb}(\hat{\mathcal{O}}_E)$ with $\sq=\mathrm{sc}_E\,\mathrm{x}_{\sV'}\,Y$ such that $Y_f\,\sk$ is unitary for the
identity form;}
\item \emph{$T$ is invertible and, as an adelic matrix, is not represented by the signs of $\sV'$.}
\end{enumerate}
\emph{Then for every identification $\se'\colon\bb\sqcup(\bb\sqcup\bn)\simeq\bm$ and every integral $\sz$,
that is, every $\bb\times\bm$ matrix $\sz$ with entries in the integers of the product of the completions at the
places of $\fc$:} \m{\cD(\stc{\chi_{\fc}}{\sz})=0}.}
\leanline{\lean{applystdColDatumeqzero}}

\statement{arith}{\phantomsection\label{t:3}\textbf{Theorem 3} (the exchange identity from a period
condition).\enspace \emph{Let $E$ be a CM field with a CM type and Weil data; let $\sX$, $\sW_1$, $\sW$, $\sV'$ be
hermitian spaces over $E$ with index sets $\bm$, $\bn$, $\ba$, $\bb$, let $\fS$ be a base-adapted splitting of $\sX$
over $\sW_1$ along $\bb$ and $\fS_1$ a base-adapted splitting of $\sW_1$ over $\sW$ along one line, and suppose
that $|\bn|$ is even. Let $\mu$ be weights, $\se'$ an identification $\bb\sqcup(\bb\sqcup\bn)\simeq\bm$, and let
$\fv$ and $\fv'$ be rapid real data for $(\sX,\sV'\oplus\sW_1)$ and $(\sX,\sV')$. Let $\mathrm{l}$ be a complex
function on the split tame characters of weight $\mu$, and $\Omega$ a real function on the real points of
$\U(\sV'\oplus\sW_1)$ which is non-negative and does not vanish near~$1$. Assume:}

\emph{(the period condition) for every split tame character $\fc$ of weight $\mu$, with
$\omega=\omega_{\fc}^{-1}\upsilon^{|\bn|/2}$, and every real point $\sh$ of $\U(\sV'\oplus\sW_1)$ there is a real
point $\mathsf{p}$ of $\U(\sV')\times\U(\sW_1)$ such that, for all $B\in\GL_{\bb}(\A_E)$, with $\det_\infty B$ the idele that is $\det B$ at the
infinite places and $1$ at the finite ones, all characters $\omega'$, finite data $F$, adelic matrices $\sx$ and points
$\sg$, $\sy$:}
\begin{multline*}
\int_{[\U(\sW_1)]}\thetaOnRealData{\sW_1}{\sV'\oplus\sW_1}{\omega}{\omega\omega'}{F\boxtimes F_{\fc}}{\tbinom{\sx}{0}}{(\sg,\sy\oplus\mathsf{u})}\\
\Bigl(\omega_\infty(1,\mathsf{p}\sh)\,\mathrm{sl}_{\fS}\bigl((\mathsf{p}\sh)^{-1}\tbinom{B}{0};\fv\bigr)\Bigr)\,d\mathsf{u}\\
=\frac{\Omega(\sh)\,\mathrm{l}(\fc)}{\omega^{\mathrm{bc}}_{\fc}(\det_\infty B)\,|\det_\infty B|^{|\bn|/2}}\;
\thetaOnRealData{\sW_1}{\sV'}{\omega}{\omega'}{F}{\sx}{(\sg,\sy)}\bigl(\mathrm{sl}_{\fS}(B;\fv')\bigr);
\end{multline*}
\emph{here $\omega^{\mathrm{bc}}_{\fc}$ is the base change of $\omega_{\fc}$, a character of the ideles of $E$, and $|\cdot|$ the idele absolute value; $F_{\fc}$ is the finite datum for $(\sW_1,\sW_1)$ named by $\stc{\chi_{\fc}}{1}$ at the places of $\fc$,
$\omega_\infty$ is the real Weil representation of weights $|\bn|/2-\mu$ and $0$, and $\mathrm{sl}_{\fS}$ is the
real slice;}

\emph{(the support or sign condition) either every non-zero real slice $\mathrm{sl}_{\fS}(\sx;\fv)$ has a Gram
matrix $\sx^{*}\operatorname{diag}(\sV'\oplus\sW_1)\sx$, with the diagonal form of the signs of $\sV'\oplus\sW_1$, that is represented by the signs of $\sV'$; or
$\sV'$ is negative definite, $\sW_1$ has a positive sign, and $2\negc(\sW_1)\le|\bn|$.}

\emph{Then}
\[\StdExchange{\sX}{\sV'\oplus\sW_1}{\sV'}{\mu}{\se'}{\fv}{\fv'}{\fc\mapsto|\mathrm{l}(\fc)|^2/\LP{\fc}{|\bn|/2}}\]
\emph{holds.}}
\leanline{\lean{stdExchangeofperiods}}

\statement{arith}{\phantomsection\label{t:4}\textbf{Theorem 4} (the inner product formula for the adjoint
lift).\enspace \emph{Let $P$ be parameters, $D$ admissible data for $P$ and $i$ a lift index, $1\le i\le m$. There
is a complex number $A\ne0$ such that for every tame character $\vartheta$ of $D$ the adjoint lift
$\theta^{*}_{D,\vartheta,i}$, from the forms of principal level on $\sW_i(D)$ to the algebraic forms on the definite
space of the $i$-th auxiliary pair $\fa_{D}(\vartheta,i)$, is injective, and for all forms $y'$, $y$ of principal
level on $\sW_i(D)$}
\[\bigl\langle\theta^{*}_{D,\vartheta,i}\,y',\ \theta^{*}_{D,\vartheta,i}\,y\bigr\rangle
=\sum_{\bar{\sg}}A\cdot\gramFactor{\fa_{D}(\vartheta,i)}\cdot\rallisTerm{D}{\vartheta}{i}{\sW_i(D)}{y'}{y}{\bar{\sg}},\]
\emph{the sum over the cosets $\bar{\sg}$ of Definition~\ref{d:gram} being unconditionally convergent. The number
$A$ does not depend on $\vartheta$.}}
\leanline{\lean{auxLiftinjectiveandhasSuminner}}

\statement{arith}{\phantomsection\label{t:5}\textbf{Theorem 5} (the norm of the lift of the trivial
character).\enspace \emph{Let $E$ be a CM field with a CM type and Weil data, $\bn$ a finite index set and $\fl$
$L$-value data for $E$ and $\bn$, with character $\omega_{\fl}$, with chosen places and twist places, and with
the set $\Sigma(\fl)$ of all these places. Let $e$ be a natural number with $2e=|\bn|$ and $e>1$. Put}
\[\rC=\frac{\operatorname{vol}[\U(\sI_{\bn})]\cdot\operatorname{vol}[\U(\sI_{1})]}{\prod_{\fp}\bigl(1-\operatorname{sg}(\fp)\Np\fp^{-e}\bigr)^{-1}},\]
\emph{the product over all finite places $\fp$ of $E$, with the signs $\operatorname{sg}(\fp)$ of Setting~\ref{set:aux}.
Then $\rC\ne0$, and the lift $\Xi(\fl)$ of the trivial character from $\U(\sI_1)$ to $\U(\sI_{\bn})$ satisfies}
\begin{multline*}
\langle\Xi(\fl),\Xi(\fl)\rangle=\rC\cdot\prod_{\fp\text{ chosen}}(\Np\fp-1)\cdot\operatorname{hook}(\fl)\cdot\operatorname{weyl}(\fl)\\
\cdot\ L^{\Sigma(\fl)\cup\overline{\Sigma(\fl)}}\bigl(e,(\omega_{\fl}^{-1})^{\mathrm{bc}}\bigr)\cdot\Bigl(\prod_{\mathfrak{q}\text{ twist}}\Np\mathfrak{q}\Bigr)^{|\bn|-1},
\end{multline*}
\emph{the products over the chosen places and over the twist places of $\fl$. The $L$-value is that of
§\ref{s:setting}: the Hecke $L$-function of the base change of $\omega_{\fl}^{-1}$ without the Euler factors at the
places of $\fl$ and at their conjugates, at the point $e$. The number $\rC$ depends only on $E$ and $|\bn|$;
$\operatorname{hook}(\fl)$ is the hook number of the weight of $\fl$ and $\operatorname{weyl}(\fl)$ the Weyl index
of Setting~\ref{set:aux}.}}
\leanline{\lean{innertrivialLiftselfeqL}}

Theorem~1 is the input of the second note of the unit, which removes hyperbolic planes from $\sX$ to reach the
exchange theorem for arbitrary complementary spaces and arbitrary partners. Theorems 2 and~3 are its proof.
Theorem~4 is used in Chapter~VIII (the determinants that measure congruences) and Theorem~5 in the second note
(the norm of the lift from $\U(1)$ is an $L$-value).

\sect{The setting}{s:setting}
\begin{conventions}\phantomsection\label{c:conventions}
The natural numbers are $0,1,2,\dots$. The letters of the handout keep their meaning: $P$ are parameters, $D$
admissible data, $m$ the number of lifts, $\vartheta$ a tame character of $D$, $\tau$ an infinite place. The
objects of this note that are local to the theory of theta series are set in other alphabets. Finite index sets
are bold ($\ba$, $\bb$, $\bm$, $\bn$, $\br$, and $\bn'$, $\br'$), $|\bn|$ is the number of elements of $\bn$, and an
$\ba\times\bn$ matrix has rows indexed by $\ba$ and columns by $\bn$; $\Mat{\ba}{\bn}(\sR)$ is the set of such
matrices over a ring $\sR$, and $\GL_{\bn}(\sR)$ the group of invertible $\bn\times\bn$ matrices. Sans-serif letters are a semilocal field $\sF$ with its subring $\sO$, a
commutative ring $\sR$, hermitian spaces, and matrices and points of groups; fraktur letters are real data ($\fv$, $\fx$), places ($\fp$), split tame
character ($\fc$), splittings ($\fS$), integral frames ($\fA$), auxiliary pairs ($\fa$) and $L$-value data ($\fl$);
upright letters ($\rc$, $\rC$) are constants. The implicit arguments of the declarations of the unit are of the
following kinds, and each statement below introduces the ones it uses: a semilocal field with its subring; finite
index sets, one of them non-empty where a statement says so; a commutative ring; a CM field $E$ with a CM type and,
where a statement says so, Weil data; signs and hermitian spaces; splittings, which are base-adapted where a
statement says so; weights and characters; natural numbers; real data, finite data and functions; sets of places and of matrices;
matrices and points; forms; parameters and admissible data with a lift index.
\end{conventions}

\subhead{Matrices over a semilocal field.} A \kw{semilocal field} is what the background library calls so: a
locally compact commutative topological ring $\sF$ with a Haar measure, a compact open subring $\sO$ of
volume~$1$ and an invertible, topologically nilpotent element $\varpi\in\sO$, such that $(\sO,\sF)$ has the
elementary divisor property of a product of valuation rings. The example is the product of the completions of $E$
at a finite set of finite places, with its integers. A matrix over $\sF$ is \kw{integral} if its entries lie in
$\sO$; a matrix $\sy\in\Mat{\bn}{\bb}(\sF)$ \kw{has a left inverse} if $\ssx\,\sy=1$ for some $\ssx$, and is
\kw{primitive} if it is integral with an integral left inverse. A square matrix is primitive exactly when it lies
in $\GL_{\bn}(\sO)$. For a set $\mathcal{U}$ of matrices, $\cS(\mathcal{U})$ is the space of locally constant,
compactly supported complex functions that vanish outside $\mathcal{U}$ (Schwartz--Bruhat functions supported in
$\mathcal{U}$); $\cS^{\mathrm{inj}}$ is this space for the matrices with a left inverse and $\cS^{\mathrm{sur}}$
for the matrices whose transpose has one.

Let $\se\colon\bb\sqcup\ba\simeq\bn$ be an identification and $\sk\in\GL_{\bn}(\sR)$. Then
$\colk{\se}{\sk}\in\Mat{\bn}{\bb}(\sR)$ consists of the columns of $\sk$ indexed by $\bb$, and
$\rowk{\se}{\sk}\in\Mat{\ba}{\bn}(\sR)$ of the rows of $\sk^{-1}$ indexed by $\ba$; so
$\rowk{\se}{\sk}\,\colk{\se}{\sk}=0$. Let $\chi$ be a multiplicative character of $\sR$ and $\sO'\subseteq\sR$ a
subring. Complex functions $\dL$ on $\Mat{\ba}{\bn}(\sR)$ and $\dS$ on $\Mat{\bn}{\bb}(\sR)$ are \kw{partners for
$\chi$ over $\sO'$ with respect to $\se$} if $\dL$ vanishes off the matrices $\rowk{\se}{\sk}$ and $\dS$ off the
matrices $\colk{\se}{\sk}$, $\sk\in\GL_{\bn}(\sO')$, and
\[\dL(\rowk{\se}{\sk})=\chi(\det\sk)\,\dS(\colk{\se}{\sk})\qquad\text{for all }\sk\in\GL_{\bn}(\sO').\]
The kernel of a surjective $\ba\times\bn$ matrix is the column space of an injective $\bn\times\bb$ matrix, and
partners are functions that correspond under this duality of Grassmannians, twisted by $\chi$.

For a multiplicative character $\chi$ of $\sF$ let $\inu{\chi}$ be the function equal to $\chi^{-1}$ on the units
of $\sO$ and to $0$ elsewhere. The character $\chi$ is \kw{tame} if there are a finite ring $\bar{\sR}$, a
surjective homomorphism $\sO\to\bar{\sR}$ with kernel $\varpi\sO$ and a primitive multiplicative character $\bar\chi$
of $\bar{\sR}$ such that $\inu{\chi}$ is $\bar\chi$ composed with the reduction $\sz\mapsto\bar{\sz}$. Here a
multiplicative character $\bar\chi$ of a finite commutative ring $\bar{\sR}$ (extended by $0$ to the non-units) is
\kw{primitive} if for every $\mathsf{a}\ne0$ in $\bar{\sR}$ there is $\mathsf{c}$ with
$\bar\chi(1+\mathsf{c}\mathsf{a})\notin\{0,1\}$, and an additive character $\bar{\mathrm{e}}$ of $\bar{\sR}$ is
\kw{primitive} if for every $\mathsf{a}\ne0$ there is $\mathsf{c}$ with $\bar{\mathrm{e}}(\mathsf{c}\mathsf{a})\ne1$,
that is, if it is non-trivial on every non-zero ideal. For an integral
$\ssx\in\Mat{\bb}{\bn}(\sF)$ the \kw{standard column datum} and, for an identification
$\se\colon\bb\sqcup\ba\simeq\bn$, the \kw{standard row datum} are the functions
\[\stc{\chi}{\ssx}(\sy)=\inu{\chi}(\det(\ssx\,\sy)),\qquad
\str{\se}{\chi}{\ssx}(\sx)=\inu{\chi}\Bigl(\det\tbinom{\ssx}{\sx}\Bigr)\]
of an integral $\sy\in\Mat{\bn}{\bb}(\sF)$ and of an $\sx\in\Mat{\ba}{\bn}(\sF)$ for which the square matrix with
the rows of $\ssx$ and of $\sx$, ordered by $\se$, is integral; they are $0$ at the other matrices. They are
partners for $\chi$ over $\sO$ with respect to $\se$ (Input~1). An \kw{integral frame} $\fA$ for
$(\bm,\bn,\bb)$ is an integral invertible $\bm\times(\bb\sqcup\bn\sqcup\bb)$ matrix with integral inverse, cut into
column blocks $(\fA_1\mid\fA_2\mid\fA_3)$, the inverse being cut into row blocks $\fA_1'$, $\fA_2'$, $\fA_3'$; for a function $G$, the
function $\mathrm{st}_{\fA}(G)$ of an $\ba\times\bm$ matrix $\sx$ is $G(\sx\,(\fA_1\mid\fA_2))$ if $\sx$ is integral and
$0$ otherwise. An additive character $\mathrm{e}_{\sF}$ of $\sF$ \kw{has conductor $\sO$} if it is continuous and
is trivial on $\sz\sO$ exactly when $\sz\in\sO$, for elements $\sz$ of $\sF$.

\subhead{Hermitian spaces, data and theta integrals.} $E$ is a CM field, $E^{+}$ its maximal totally real
subfield, $\A_E$ its ring of adeles and $\hat{\mathcal{O}}_E$ the ring of integral finite adeles. A \kw{CM type}
chooses one complex embedding above each infinite place $\tau$. \kw{Weil data} for $E$ consist of: the
hypothesis that $E$ contains a primitive root of unity of order $5040=7!$ (the formal text calls such an $E$
\emph{stably unramified}; then $\sqrt{-1}\in E$, so $E=E^{+}(\sqrt{-1})$ is unramified over $E^{+}$ at every odd
place); an element $\jmath\in E$ with $\bar\jmath=-\jmath$, with positive imaginary part at every embedding of the
CM type and with $(1-\jmath^{-1})/2$ integral; a finite idele that makes the integral adeles self-dual for the
standard additive character $\mathrm{e}_{\A}$ twisted by $\jmath$ (from it come the scale idele $\mathrm{sc}_E$ and
the idele of the inverse different); and a Hecke character $\upsilon$ of $E$ (a continuous unitary character of
the ideles, trivial on $E^{\times}$), the \kw{splitting character}, with
$\upsilon(\mathsf{x})=\prod_\tau\mathsf{x}_\tau/|\mathsf{x}_\tau|$ when the finite part of the idele $\mathsf{x}$ is a
unit, with $\upsilon(\bar{\mathsf{x}})=\upsilon(\mathsf{x})^{-1}$, and with $\upsilon(\varpi_{\fp})=-1$ for a
uniformiser $\varpi_{\fp}$ at every finite place $\fp=\bar{\fp}$. The same letter $\upsilon$ denotes its
restriction to the norm-one ideles $\U(1)(\A)$; a power of $\upsilon$ next to a character of $\U(1)(\A)$ is a
power of the restriction, and the local character of a power of $\upsilon$ (below) is formed with the Hecke
character. A
\kw{hermitian space} $\sW$ over $E$ with index set $\bn$ and signs (a function from $\bn$ to $\{\pm1\}$) is given by a hermitian
Gram matrix $\mathrm{J}_{\sW}\in\GL_{\bn}(E)$ and an adelic frame $\mathrm{x}_{\sW}\in\GL_{\bn}(\A_E)$ in which the form is the
diagonal form of the signs; $\negc(\sW)$ is the number of negative signs, $\U(\sW)$ the unitary group,
$[\U(\sW)]=\U(\sW)(E^{+})\backslash\U(\sW)(\A)$, and $\mathcal{K}_f$ the stabilizer of the integral lattice in the
finite adelic points. The measure on $[\U(\sW)]$ is the one induced by the following Haar measure of
$\U(\sW)(\A)$: the product of the measure of the formal text at the infinite places (it calls it the compact-dual
measure) and of the Haar measure of the finite adelic points that gives $\mathcal{K}_f$ the mass~$1$; integrals over
the quotient are computed with a cut-off for the counting measure on the rational points.
$\operatorname{vol}[\U(\sW)]$ is the total mass. All constants of this note ($\mathrm{R}$ in the proof of
Theorem~3, the volumes in Theorem~5 and in Input~6) refer to this normalisation. $\sI_{\bn}$ is the standard positive definite space on $\bn$, and
$\sV'\oplus\sW$ the orthogonal sum.

For hermitian spaces $\sW$, $\sV$ and characters $\omega$, $\omega'$ of $\U(1)(\A)$ the Weil representation
$\omega_{\sW,\sV}(\omega,\omega')$ of $\U(\sW)(\A)\times\U(\sV)(\A)$ acts on the space of \kw{data} of the pair, the
restricted tensor product of the models of the Heisenberg representation of $\operatorname{Hom}(\sW,\sV)$. A datum
of the form $\fv\otimes F$ has a \kw{real datum} $\fv$, a vector of the Fock-type model at the infinite places, and
a \kw{finite datum} $F$; $\fv$ is \kw{rapid} if all its weighted norms are finite. With the theta functional
$\theta$, the \kw{theta kernel} is $\theta_{\sW,\sV}(\omega,\omega';f)(\sg,\sh)=\theta(\omega_{\sW,\sV}(\omega,\omega')(\sg,\sh)f)$,
and the \kw{theta integral} of a datum $f$ is the function on $\U(\sW)(\A)$
\[\thI{\sW}{\sV}{\omega}{\omega'}{f}(\sg)=\int_{[\U(\sV)]}\theta_{\sW,\sV}(\omega,\omega';f)(\sg,\sh)\,d\sh.\]

A \kw{split tame character} $\fc$ of weight $\mu$ ($\mu(\tau)>0$ an integer for each infinite place $\tau$) is a
character $\omega_{\fc}$ of $\U(1)(E^{+})\backslash\U(1)(\A)$ of weight $\mu$ whose tame part is non-trivial exactly
at the places of a non-empty finite set $\Sigma(\fc)$ of finite places of $E$ that contains the conjugate of none of
its members. Its places split in $E/E^{+}$, and $\sF_{\fc}=\prod_{\fp\in\Sigma(\fc)}E_{\fp}$ is a semilocal field.
The \kw{base change} of a character $\omega$ of $\U(1)(\A)$ is the Hecke character
$\omega^{\mathrm{bc}}(\mathsf{x})=\omega(\mathsf{x}/\bar{\mathsf{x}})$ of $E$; we write $\omega^{\mathrm{bc}}_{\fc}$ for
that of $\omega_{\fc}$. For a Hecke character $\xi$ of $E$, a finite set $\Sigma$ of finite places and an integer
$e$, the \kw{local character of $\xi$ in rank $e$} is the multiplicative character of $\sF_\Sigma=\prod_{\fp\in\Sigma}E_{\fp}$
that is $\sz\mapsto\xi(\sz)\,|\sz|^{e/2}$ on the invertible $\sz$, read as ideles supported on
$\Sigma$, with $|\sz|$ the product of the normalised absolute values, and $0$ on the others. Then
$\chi_{\fc}$ is the local character of $\omega^{\mathrm{bc}}_{\fc}$ in rank $0$ at $\Sigma(\fc)$, and
$\chi^{(e)}_{\fc}(\sz)=\chi_{\fc}(\sz)\,|\sz|^{e/2}$ is its local character in rank $e$. All
$\chi^{(e)}_{\fc}$ are tame, and they agree with $\chi_{\fc}$ on the units of the integers of $\sF_{\fc}$. For
a function $\dL$ on matrices over $\sF_{\fc}$ and a real datum $\fv$, the \kw{local datum} $\loc{\fv}{[\dL]}$ has
real part $\fv$, is the function $\dL$ of the components at $\Sigma(\fc)$, and is the vacuum at the other finite
places; $\loc{\fv}{[\dL]^{c}}$ is the same with the components at the conjugate places, read through complex
conjugation. For hermitian spaces $\sW$ (index set $\bn$) and $\sV$ (index set $\ba$) and $\dL$ on
$\Mat{\ba}{\bn}(\sF_{\fc})$,
\[\thL{\sW}{\sV}{\fc}{\fv}{\dL}=\thI{\sW}{\sV}{\omega_{\fc}^{-1}\upsilon^{\lfloor|\bn|/2\rfloor}}{\omega_{\fc}^{-1}\upsilon^{\lfloor|\bn|/2\rfloor}}{\loc{\fv}{[\dL]}}\]
is the theta integral \kw{on the large side}.

\subhead{The partial $L$-function.} For a set $\Sigma$ of finite places of $E$, with $\bar\Sigma$ the set of the
conjugate places, and a character $\omega$ of $\U(1)(\A)$, $L^{\Sigma\cup\bar\Sigma}(s,\omega^{\mathrm{bc}})$ is the
Hecke $L$-function of the base change $\omega^{\mathrm{bc}}$ without the Euler factors at the places of
$\Sigma\cup\bar\Sigma$: for real $s>1$ the Euler product over the finite places
$\fp\notin\Sigma\cup\bar\Sigma$ of $(1-\omega^{\mathrm{bc}}(\varpi_{\fp})\Np\fp^{-s})^{-1}$, with $\varpi_{\fp}$ the
chosen uniformiser at $\fp$ read as an idele, and at $s=1$ its limit from the right. This is the only $L$-function of the three notes of the unit, and it is always written so. For a
finite $\Sigma$ and a natural number $e$ put
\[\Pi(\Sigma,\omega,e)=\Bigl(\prod_{\fp\in\Sigma}\Np\fp\Bigr)^{e^2}\,\prod_{i=1}^{e}\bigl|L^{\Sigma\cup\bar\Sigma}(i,\omega^{\mathrm{bc}})\bigr|^{2}\]
(the norms run over $\Sigma$ only; both products are $1$ for $e=0$), and for an integer $e$
\[\LP{\fc}{e}=\Pi(\Sigma(\fc),\omega_{\fc},e)\ \ (e\ge0),\qquad
\LP{\fc}{e}=\Pi(\Sigma(\fc),\omega_{\fc},-e)^{-1}\ \ (e<0).\]

\subhead{Splittings and centre coefficients.} A \kw{splitting} $\fS$ of $\sX$ over $\sW_1$ along $\bb$ writes
$\sX$ as the orthogonal sum of $|\bb|$ hyperbolic planes, spanned by two isotropic frames in duality indexed by
$\bb$, and of a copy of $\sW_1$; it is \kw{base-adapted} when its frames are the standard ones in the adelic
frames of the two spaces. The \kw{opposite} splitting exchanges the two isotropic frames. The stabilizer of the
first frame is a parabolic subgroup with radical elements $\operatorname{rad}_{\fS}(\mathsf{c},\mathsf{w})$ ($\mathsf{w}$ in the centre of the unipotent
radical, a space of hermitian $\bb\times\bb$ matrices) and Levi elements $\operatorname{lev}_{\fS}(\sq,\sg)$,
$\sq\in\GL_{\bb}(\A_E)$, $\sg\in\U(\sW_1)(\A)$. For a function $f$ on $\U(\sX)(\A)$ and a $\bb\times\bb$ matrix
$T$, the \kw{centre coefficient} $\cf{\fS}{T}{f}$ is the Fourier coefficient of $f$ along the centre of the
radical at the character $\mathsf{w}\mapsto\mathrm{e}_{\A}(\tr(T\mathsf{w}))$. In the mixed model of $\fS$ the theta kernel of
$(\sX,\sV)$ is a sum, over the rational matrices $\mathsf{v}$ that are the images of the first frame, of \kw{theta terms},
and the datum enters a term through its \kw{real slice} $\mathrm{sl}_{\fS}(\sx;\fv)$, a real datum of the smaller
pair. A reading of the components at a set of places $\Sigma$ attaches to a base-adapted $\fS$ an integral frame,
the \kw{adapted frame}. An adelic matrix $T$ is \kw{represented by} signs $\mathsf{t}$ if at every infinite place it has the
form $\mathsf{c}^{*}\operatorname{diag}(\mathsf{t})\,\mathsf{c}$.

\subhead{Equivariant lifts.} Let $\fc$ be a split tame character, $\sX$ a hermitian space with index set $\bm$
and $\chi'$ a multiplicative character of $\sF_{\fc}$. A pair $(\cM,\cD)$ of a space $\cM$ of complex functions on
$\Mat{\bm}{\bb}(\sF_{\fc})$ and a linear map $\cD$ from such functions to functions on $\U(\sX)(\A)$ is an
\kw{equivariant lift} for $\fc$, $\sX$ and $\chi'$ if $\cM\subseteq\cS^{\mathrm{inj}}$; each $\cD(\dS)$, $\dS\in\cM$,
is an automorphic function; for every finite adelic point $\sg$ supported at $\Sigma(\fc)$ and $\dS\in\cM$ the
translate $\dS(\sg_{\Sigma}\,\cdot\,)$ lies in $\cM$ and
$\cD(\dS(\sg_{\Sigma}\,\cdot\,))(\sx\sg)=\chi'(\det\sg_{\Sigma})\,\cD(\dS)(\sx)$; and every $\cD(\dS)$, $\dS\in\cM$,
is invariant under the elements of $\mathcal{K}_f$ that are trivial at $\Sigma(\fc)$.

\statement{inksoft}{\textbf{Inputs.} The proofs below take the following for granted. The numbers are the units of
the reading; the first eight are units of the background library, whose results this write-up quotes, the last
four are units of the argument. References are those of the handout's catalogue; none of their places has been
confirmed yet.

\vspace{3pt}\phantomsection\label{i:gauss}\textbf{Input 1} (unit~103, Gauss sums and standard data).\enspace
The standard row and column data of an integral $\ssx$ are partners; the blocks of an integral frame satisfy
$\fA_1'\fA_1=1$, $\fA_2'\fA_2=1$, $\fA_3'\fA_3=1$ and the six mixed products vanish; the \emph{matrix Gauss sum}: for a primitive
multiplicative character $\bar\chi$ of a finite ring $\bar{\sR}$ and a primitive additive character
$\bar{\mathrm{e}}$ and every square matrix $\bar T$ over $\bar{\sR}$,
\[\sum_{\bar{\sz}}\bar\chi(\det\bar{\sz})\,\bar{\mathrm{e}}(\tr(\bar T\bar{\sz}))
=\bar\chi^{-1}(\det\bar T)\sum_{\bar{\sz}}\bar\chi(\det\bar{\sz})\,\bar{\mathrm{e}}(\tr\bar{\sz})\]
(relation to the literature unresolved: T.~Kondo, J.~Math.\ Soc.\ Japan 15 (1963), pp.~244--245, has the law over
a finite field at invertible $\bar T$; the statement over a finite ring and at all $\bar T$ is wider than the
source); invariance of functionals under shears.

\vspace{3pt}\phantomsection\label{i:la}\textbf{Input 2} (units 112, 120 and 105, linear algebra, Witt's
theorem and a fact on compact groups).\enspace A matrix is primitive exactly when it is $\colk{\se}{\sk}$ for some $\sk\in\GL_{\bn}(\sO)$;
Jacobi's identity for complementary minors,
$\det\tbinom{\ssx}{\rowk{\se}{\sk}}\det\sk=\det(\ssx\,\colk{\se}{\sk})$; a matrix with a left inverse becomes
primitive after right multiplication by an invertible matrix; Witt's extension theorem over a field with
involution and over $\hat{\mathcal{O}}_E$ (U.~First, J.~Pure Appl.\ Algebra 219 (2015); the theorem, and a special
case); for stably unramified $E$ (a hypothesis that is part of the Weil data), a unimodular hermitian matrix
$\mathsf{b}$ over $\hat{\mathcal{O}}_E$ becomes the identity after a change of basis:
$\sk^{*}\mathsf{b}\,\sk=1$ for some $\sk\in\GL_{\bb}(\hat{\mathcal{O}}_E)$ (R.~Jacobowitz, Amer.\ J.~Math.\ 84
(1962); a special case; without the hypothesis the statement fails at places where $E/E^{+}$ is ramified); the
integral over a compact group of a locally constant family of elements of a space of functions is a finite
combination of its members (unit 105).

\vspace{3pt}\phantomsection\label{i:dual}\textbf{Input 3} (units 43, 38 and 37, dual pairs and the mixed
model).\enspace Theta integrals of target-bounded data are automorphic and depend linearly on the datum; local
data of rapid real data and Schwartz--Bruhat functions on surjective matrices are target-bounded;
$\thI{\sW}{\sV}{\omega}{\omega'}{\omega_{\sW,\sV}(\omega,\omega')(\sg,\sh)f}=\thI{\sW}{\sV}{\omega}{\omega'}{f}(\,\cdot\,\sg)$;
the action of $\mathcal{K}_f$ and of points supported at $\Sigma$ on local data; the centre coefficient of a theta integral
at $\jmath T$ is the sum over the rational $\mathsf{v}$ with $\mathsf{v}^{*}\mathrm{J}\mathsf{v}=T$ of the integrals of the theta terms, and for
invertible $T$ an integral over the quotient by the stabilizer of one $\mathsf{v}$; the laws of theta terms under
radical and Levi elements; existence of sections in the mixed model; the linear functional ``centre coefficient of
$\cD(\dS)$'' of an equivariant lift and its behaviour under shears.}

\statement{inksoft}{\textbf{Inputs (continued).}

\vspace{3pt}\phantomsection\label{i:singular}\textbf{Input 4} (units 42, 98, 10 and 123, singular forms).\enspace
\emph{Freitag's lemma}: under invariance, radical and base conditions, the centre coefficient of an automorphic
function at an invertible index vanishes at points supported at $\Sigma$ once it does at unimodular indices
(E.~Freitag, LNM~1487 (1991); J.-S.~Li, Amer.\ J.~Math.\ 119 (1997); not a literal theorem of these sources);
\emph{Landherr's theorem} for spaces with a point and \emph{weak approximation} for unitary groups (W.~Landherr,
Abh.\ Math.\ Sem.\ Hamburg 11 (1936); M.~Kneser, TIFR lectures 47 (1969), §5.5; special cases); a continuous
periodic function on adelic matrices with vanishing Fourier coefficients is zero.

\vspace{3pt}\phantomsection\label{i:tower}\textbf{Input 5} (units 131, 17 and 114, the tower).\enspace The
\emph{tower property}: the constant term of a theta integral on the large side along an isotropic line is a
multiple of the theta integral of the smaller pair (after S.~Rallis, Compositio Math.\ 51 (1984); an analogue);
the \emph{constancy criterion}, a form of strong approximation: let $\sW$ be a hermitian space with a splitting
$\fS_1$ over a space $\sW_0$ along one line, $\Sigma$ a non-empty set of finite places, and $f$ a continuous
automorphic function on $\U(\sW)(\A)$ that is invariant on the right under the finite adelic points supported at
$\Sigma$ and under the elements of $\mathcal{K}_f$ that are trivial at $\Sigma$; suppose that, if $f$ is invariant on
the left under all radical elements $\operatorname{rad}_{\fS_1}(\mathsf{c},\mathsf{w})$, then
$f(\operatorname{lev}_{\fS_1}(\sq,\sg))$ does not depend on the torus coordinate $\sq$; then $f$ is constant
(M.~Kneser, J.~reine angew.\ Math.\ 218 (1965), §1 and Satz~2; a weaker form, a consequence of strong
approximation for the special unitary group with the hypothesis on the torus); \emph{unfolding in stages}: let
$\mathcal{G}_1\times\mathcal{G}_2\to\mathcal{G}$ be a closed embedding of a product of unimodular groups in a unimodular
group, written $(\mathsf{l},\mathsf{a})\mapsto\mathsf{l}\mathsf{a}$, let $\mathcal{G}_0\subset\mathcal{G}_2$ be a discrete
subgroup and $\Omega$ a non-negative function on $\mathcal{G}$ that does not vanish near some point; then there is
$\mathrm{R}\ne0$ such that $\int_{\mathcal{G}_0\backslash\mathcal{G}}f=\mathrm{R}\,\rc$ for every continuous integrable
function $f$ on $\mathcal{G}$, invariant on the left under $\mathcal{G}_0$, and every number $\rc$ with the following
property: for every $\sg\in\mathcal{G}$ there is a shift $\sk\in\mathcal{G}_1\times\mathcal{G}_2$ with
$\int_{\mathcal{G}_1}\int_{\mathcal{G}_0\backslash\mathcal{G}_2}f(\mathsf{l}\,\mathsf{a}\,\sk\,\sg)\,d\mathsf{a}\,d\mathsf{l}=\Omega(\sg)\,\rc$
(reference to be confirmed); the averaging criterion for the vanishing of an integral over a quotient;
$\LP{\fc}{e}>0$.}

\statement{inksoft}{\textbf{Inputs (continued).}

\vspace{3pt}\phantomsection\label{i:rallis}\textbf{Input 6} (units 26 and 27, the inner product formula).\enspace
The inner product formula for theta lifts (the \emph{Rallis inner product formula}) in two forms. Doubling form:
for $\operatorname{rk}\sV>2\operatorname{rk}\sW$, a clean datum $f$, any datum $f'$ and rapidly decreasing
automorphic $y$, $y'$,
\[\langle\theta_{f'}y',\theta_fy\rangle=\operatorname{vol}[\U(\sV)]\int_{\U(\sW)(\A)}\langle f',\omega(\sg,1)f\rangle\,\pet{y'}{R(\sg)y}\,d\sg\]
(relation to the literature unresolved: the identity and its rank range are those of J.-S.~Li, J.~reine angew.\
Math.\ 428 (1992), §2, Thm~2.1, (24)--(27), p.~184, for cuspidal vectors and Schwartz data; the test spaces and
the normalisation of the volume differ; S.~Rallis, LNM~1245 (1987), is background). Euler form: for an edge and
$2\operatorname{rk}\sW<\operatorname{rk}\sV$ there is $A_0\ne0$ such that for all characters of the weights of the
edge and every finite datum $F$ that is sharp of a level with constant $\rc_F$, the lift of forms of that level
is injective and $\langle\theta y',\theta y\rangle=\sum_{\bar{\sg}}A_0\,\rc_F\,\mathrm{ct}(\bar{\sg})$, a sum of
coset terms (relation to the literature unresolved: S.~Rallis, Compositio Math.\ 52 (1984), §1 and Thm~1.6,
treats another dual pair and is an analogue only; the unitary inner product in this rank range is J.-S.~Li,
loc.\ cit., §2, Thm~2.1, eq.~(26), p.~184; the constant uniform in the characters and the sharp data, the
evaluation as a sum of coset terms and the injectivity go beyond both). Sharp data: a datum $F\otimes$vacuum
whose function $F$ modulo a modulus has a fixed left inverse on its support is sharp, with constant
$\sum|F|^2$ divided by the numbers of matrices and of invertible matrices modulo the modulus; the local doubling
sum at one place; partial integrals.

\vspace{3pt}\phantomsection\label{i:hecke}\textbf{Input 7} (units 44, 25, 96, 104 and 102, Hecke transfer,
lifts and forms).\enspace \emph{Hecke transfer}: an equivariant map from data to functions sends the unramified vacuum at a
split place to a spherical vector, and the Hecke polynomials of the two sides are related by
$T\mapsto\mathrm{tr}_{\fp,e}(\mathrm{a}_{\fp}T)\cdot\mathrm{pol}(\Np\fp^{e}\mathrm{a}_{\fp}T/\mathrm{b}_{\fp})$, with the Frobenius values $\mathrm{a}_{\fp}$, $\mathrm{b}_{\fp}$ of the two characters (relation to the literature
unresolved: A.~Mínguez, Ann.\ Sci.\ Éc.\ Norm.\ Sup.\ 41 (2008), Thm~1 and Cor.~6.3, gives the correspondence of the
parameters at split places; the identity of operators and its normalisation by the splitting character need a
further comparison); an edge gives arrows in both directions, and the adjoint of the lift along an arrow
is the conjugate of the descent; Euler products converge for exponents $>1$; a matrix over the residue ring of a
product of distinct places is the family of its reductions at these places (unit 104); a form is spherical with a
given Hecke polynomial exactly when its algebraic form is, and the Petersson product is invariant under
translation (unit 102).

\vspace{3pt}\phantomsection\label{i:new}\textbf{Input 8} (units 134 Forms, 138 Selmer, 140 Trace, 141 Seesaw, and
the second note of this unit).\enspace Auxiliary pairs, their data and class-sum lifts, and the facts on them that
Setting~\ref{set:aux} and the proofs of §\ref{s:norms} name; the chain of lifts and its arrows; $L$-value data and
the lift $\Xi(\fl)$ with its elementary properties (the second note).}

\sect{Partners at the split places}{s:local}
In this section $\sF$ is a semilocal field with subring $\sO$. The functions on the two sides of the exchange
identity live on $\ba\times\bn$ and on $\bn\times\bb$ matrices; they are compared through their averages over
$\GL_{\ba}(\sF)$ and $\GL_{\bb}(\sF)$, which for want of a convergent integral are defined by a choice of coset
representatives.

\begin{lemma}[coset representatives]\label{l:reps}
Let $\sF$ be a semilocal field with subring $\sO$ and let $\bb$, $\bn$ be finite index sets. There are a linear
functional $\cL_0$ on the space of all complex functions on $\GL_{\bb}(\sF)$ and a set
$\mathcal{R}\subset\GL_{\bb}(\sF)$ of representatives of the cosets of $\GL_{\bb}(\sF)$ modulo $\GL_{\bb}(\sO)$ such
that
\begin{enumerate}[label=\textup{(\arabic*)},leftmargin=2.3em,itemsep=1pt]
\item for every finite subset $\mathcal{R}_0\subseteq\mathcal{R}$ and every function $G$ that vanishes on
$\mathcal{R}\smallsetminus\mathcal{R}_0$: $\cL_0(G)=\sum_{\mathsf{r}\in\mathcal{R}_0}G(\mathsf{r})$;
\item for every $\sy\in\Mat{\bn}{\bb}(\sF)$ with a left inverse there is exactly one $\mathsf{r}\in\mathcal{R}$
such that $\sy\,\mathsf{r}$ is primitive, and exactly one $\mathsf{r}\in\mathcal{R}$ such that $\sy\,\mathsf{r}^{-1}$
is primitive.
\end{enumerate}
\end{lemma}
\begin{proof}
By the elementary divisor property, every $\sg\in\GL_{\bb}(\sF)$ has a $\sk\in\GL_{\bb}(\sO)$ with
$\sg^{\mathsf{T}}=\sk^{-\mathsf{T}}\sg\,\sk$, where ${}^{-\mathsf{T}}$ is the inverse transpose (Input~2); that is,
$(\sg\sk)^{\mathsf{T}}=\sg\sk$: the matrix $\sg\sk$ is symmetric. Choose an element $\sg$ of each coset and such a
$\sk$ for it, and let $\mathcal{R}$ consist of the products $\sg\sk$; so every $\mathsf{r}\in\mathcal{R}$ is
symmetric, and $\mathsf{r}^{-1}=\mathsf{r}^{-\mathsf{T}}$. The functional ``sum of the values on $\mathcal{R}$'' is defined on the functions that vanish
at all but finitely many points of $\mathcal{R}$; extend it linearly to all functions. This gives~(1). For~(2): a
matrix with a left inverse becomes primitive after right multiplication by some $\mathsf{u}\in\GL_{\bb}(\sF)$, and
if $\sy\mathsf{u}$ is primitive then $\sy\mathsf{u}\mathsf{u}'$ is primitive exactly when
$\mathsf{u}'\in\GL_{\bb}(\sO)$ (Input~2). So the $\mathsf{u}$ with $\sy\mathsf{u}$ primitive form one coset, which
contains one element of $\mathcal{R}$. The same holds for the inverses: the inverse transpose is an automorphism of
$\GL_{\bb}(\sF)$ that preserves $\GL_{\bb}(\sO)$, so it permutes the cosets, and the matrices
$\mathsf{r}^{-\mathsf{T}}=\mathsf{r}^{-1}$, $\mathsf{r}\in\mathcal{R}$, are again representatives of all cosets, one
for each.
\end{proof}
\leanline{\lean{existscommoncosetreps}}

\begin{definition}[local section]\label{d:locsec}
Let $\sF$ be a semilocal field, $\bb$, $\bn$ finite index sets, $\chi$ a multiplicative character of $\sF$ and
$\sy\in\Mat{\bn}{\bb}(\sF)$. With the functional $\cL_0$ and the set $\mathcal{R}$ of Lemma~\ref{l:reps}, chosen
once for $(\sF,\bb,\bn)$, $\localSection{\chi}{\sy}$ is the linear functional
\[\localSection{\chi}{\sy}(\dS)=\cL_0\bigl(\sg\mapsto\chi(\det\sg)\,\dS(\sy\,\sg)\bigr)\]
on the complex functions $\dS$ on $\Mat{\bn}{\bb}(\sF)$. By Lemma~\ref{l:reps}\,(1) it equals the sum
$\sum_{\mathsf{r}\in\mathcal{R}}\chi(\det\mathsf{r})\dS(\sy\,\mathsf{r})$ whenever only finitely many of its terms
are non-zero: it stands for the integral of $\chi(\det\sg)\dS(\sy\sg)$ over $\GL_{\bb}(\sF)$.
\end{definition}
\leanline{\lean{localSection}}

\begin{definition}[matched pairs]\label{d:match}
Let $\sF$ be a semilocal field, $\ba$, $\bb$, $\bn$ finite index sets, $\chi$ a multiplicative character of $\sF$
and $\se\colon\bb\sqcup\ba\simeq\bn$ an identification. $\matchingPairs{\chi}{\se}$ is the space of the pairs
$(\dL,\dS)\in\cS^{\mathrm{sur}}\times\cS^{\mathrm{inj}}$ of Schwartz--Bruhat functions, on $\Mat{\ba}{\bn}(\sF)$
and on $\Mat{\bn}{\bb}(\sF)$, such that for every $\sk\in\GL_{\bn}(\sF)$
\[\localSection{\chi}{\rowk{\se}{\sk}^{\mathsf{T}}}\bigl(\dL\circ{}^{\mathsf{T}}\bigr)
=\chi(\det\sk)\cdot\localSection{\chi}{\colk{\se}{\sk}}(\dS).\]
On the left stands the local section for the index sets $(\ba,\bn)$, applied to the function
$\sx\mapsto\dL(\sx^{\mathsf{T}})$ of $\bn\times\ba$ matrices.
\end{definition}
\leanline{\lean{matchingPairs}}

\begin{lemma}[translation]\label{l:translate}
Let $\sF$ be a semilocal field and $\ba$, $\bb$, $\bn$ finite index sets.
\begin{enumerate}[label=\textup{(\alph*)},leftmargin=2.3em,itemsep=1pt]
\item For an identification $\se\colon\bb\sqcup\ba\simeq\bn$ and $\sg,\sk\in\GL_{\bn}(\sF)$:
$\sg\,\colk{\se}{\sk}=\colk{\se}{\sg\sk}$.
\item Let $\chi$ be a multiplicative character of $\sF$ and $\se\colon\bb\sqcup\ba\simeq\bn$. If the pair
$(\dL,\dS)$ lies in $\matchingPairs{\chi}{\se}$ and $\sg\in\GL_{\bn}(\sF)$, then the pair
$\bigl(\chi(\det\sg)\,\dL(\,\cdot\,\sg),\ \dS(\sg^{-1}\,\cdot\,)\bigr)$ lies in $\matchingPairs{\chi}{\se}$.
\end{enumerate}
\end{lemma}
\begin{proof}
(a) holds by definition. (b) Translation by an invertible matrix is a homeomorphism that preserves the matrices
with a left inverse, so both functions stay Schwartz--Bruhat with the required supports. Since
$\rowk{\se}{\sg^{-1}\sk}=\rowk{\se}{\sk}\,\sg$ and, by (a), $\colk{\se}{\sg^{-1}\sk}=\sg^{-1}\colk{\se}{\sk}$,
the two sides of the matching identity of the new pair at $\sk$ are $\chi(\det\sg)$ times those of the old pair
at $\sg^{-1}\sk$, the factor $\chi(\det(\sg^{-1}\sk))$ becoming $\chi(\det\sk)$.
\end{proof}
\leanline{\lean{mulfirstCols}\sep \lean{translatemem}}

\begin{lemma}[local sections of primitive supports]\label{l:section}
Let $\sF$ be a semilocal field, $\bb$, $\bn$ finite index sets, $\chi$ a multiplicative character of $\sF$ and
$\sy\in\Mat{\bn}{\bb}(\sF)$ a matrix that has a left inverse. There is $\sg\in\GL_{\bb}(\sF)$ such that $\sy\sg$
is primitive and, for every function $\dS$ that vanishes off the primitive matrices,
$\localSection{\chi}{\sy}(\dS)=\chi(\det\sg)\,\dS(\sy\sg)$.
\end{lemma}
\begin{proof}
Take for $\sg$ the unique $\mathsf{r}\in\mathcal{R}$ with $\sy\mathsf{r}$ primitive (Lemma~\ref{l:reps}\,(2)). The
function $\sh\mapsto\chi(\det\sh)\dS(\sy\sh)$ vanishes at the other elements of $\mathcal{R}$, so
Lemma~\ref{l:reps}\,(1) with the set $\{\mathsf{r}\}$ gives the value.
\end{proof}
\leanline{\lean{localSectionofprimitive}}

\begin{lemma}[primitive matrices form a clopen set]\label{l:clopen}
Let $\sF$ be a semilocal field, $\ba$, $\bb$, $\bn$ finite index sets, $\chi$ a multiplicative character of
$\sF$ that is tame, and $\se\colon\bb\sqcup\ba\simeq\bn$. Then the set of primitive matrices in
$\Mat{\bn}{\ba}(\sF)$ is closed and open.
\end{lemma}
\begin{proof}
It is the image of the compact group $\GL_{\bn}(\sO)$ under the continuous map that takes the columns indexed by
$\ba$ (Input~2), hence closed. Let $\sy_0$ be primitive with integral left inverse $\ssx$. The standard
datum $\stc{\chi}{\ssx}$ is locally constant, its value at $\sy_0$ is $\inu{\chi}(1)=1$, and every matrix at which
it does not vanish is primitive (Input~1). So a neighbourhood of $\sy_0$ consists of primitive
matrices. The tame character is needed only for the standard datum to be locally constant.
\end{proof}
\leanline{\lean{isClopenprimitive}}

\begin{proposition}[the large side factors through the partner]\label{p:factor}
Let $\sF$ be a semilocal field, $\ba$, $\bb$, $\bn$ finite index sets, $\chi$ a multiplicative character of $\sF$
that is tame, $\se\colon\bb\sqcup\ba\simeq\bn$, and $\mathcal{X}$ a set. Let $\cL$ be a linear map from the complex
functions on $\Mat{\ba}{\bn}(\sF)$ to the complex functions on $\mathcal{X}$ such that
$\chi(\det\sg)\cdot\cL(\dL(\sg\,\cdot\,))=\cL(\dL)$ for every $\sg\in\GL_{\ba}(\sF)$ and every
$\dL\in\cS^{\mathrm{sur}}$. Then there is a linear map $\cL'$ from the complex functions on $\Mat{\bn}{\bb}(\sF)$
to the complex functions on $\mathcal{X}$ with $\cL'(\dS)=\cL(\dL)$ for every pair
$(\dL,\dS)\in\matchingPairs{\chi}{\se}$.
\end{proposition}
\begin{proof}
A matrix $\sx\in\Mat{\ba}{\bn}(\sF)$ has a primitive transpose exactly when $\sx=\rowk{\se}{\sk}$ for some
$\sk\in\GL_{\bn}(\sO)$ (Input~2); choose such a $\sk(\sx)$ for each of them. Define $\cL'(\dS)$ as $\cL$
of the function that is $\chi(\det\sk(\sx))\cdot\localSection{\chi}{\colk{\se}{\sk(\sx)}}(\dS)$ at the $\sx$ with
primitive transpose and $0$ elsewhere. It is linear in $\dS$.

Let $(\dL,\dS)$ be a matched pair, and let $\mathcal{R}$ be the set of Lemma~\ref{l:reps} for $(\ba,\bn)$. For
$\mathsf{r}\in\mathcal{R}$ let $\dL_{\mathsf{r}}$ be the product of $\dL$ with the indicator of the set of the
$\sx$ for which $\sx^{\mathsf{T}}\mathsf{r}^{-1}$ is primitive. This set is closed and open by
Lemma~\ref{l:clopen}, so $\dL_{\mathsf{r}}\in\cS^{\mathrm{sur}}$. Every $\sx$ in the support of $\dL$ has a
transpose with a left inverse, so it lies in exactly one of these sets (Lemma~\ref{l:reps}\,(2)). The support is
compact, hence covered by the sets of a finite subset $\mathcal{R}_0\subseteq\mathcal{R}$, and
$\dL=\sum_{\mathsf{r}\in\mathcal{R}_0}\dL_{\mathsf{r}}$. By the invariance of $\cL$ under the transposes of the
elements of $\mathcal{R}_0$,
\[\cL(\dL)=\sum_{\mathsf{r}\in\mathcal{R}_0}\chi(\det\mathsf{r})\,\cL\bigl(\dL_{\mathsf{r}}(\mathsf{r}^{\mathsf{T}}\,\cdot\,)\bigr)
=\cL\Bigl(\sx\mapsto[\sx^{\mathsf{T}}\text{ primitive}]\sum_{\mathsf{r}\in\mathcal{R}_0}\chi(\det\mathsf{r})\,\dL(\mathsf{r}^{\mathsf{T}}\sx)\Bigr).\]
If $\sx^{\mathsf{T}}$ is primitive and $\dL(\mathsf{r}^{\mathsf{T}}\sx)\ne0$, then $\mathsf{r}\in\mathcal{R}_0$,
because $\mathsf{r}$ is the representative that belongs to the point $\mathsf{r}^{\mathsf{T}}\sx$ of the support.
So by Lemma~\ref{l:reps}\,(1) the inner sum is $\localSection{\chi}{\sx^{\mathsf{T}}}(\dL\circ{}^{\mathsf{T}})$, and
by the matching identity at $\sk(\sx)$ it equals
$\chi(\det\sk(\sx))\cdot\localSection{\chi}{\colk{\se}{\sk(\sx)}}(\dS)$. This is the function that defines
$\cL'(\dS)$.
\end{proof}
\leanline{\lean{existslinearfactor}}

\begin{lemma}[partners from a function of a block]\label{l:block}
Let $\ba$, $\bb$, $\bn$ be finite index sets, $\sR$ a commutative ring, $\se\colon\bb\sqcup\ba\simeq\bn$ an
identification, $\chi$ a multiplicative character of $\sR$ and $G$ a complex function on $\Mat{\ba}{\bb}(\sR)$. Cut
an $\ba\times\bn$ matrix $\sx$ by $\se$ into column blocks $(\sx_1\mid\sx_2)$ with $\sx_2$ square, and an
$\bn\times\bb$ matrix $\sy$ into row blocks $\sy_1$ (square) and $\sy_2$. Then the functions
\[\dL(\sx)=\chi^{-1}(\det\sx_2)\,G(\sx_2^{-1}\sx_1),\qquad\dS(\sy)=\chi^{-1}(\det\sy_1)\,G(-\sy_2\sy_1^{-1})\]
are partners for $\chi$ over $\sR$ with respect to $\se$.
\end{lemma}
\begin{proof}
If $\dL(\sx)\ne0$ then $\sx_2$ is invertible and $\sx=\rowk{\se}{\sk}$ for the inverse $\sk$ of the block matrix
with rows $(1\ 0)$ and $(\sx_1\ \sx_2)$; if $\dS(\sy)\ne0$ then $\sy_1$ is invertible and $\sy=\colk{\se}{\sk}$ for
the block matrix $\sk$ with columns $\sy$ and $\binom{0}{1}$. For $\sk\in\GL_{\bn}(\sR)$ write
$\rowk{\se}{\sk}=(\sx_1\mid\sx_2)$ and $\colk{\se}{\sk}=\binom{\sy_1}{\sy_2}$. Jacobi's identity
(Input~2, with the rows $(1\ 0)$) gives $\det\sx_2\cdot\det\sk=\det\sy_1$, and
$\rowk{\se}{\sk}\colk{\se}{\sk}=0$ gives $\sx_1\sy_1+\sx_2\sy_2=0$. So $\sx_2$ is invertible exactly when $\sy_1$
is; then $\chi^{-1}(\det\sx_2)=\chi(\det\sk)\chi^{-1}(\det\sy_1)$ and $\sx_2^{-1}\sx_1=-\sy_2\sy_1^{-1}$, and
otherwise both values vanish.
\end{proof}
\leanline{\lean{partnerblock}}

\begin{lemma}[partners are matched]\label{l:matched}
Let $\ba$, $\bb$, $\bn$ be finite index sets, $\sF$ a semilocal field with subring $\sO$, $\chi$ a multiplicative
character of $\sF$ that is tame, $\se\colon\bb\sqcup\ba\simeq\bn$, and let $\dL$, $\dS$ be partners for $\chi$ over
$\sO$ with respect to $\se$ such that $\dS$ is locally constant. Then $(\dL,\dS)\in\matchingPairs{\chi}{\se}$.
\end{lemma}
\begin{proof}
$\dS$ vanishes off the primitive matrices, and so does $\dL\circ{}^{\mathsf T}$; $\dL$ is locally constant because
$\dS$ is and $\GL_{\bn}(\sO)$ is compact. A locally constant function supported on the primitive matrices is a
Schwartz--Bruhat function supported on matrices with a left inverse (Input~1). Let
$\sk\in\GL_{\bn}(\sF)$. By Lemma~\ref{l:section} there are $\sg$ and $\sh$ with $\colk{\se}{\sk}\sg$ and
$\rowk{\se}{\sk}^{\mathsf T}\sh$ primitive, and the two local sections of the matching identity are
$\chi(\det\sg)\dS(\colk{\se}{\sk}\sg)$ and $\chi(\det\sh)\dL(\sh^{\mathsf T}\rowk{\se}{\sk})$. Write
$\colk{\se}{\sk}\sg=\colk{\se}{\mathsf{j}_1}$ with $\mathsf{j}_1\in\GL_{\bn}(\sO)$. Then $\sk^{-1}\mathsf{j}_1$ is
block triangular for $\se$ with diagonal blocks $\sg$ and some $\mathsf{f}^{-1}$, so
$\rowk{\se}{\sk}=\mathsf{f}^{-1}\rowk{\se}{\mathsf{j}_1}$ and $\det\sk\det\sg=\det\mathsf{j}_1\det\mathsf{f}$;
as $\rowk{\se}{\mathsf{j}_1}^{\mathsf T}$ is primitive, $\sh^{\mathsf T}\mathsf{f}^{-1}$ lies in $\GL_{\ba}(\sO)$.
Hence $\sh^{\mathsf T}\rowk{\se}{\sk}=\rowk{\se}{\mathsf{j}}$ for some $\mathsf{j}\in\GL_{\bn}(\sO)$ with
$\colk{\se}{\mathsf{j}}=\colk{\se}{\mathsf{j}_1}$ and
$\det\mathsf{j}\cdot\det(\sh^{\mathsf T}\mathsf{f}^{-1})=\det\mathsf{j}_1$ (Input~2). The partner
identity at $\mathsf{j}$ now gives
$\chi(\det\sh)\dL(\rowk{\se}{\mathsf{j}})=\chi(\det\sh\det\mathsf{j})\dS(\colk{\se}{\mathsf{j}_1})
=\chi(\det\sk)\chi(\det\sg)\dS(\colk{\se}{\sk}\sg)$.
\end{proof}
\leanline{\lean{memmatchingPairs}}

\begin{lemma}[the standard row datum in an integral frame]\label{l:rowstd}
Let $\sF$ be a semilocal field, $\bb$, $\bm$, $\bn$ finite index sets, $\chi$ a multiplicative character of $\sF$
that is tame, $\fA$ an integral frame over $\sF$ whose matrix $(\fA_1\mid\fA_2\mid\fA_3)$ has rows indexed by $\bn$ and
column blocks indexed by $\bb$, $\bm$, $\bb$, and $\se'\colon\bb\sqcup(\bb\sqcup\bm)\simeq\bn$ an identification.
There is a complex number $\rc$ with $|\rc|=1$ and
\[\str{\se'}{\chi}{\fA_3'}=\mathrm{st}_{\fA}\bigl(\sy\mapsto\rc\cdot\inu{\chi}(\det\sy)\bigr).\]
\end{lemma}
\begin{proof}
Multiply the square matrix with the rows of $\fA_3'$ and of $\sx$ by the integral invertible matrix
$(\fA_3\mid\fA_1\mid\fA_2)$. By the relations $\fA_3'\fA_3=1$, $\fA_3'\fA_1=0$, $\fA_3'\fA_2=0$ the product is block triangular with diagonal
blocks $1$ and $\sx\,(\fA_1\mid\fA_2)$; the first matrix is integral exactly when $\sx$ is, and
its determinant is a unit exactly when $\det(\sx(\fA_1\mid\fA_2))$ is. So the two functions agree up to the factor
$\chi$ of the unit $\det(\fA_3\mid\fA_1\mid\fA_2)$, reordered by $\se'$, which has absolute value~$1$.
\end{proof}
\leanline{\lean{existsstdRowDatumeqstdDatum}}

\begin{setting}[a stable space and a partner]\label{set:stable}
$\sF$ is a semilocal field with subring $\sO$ and element $\varpi$; $\ba$, $\bb$, $\bm$, $\bn$ are finite index
sets; $\chi$ is a multiplicative character of $\sF$ that is tame; $\se\colon\bb\sqcup\ba\simeq\bn$ is an
identification, $\mathcal{U}$ a set of $\bn\times\bb$ matrices, and $\cP$ a space of complex functions on
$\Mat{\bn}{\bb}(\sF)$. The hypotheses are: \textup{(i)} $\cP$ is stable under translation: $G(\sg\,\cdot\,)\in\cP$
for all $G\in\cP$ and $\sg\in\GL_{\bn}(\sF)$; \textup{(ii)} $\cP\subseteq\cS(\mathcal{U})$; \textup{(iii)} $\dS$ is
a locally constant function on $\Mat{\bn}{\bb}(\sF)$; \textup{(iv)} $\dL$ and $\dS$ are partners for $\chi$ over
$\sO$ with respect to $\se$.
\end{setting}

\begin{lemma}[a stable space that contains a multiple of a standard datum]\label{l:stable}
In Setting~\ref{set:stable}, let $G\in\cP$, let $\sx\in\Mat{\bn}{\bb}(\sF)$ be integral and
$\ssx\in\Mat{\bb}{\bn}(\sF)$ integral with $\ssx\,\sx=1$, and let $\mathrm{n}$ be a function on
$\Mat{\bn}{\bb}(\sF)$ with values in the natural numbers such that $G(\sy)=\mathrm{n}(\sy)\cdot\stc{\chi}{\ssx}(\sy)$
for all $\sy$ and $\mathrm{n}(\sx)\ne0$. Then $\dS\in\cP$.
\end{lemma}
\begin{proof}
The partner $\dS$ is locally constant and vanishes off the compact set of primitive matrices, so there is an
exponent $i$ with $\dS(\sy+\varpi^{i}\sz)=\dS(\sy)$ for all $\sy$ and all integral $\sz$. Put
$\mathsf{d}=\varpi^{-i}\bigl(1-\sx(1-\varpi^{i})\ssx\bigr)$. It is invertible, with integral inverse
$\varpi^{i}+(1-\varpi^{i})\sx\ssx$, and $\ssx\mathsf{d}=\ssx$, $\mathsf{d}\sx=\sx$. The function
\[G'(\sy)=\int_{\GL_{\bn}(\sO)}\dS(\sk^{-1}\sx)\,G(\mathsf{d}\,\sk\,\sy)\,d\sk\]
lies in $\cP$: the integrand is a locally constant family of elements of $\cP$ over a compact group, by (i) and
(ii), so the integral is a finite combination of its members (unit 105). We show $G'=\mathrm{K}_0\dS$ with
$\mathrm{K}_0=\int|G(\mathsf{d}\sk\sx)|\,d\sk$, which is positive because the integrand is continuous,
non-negative and equal to $\mathrm{n}(\sx)\ne0$ at $\sk=1$. If $G(\mathsf{d}\sk\sy)\ne0$ then $\mathsf{d}\sk\sy$ is
integral and $\ssx\sk\sy=\ssx\mathsf{d}\sk\sy$ has a unit determinant, so $\sy$ is primitive; hence both sides
vanish at a $\sy$ that is not primitive. Let $\sy=\colk{\se}{\mathsf{l}}$ and $\sx=\colk{\se}{\mathsf{j}}$ with
$\mathsf{l},\mathsf{j}\in\GL_{\bn}(\sO)$ (Input~2). Replacing $\sk$ by $\sk\,\mathsf{j}\mathsf{l}^{-1}$,
\[G'(\sy)=\int\dS(\mathsf{l}\mathsf{j}^{-1}\sk^{-1}\sx)\,G(\mathsf{d}\sk\sx)\,d\sk .\]
Where the integrand is non-zero, $\sg=\ssx\sk\sx$ lies in $\GL_{\bb}(\sO)$ and
$\sk\sx=\sx\sg+\varpi^{i}\sz$ with $\sz$ integral, so $\mathsf{l}\mathsf{j}^{-1}\sk^{-1}\sx=\sy\sg^{-1}$ up to
$\varpi^{i}$ times an integral matrix. By the choice of $i$ and the law $\dS(\sy\sg)\chi(\det\sg)=\dS(\sy)$ of a
partner (unit 112), $\dS(\mathsf{l}\mathsf{j}^{-1}\sk^{-1}\sx)=\chi(\det\sg)\dS(\sy)$, while
$G(\mathsf{d}\sk\sx)=\mathrm{n}(\mathsf{d}\sk\sx)\chi^{-1}(\det\sg)$. The product is
$\dS(\sy)\,|G(\mathsf{d}\sk\sx)|$, since $|\chi|=1$ on units. So $G'=\mathrm{K}_0\dS$ and $\dS\in\cP$.
\end{proof}
\leanline{\lean{partnermemofstdColDatum}}

\begin{proposition}[a criterion by unimodular indices]\label{p:unimod}
In Setting~\ref{set:stable}, let $\fA$ be an integral frame over $\sF$ whose matrix $(\fA_1\mid\fA_2\mid\fA_3)$ has rows
indexed by $\bn$ and column blocks indexed by $\bb$, $\bm$, $\bb$, let $\mathrm{e}_{\sF}$ be an additive character
of $\sF$ that has conductor $\sO$, and let $\cM$ be a space of complex functions on $\Mat{\bn}{\bb}(\sF)$. For a
$\bb\times\bb$ matrix $\sz$ over $\sF$ put
\[S_{\sz}(\sy)=\stc{\chi}{\fA_1'}\bigl(\sy+\fA_3\,\sz\,\fA_1'\sy\bigr),\qquad\text{so that }S_0=\stc{\chi}{\fA_1'}.\]
Suppose that every $S_{\sz}$ lies in $\cM$, and that $\cP$ contains every $\dS'\in\cM$ with the following property:
$\cL(\dS')=0$ for every $\bb\times\bb$ matrix $T$ over $\sF$ and every linear functional $\cL$ such that
$\cL(S_{\sz})=\mathrm{e}_{\sF}(\tr(T\sz))\,\cL(S_0)$ for all $\sz$, and such that $\cL(S_0)=0$ if $T$ is primitive.
Then $\dS\in\cP$.
\end{proposition}
\begin{proof}
Let $\bar{\sR}$, the reduction $\sz\mapsto\bar{\sz}$ and $\bar\chi$ be as in the definition of a tame character, and
let $\mathcal{Z}$ be a set of integral $\bb\times\bb$ matrices, containing $1$, that maps bijectively to
$\Mat{\bb}{\bb}(\bar{\sR})$. Since $\mathrm{e}_{\sF}$ has conductor $\sO$,
$\bar{\mathrm{e}}(\bar{\sz})=\mathrm{e}_{\sF}(\varpi^{-1}\sz)$ is a primitive additive character of $\bar{\sR}$ (Input~1).
Consider the Gauss sum of shears
\[\Gamma_0=\sum_{\sz\in\mathcal{Z}}\bar\chi(\det\bar{\sz})\,S_{\varpi^{-1}\sz}\in\cM .\]
\emph{$\Gamma_0\in\cP$.} Let $T$ and $\cL$ be as in the hypothesis, and suppose $\cL(\Gamma_0)\ne0$. Then
$\cL(S_0)\ne0$, and $T$ is integral: for an integral $\ssx$ and a matrix unit $\mathsf{u}$,
$S_{\sz+\ssx\mathsf{u}}=S_{\sz}$, because the standard datum does not see an integral shear; hence
$\mathrm{e}_{\sF}(\tr(T\ssx\mathsf{u}))=1$ for all integral $\ssx$, and the conductor gives the entries of $T$ in
$\sO$. Now $\cL(S_{\varpi^{-1}\sz})=\mathrm{e}_{\sF}(\tr(T\varpi^{-1}\sz))\,\cL(S_0)=\bar{\mathrm{e}}(\tr(\bar T\bar{\sz}))\,\cL(S_0)$, so
$\cL(\Gamma_0)=\bigl(\sum_{\sz\in\mathcal{Z}}\bar\chi(\det\bar{\sz})\,\bar{\mathrm{e}}(\tr(\bar T\bar{\sz}))\bigr)\cL(S_0)$,
and by the matrix Gauss sum (Input~1) the sum is $\bar\chi^{-1}(\det\bar T)$ times its value at $1$. It is
non-zero, so $\det T$ is a unit and $T$ is primitive; but then $\cL(S_0)=0$, a contradiction. So every such $\cL$
kills $\Gamma_0$, and $\Gamma_0\in\cP$ by hypothesis.

\emph{Conclusion.} Apply Lemma~\ref{l:stable} to the translate $G=\Gamma_0(-\varpi^{-1}\,\cdot\,)\in\cP$, to
$\sx=\fA_3-\varpi\fA_1$ and to $\ssx=\fA_3'$; both are integral, and the frame relations (Input~1) give
$\fA_3'\sx=\fA_3'\fA_3-\varpi\fA_3'\fA_1=1$. It remains to show that $G=\mathrm{n}\cdot\stc{\chi}{\fA_3'}$ with a
function $\mathrm{n}$ with values in $\{0,1\}$ and $\mathrm{n}(\sx)=1$. Fix $\sy$ and write
$\mathsf{u}=\fA_1'\sy$, $\mathsf{v}=\fA_2'\sy$, $\mathsf{w}=\fA_3'\sy$. The term of $\sz\in\mathcal{Z}$ in $G(\sy)$ is
\[\bar\chi(\det\bar{\sz})\,\stc{\chi}{\fA_1'}\bigl(-\varpi^{-1}\sy-\varpi^{-2}\fA_3\,\sz\,\mathsf{u}\bigr).\]
A matrix is integral exactly when its three products with $\fA_1'$, $\fA_2'$, $\fA_3'$ are, the frame being
integral with integral inverse. By the frame relations these products of the argument are
$-\varpi^{-1}\mathsf{u}$, $-\varpi^{-1}\mathsf{v}$ and $-\varpi^{-1}(\mathsf{w}+\varpi^{-1}\sz\mathsf{u})$. So the term is
non-zero only if $\mathsf{u}=\varpi\mathsf{u}_1$ with $\mathsf{u}_1\in\GL_{\bb}(\sO)$, $\mathsf{v}$ is $\varpi$ times an
integral matrix, and $\mathsf{w}\equiv-\sz\mathsf{u}_1$ modulo $\varpi$; then $\mathsf{w}$ and $\sy$ are integral,
$\bar{\sz}=-\bar{\mathsf{w}}\,\bar{\mathsf{u}}_1^{-1}$, and the term equals
\[\inu{\chi}\bigl(\det(-\mathsf{w}\mathsf{u}_1^{-1})\bigr)\,\inu{\chi}\bigl(\det(-\mathsf{u}_1)\bigr)=\inu{\chi}(\det\mathsf{w})=\stc{\chi}{\fA_3'}(\sy),\]
because $\bar\chi(\det\bar{\sz})=\inu{\chi}(\det\sz)$ depends only on the reduction of $\sz$, and $\inu{\chi}$ is
multiplicative on the integral elements. At most one $\sz\in\mathcal{Z}$ satisfies the
congruence. So $G(\sy)=\mathrm{n}(\sy)\,\stc{\chi}{\fA_3'}(\sy)$, where $\mathrm{n}(\sy)\in\{0,1\}$ is the number of
non-zero terms. For $\sy=\sx$ one has $\mathsf{u}_1=-1$, $\mathsf{v}=0$, $\mathsf{w}=1$, and the element
$\sz=1$ of $\mathcal{Z}$ contributes: $\mathrm{n}(\sx)=1$.
\end{proof}
\leanline{\lean{partnermemofunimodularvanishing}}

\sect{The difference of the two theta integrals}{s:lift}
\begin{definition}[the theta integral on the small side]\label{d:thsm}
Let $E$ be a CM field with a CM type and Weil data, $\sW$ and $\sV$ hermitian spaces over $E$ with index sets
$\bn$ and $\ba$, $\fc$ a split tame character, $\fv$ a real datum for $(\sW,\sV)$ and $\dS$ a complex function on
$\Mat{\bn}{\ba}(\sF_{\fc})$. Then
\[\thetaSmall{\sW}{\sV}{\fc}{\fv}{\dS}
=\thI{\sW}{\sV}{\omega_{\fc}\upsilon^{\lfloor|\bn|/2\rfloor}}{\upsilon^{|\ba|}}{\loc{\fv}{[\dS\circ{}^{\mathsf T}]^{c}}},\]
a function on $\U(\sW)(\A)$: the theta integral over $[\U(\sV)]$ whose finite datum is the function
$\sy\mapsto\dS(\sy^{\mathsf T})$ of the components at the places conjugate to those of $\fc$.
\end{definition}
\leanline{\lean{thetaSmall}\sep the integral on the large side is the declaration {\leanfont thetaLarge} of unit 43}

\begin{proposition}[the difference is an equivariant lift]\label{p:lift}
Let $E$ be a CM field with a CM type and Weil data; let $\sW$, $\sV$, $\sV'$ be hermitian spaces over $E$ with
index sets $\bn$, $\ba$, $\bb$; let $\fc$ be a split tame character and $\se\colon\bb\sqcup\ba\simeq\bn$ an
identification. Suppose that $|\bn|$ is even, let $\fv$ and $\fv'$ be rapid real data for $(\sW,\sV)$ and
$(\sW,\sV')$, and let $\rc$ be a complex number. There are a space $\cM$ of complex functions on
$\Mat{\bn}{\bb}(\sF_{\fc})$ and a linear map $\cD$ from such functions to functions on $\U(\sW)(\A)$ such that
$(\cM,\cD)$ is an equivariant lift for $\fc$, $\sW$ and the inverse of the local character of $\upsilon^{|\bb|}$ in
rank $|\bb|$, and such that for all partners $\dL$, $\dS$ for $\chi_{\fc}$ over the integers of $\sF_{\fc}$ with
respect to $\se$, with $\dS$ locally constant:
\[\dS\in\cM\qquad\text{and}\qquad
\cD(\dS)=\thL{\sW}{\sV}{\fc}{\fv}{\dL}-\rc\cdot\thetaSmall{\sW}{\sV'}{\fc}{\fv'}{\dS}.\]
\end{proposition}
\begin{proof}
Both theta integrals extend to linear functionals of the datum (Input~3). Write $\chi^{\circ}=\chi^{(|\bn|)}_{\fc}$
for the local character of $\omega^{\mathrm{bc}}_{\fc}$ in rank $|\bn|$ (§\ref{s:setting}); it is tame. The functional
$\cL\colon\dL\mapsto\thL{\sW}{\sV}{\fc}{\fv}{\dL}$ has the invariance required in Proposition~\ref{p:factor} for the
character $\chi^{\circ}$, and not for $\chi_{\fc}$: write $\sg\in\GL_{\ba}(\sF_{\fc})$ as the components of a
finite adelic point $\sk$ of $\U(\sV)$ supported at $\Sigma(\fc)$. Moving the datum by $\sk^{-1}$ on the side of
$\sV$ does not change an integral over $[\U(\sV)]$, and the Weil representation of $\sk^{-1}$ sends the local
datum of $\dL$ to $\chi^{\circ}(\det\sg)$ times the local datum of $\dL(\sg\,\cdot\,)$ (Input~3; here $|\bn|$ even
makes the character of the integral the base change of $\omega_{\fc}^{-1}$, and the factor $|\det\sg|^{|\bn|/2}$ of
the Weil representation is the reason for the rank $|\bn|$). So
$\chi^{\circ}(\det\sg)\,\cL(\dL(\sg\,\cdot\,))=\cL(\dL)$, which is the hypothesis of Proposition~\ref{p:factor} with
$\chi=\chi^{\circ}$, and there is a linear $\cL'$ with $\cL'(\dS)=\thL{\sW}{\sV}{\fc}{\fv}{\dL}$ for every pair
$(\dL,\dS)\in\matchingPairs{\chi^{\circ}}{\se}$. Let $\cM$ be the space of second components of the pairs in
$\matchingPairs{\chi^{\circ}}{\se}$, a subspace of $\cS^{\mathrm{inj}}$, and
$\cD=\cL'-\rc\cdot\thetaSmall{\sW}{\sV'}{\fc}{\fv'}{\,\cdot\,}$. Partners for $\chi_{\fc}$ over the integers of
$\sF_{\fc}$ are partners for $\chi^{\circ}$ over the integers, because the definition of partners over the integers
involves the character only at units and the two characters agree there; so, when $\dS$ is locally constant, they
lie in $\matchingPairs{\chi^{\circ}}{\se}$ (Lemma~\ref{l:matched} for $\chi^{\circ}$), which gives the last assertion. For $\dS\in\cM$, $\cD(\dS)$ is a difference of theta
integrals of target-bounded data, hence automorphic (Input~3). For a point $\sg$ supported at
$\Sigma(\fc)$ the translated pair lies in $\matchingPairs{\chi^{\circ}}{\se}$ (Lemma~\ref{l:translate}), translation of
a theta integral is the theta integral of the moved datum, and the Weil representation of $\sg$ acts on both local
data by the translation of the function and a character of $\det\sg_{\Sigma}$; with $|\bn|=|\ba|+|\bb|$ the two
characters agree and equal the one in the statement, the inverse of the local character of the Hecke character
$\upsilon^{|\bb|}$ in rank $|\bb|$. For $\sk\in\mathcal{K}_f$ trivial at $\Sigma(\fc)$ both local data are fixed.
\end{proof}
\leanline{\lean{existsisEquivariantLift}}

\begin{proof}[Proof of Theorem 2]
Let $\fA$ be the frame of $\fS$ adapted to the places of $\fc$, as in the statement, and $\fA^{\mathrm{op}}$ the
adapted frame of the splitting opposite to $\fS$. The opposite splitting exchanges the two isotropic frames, so
the first and the third blocks change places: the block $(\fA^{\mathrm{op}})_1'$ is the block $\fA_3'$ of $\fA$, and
the standard datum $S_0$ of Proposition~\ref{p:unimod} for the frame $\fA^{\mathrm{op}}$ is
$\dS_0=\stc{\chi_{\fc}}{\fA_3'}$. Put $\cP=\cM\cap\ker\cD$. It is stable under translation by
$\GL_{\bm}(\sF_{\fc})$ (Input~3) and lies in $\cS^{\mathrm{inj}}$. The standard data $\str{\se'}{\chi_{\fc}}{\sz}$
and $\stc{\chi_{\fc}}{\sz}$ are partners and $\stc{\chi_{\fc}}{\sz}$ is locally constant. So
Proposition~\ref{p:unimod}, for the frame $\fA^{\mathrm{op}}$ and the additive character of $\sF_{\fc}$, shows
$\stc{\chi_{\fc}}{\sz}\in\cP$, which is the assertion, once its two hypotheses are verified. The shears $S_{\sz}$ of
$\dS_0$ lie in $\cM$ by Input~3 and the hypothesis $\dS_0\in\cM$. It remains to take $\dS'\in\cM$ that is killed by
every admissible functional and to show $\cD(\dS')=0$.

The function $\cD(\dS')$ is automorphic. By weak approximation (Input~4) it suffices to show that
it vanishes at the points that are trivial at $\Sigma(\fc)$, and a continuous function on the adelic hermitian
matrices that is periodic under the rational ones vanishes when all its Fourier coefficients do. So it suffices
that every centre coefficient $\cf{\fS}{T'}{\cD(\dS')}$, $T'$ a $\bb\times\bb$ matrix over $E$, vanishes at such
points. For fixed $T'$ and a fixed point, this coefficient is a linear functional $\cL$ of $\dS'\in\cM$ which
transforms under the shears $S_{\sz}$ by $\mathrm{e}_{\sF}(\tr(T'_{\Sigma}\sz))$, $T'_\Sigma$ being the matrix of the
components of $T'$ at $\Sigma(\fc)$ (Input~3). So $\cL$ is admissible, and kills $\dS'$, provided
$\cL(\dS_0)=0$ whenever $T'_{\Sigma}$ is primitive.

Let $T'_\Sigma$ be primitive. Then $\det T'\ne0$, and $T'$ may be taken skew-hermitian, the other coefficients
being zero; write $T'=\jmath T$ with $T$ hermitian and invertible, so that the coefficient is the one at
$\jmath T$ of the statement. Freitag's lemma (Input~4), applied to $\cD(\dS_0)$ with a base-adapted stage of the
composed splitting of $\sX$ over $\sW$, reduces the vanishing of the coefficient to its vanishing at the points
$\operatorname{rad}_{\fS}(\mathsf{c},\mathsf{w})\operatorname{lev}_{\fS}(\sq,\sg)$ for which the hermitian matrix built
from $T$ and $\sq$, divided by the inverse different, is unimodular over $\hat{\mathcal{O}}_E$. For such a point
there is $\sk\in\GL_{\bb}(\hat{\mathcal{O}}_E)$ making this matrix the identity (Input~2; this is where the Weil
data must contain the hypothesis that $E$ is stably unramified). If $T$ is represented by the signs of $\sV'$ at
every infinite place, it is the Gram matrix of a space with these signs and a point, and Landherr's theorem
(Input~4) gives a rational change of coordinates $\mathsf{c}_0$ with
$\mathsf{c}_0^{*}\mathrm{J}_{\sV'}\mathsf{c}_0=T$. Moving the rational Levi element of $\mathsf{c}_0$ through the
coefficient turns it into a coefficient at $\jmath\mathrm{J}_{\sV'}$ at a point as in alternative~(i), with
$Y=\mathrm{x}_{\sV'}^{-1}\mathsf{c}_0\sq$ up to the scale idele. If $T$ is not represented, alternative~(ii)
applies. In both cases the coefficient vanishes by hypothesis.
\end{proof}

\sect{Fourier coefficients of theta integrals}{s:coeff}
\begin{definition}[the theta functional on real data]\label{d:threal}
Let $E$ be a CM field with a CM type and Weil data, $\sW_1$ and $\sV$ hermitian spaces with index sets $\bn$ and
$\ba$, $\omega$ and $\omega'$ characters of $\U(1)(\A)$, $F$ a finite datum for $(\sW_1,\sV)$, $\sx$ an adelic
$\ba\times\bn$ matrix and $\mathsf{p}$ a point of $\U(\sW_1)(\A)\times\U(\sV)(\A)$. Then
\[\thetaOnRealData{\sW_1}{\sV}{\omega}{\omega'}{F}{\sx}{\mathsf{p}}\colon\quad
\fv\longmapsto\theta\bigl(\operatorname{heis}(\sx)\,\omega_{\sW_1,\sV}(\omega,\omega')(\mathsf{p})(\fv\otimes F)\bigr)\]
is a continuous linear functional on the real data of the pair; $\operatorname{heis}(\sx)$ is the operator of
the Heisenberg group attached to the vector of $\sx$.
\end{definition}
\leanline{\lean{thetaOnRealData}}

\begin{proposition}[coefficients at indices that are not negative]\label{p:nonneg}
Let $E$ be a CM field with a CM type and Weil data, with element $\jmath$ and splitting character $\upsilon$. Let
$\sX$, $\sY$, $\sZ$ be hermitian spaces with index sets $\bm$, $\bn$, $\ba$, and $\fS$ a splitting of $\sX$ over
$\sZ$ along $\bb$. Let $\omega$ be a character of $\U(1)(\A)$ with weights $\mu$, $\delta$ a character, $f$ a
datum of the pair $(\sX,\sY)$, $\sg$ a point of $\U(\sX)(\A)$ and $T$ a $\bb\times\bb$ matrix over $E$. Assume:
there is an index $i$ of $\sY$ with positive sign; $\mu(\tau)$ is smaller than the number of positive signs of
$\sZ$ for every infinite place $\tau$; $|\bb|+\negc(\sZ)=\negc(\sY)$; $\delta=\omega\,\upsilon^{|\bb|}$; $T$ is
invertible; $T$, as an adelic matrix, is not represented by the signs $-1$, that is, at some infinite place it is
not of the form $-\mathsf{c}^{*}\mathsf{c}$; and the theta terms of $\fS$ for $(\sY,\delta,\delta,f)$ at $\sg$ are integrable. Then
\[\cf{\fS}{\jmath T}{\thI{\sX}{\sY}{\delta}{\delta}{f}}(\sg)=0 .\]
\end{proposition}
\begin{proof}
The coefficient is the sum, over the rational $\bn\times\bb$ matrices $\mathsf{v}$ with
$\mathsf{v}^{*}\mathrm{J}_{\sY}\mathsf{v}=T$, of integrals of theta terms (Input~3); if there is no such
$\mathsf{v}$ it is zero. Fix one, and let $\mathsf{v}'=\mathrm{x}_{\sY}^{-1}\mathsf{v}$ be its coordinates in the
adelic frame. As $T$ is invertible, the coefficient is the integral over the quotient of $\U(\sY)(\A)$ by the
stabilizer $\operatorname{Stab}$ of $\mathsf{v}'$ of the theta term at $\mathsf{v}'$ (Input~3). Let $\tau$ be an
infinite place where $T_\tau$ is not of the form $-\mathsf{c}^{*}\mathsf{c}$, and $\mathsf{w}_0$ the component of $\mathsf{v}'$ at
$\tau$. The hermitian matrix $\mathsf{w}_0^{*}\operatorname{diag}(\sY)\mathsf{w}_0$ is not negative semidefinite, so
some vector $\mathsf{a}_0$ has $\mathsf{a}_0^{*}\mathsf{w}_0^{*}\operatorname{diag}(\sY)\mathsf{w}_0\mathsf{a}_0=1$;
by Witt's theorem (Input~2) an element $\mathsf{u}_0$ of the real unitary group carries
$\mathsf{w}_0\mathsf{a}_0$ to the basis vector of the index $i$.

\emph{No invariant functional.} Let $e_0\ge0$ be the integer with (number of positive signs of
$\sZ)=\mu(\tau)+e_0+1$. We claim that a continuous linear functional $\cL_1$ on the Fock model of the pair
$(\sZ,\sY)$ at the infinite places vanishes if it is invariant under the real Weil representation, of weights $\mu$
and $0$, of all elements of the maximal compact subgroup at $\tau$ that fix $\mathsf{u}_0\mathsf{w}_0$. It suffices
to show $\cL_1(\operatorname{heis}(\mathsf{y}_0)\operatorname{vac})=0$ for every $\mathsf{y}_0$, the Heisenberg
translates of the vacuum being total. Counting dimensions with $|\bb|+\negc(\sZ)=\negc(\sY)$ and the positive index
$i$, one finds a unit row vector $\mathsf{r}_0$ supported on the negative indices of $\sY$, orthogonal to the
columns of $\mathsf{u}_0\mathsf{w}_0$ and to the negative coordinates of $\mathsf{y}_0$ at $\tau$. For a complex
number $\varkappa$ with $|\varkappa|=1$ the rotation
$\mathsf{c}_{\varkappa}=1+(\bar\varkappa-1)\bar{\mathsf{r}}_0\mathsf{r}_0^{\mathsf T}$ of the line of $\mathsf{r}_0$
lies in the maximal compact subgroup and fixes $\mathsf{u}_0\mathsf{w}_0$. It moves
$\operatorname{heis}(\mathsf{y}_0)\operatorname{vac}$ to $\varkappa^{e_0+1}$ times the translate of the vacuum by
$\mathsf{y}_0+(\bar\varkappa-1)\mathsf{z}_0$, $\mathsf{z}_0$ the part of $\mathsf{y}_0$ on the line of
$\mathsf{r}_0$: the vacuum is multiplied by the minor $\det\mathsf{c}_{\varkappa}=\bar\varkappa$ to the power given
by the weight, which is $\varkappa^{e_0+1}$. So $\cL_1(\operatorname{heis}(\mathsf{y}_0)\operatorname{vac})$ is the
mean over the circle of a function $\varkappa\mapsto\varkappa^{e_0+1}G_0(\varkappa)$ with $G_0$ holomorphic on the
closed disc (the translate depends holomorphically on $\varkappa$, as $\mathsf{z}_0$ lies in the holomorphic
coordinates), and by the mean value property this mean is the value at $\varkappa=0$, which is $0$ because
$e_0+1\ge1$.

\emph{Averaging.} Let $\mathcal{C}$ be the intersection of $\operatorname{Stab}$ with the conjugate by
$\mathsf{u}_0^{-1}$, at $\tau$, of the maximal compact subgroup; it is compact. For $\sk\in\operatorname{Stab}$ the
theta term at $\mathsf{v}'$ of the datum moved by $\sk$ is the theta term moved through the pair $(\sZ,\sY)$ (law
of the Levi elements, Input~3, with $\delta=\omega\upsilon^{|\bb|}$). The average of the theta term over
$\mathcal{C}$, for its Haar probability measure, is therefore a functional of the kind just excluded, applied to
the datum: it is zero. Right multiplication by $\mathcal{C}$ preserves the Haar measure of $\operatorname{Stab}$.
The averaging criterion (Input~5: if averaging multiplies the integrand by a number different from the modulus,
the integral over the quotient vanishes) gives the result.
\end{proof}
\leanline{\lean{centreCoeffeqzeroofnotnegative}}

\begin{proposition}[theta terms of standard data]\label{p:terms}
Let $E$ be a CM field with a CM type and Weil data; let $\sX$, $\sW_1$, $\sV$ be hermitian spaces with index sets
$\bm$, $\bn$, $\ba$, and $\fS$ a base-adapted splitting of $\sX$ over $\sW_1$ along $\bb$.
Let $\Sigma$ be a finite set of finite places whose product of completions $\sF_\Sigma$ is not the zero ring, with
a reading of the components at $\Sigma$, and let $\fA$ be the adapted frame of $\fS$. Let $\omega'$ be a
character, $\chi$ a multiplicative character of $\sF_\Sigma$, $\rc_0$ a complex number and $\mathsf{n}$ a
$(\bb\sqcup\bn)\times\ba$ matrix over $\sF_\Sigma$; let $\dL$ be a function with $\dL\in\cS(\mathcal{U})$ for a
set $\mathcal{U}$ and $\dL=\mathrm{st}_{\fA}\bigl(\sy\mapsto\rc_0\,\inu{\chi}(\det(\sy\mathsf{n}))\bigr)$; let
$\fv$ be a rapid real datum; let $\omega$, $\omega''$ be characters with $\omega''=\omega\,\upsilon^{|\bb|}$; and
let $\mathsf{c}$, $\mathsf{w}$, $\sq$, $\sg$ be the coordinates of a point. Put $f=\loc{\fv}{[\dL]}$ and
$\mathsf{p}=\operatorname{rad}_{\fS}(\mathsf{c},\mathsf{w})\,\operatorname{lev}_{\fS}(\sq,\sg)$. Then:
\begin{enumerate}[label=\textup{(\arabic*)},leftmargin=2.3em,itemsep=1pt]
\item the theta terms of $\fS$ for $(\sV,\omega'',\omega',f)$ at $\mathsf{p}$ are integrable;
\item for every rational $\ba\times\bb$ matrix $\mathsf{v}$, with $\mathsf{v}'=\mathrm{x}_{\sV}^{-1}\mathsf{v}$,
every adelic $\ba\times\bb$ matrix $Y$ with $\mathsf{v}'\sq=\mathrm{sc}_EY$ and every point $\sk$ of $\U(\sV)(\A)$,
with $\sy=\sk_f^{-1}Y_f$, the theta term at $\mathsf{v}$ and $\sk$ equals
\begin{multline*}
\mathrm{e}_{\A}\bigl(\tfrac12\tr(\mathsf{w}\,\mathsf{v}^{*}\mathrm{J}_{\sV}\mathsf{v})\bigr)\,\omega'_{\ba}(\det\sq)^{-1}\,|\det\sq|^{|\ba|/2}\,
[\sy\text{ integral}]\cdot\operatorname{vac}_{\fS}(\sV)\\
\cdot\ \thetaOnRealData{\sW_1}{\sV}{\omega}{\omega'}{F_{\sy}}{\mathsf{v}'\mathsf{c}^{*}\operatorname{diag}(\sW_1)}{(\sg,\sk)}\bigl(\mathrm{sl}_{\fS}(\sk^{-1}Y;\fv)\bigr),
\end{multline*}
where $\omega'_{\ba}$ is the character attached to $\omega'$ in rank $|\ba|$, $\operatorname{vac}_{\fS}(\sV)$ the
vacuum factor of the mixed model, $\operatorname{diag}(\sW_1)$ the diagonal form of the signs of $\sW_1$, and
$F_{\sy}$ the finite datum named at $\Sigma$ by the function
$\sy'\mapsto\rc_0\,\inu{\chi}(\det((\bar{\sy}\mid\sy')\mathsf{n}))$ of the integral $\sy'$, $\bar{\sy}$ being the
components of $\sy$ at $\Sigma$;
\item for all signs $\mathsf{t}'$ on $\bb$ and every $\bb\times\bb$ matrix $T$ over $E$ that, as an adelic matrix, is not represented
by $\mathsf{t}'$: if every non-zero real slice $\mathrm{sl}_{\fS}(\sx;\fv)$ has a Gram matrix
$\sx^{*}\operatorname{diag}(\sV)\sx$ that is represented by $\mathsf{t}'$, then
$\cf{\fS}{\jmath T}{\thI{\sX}{\sV}{\omega''}{\omega'}{f}}(\mathsf{p})=0$.
\end{enumerate}
\end{proposition}
\begin{proof}
(1) The datum is named at a set of places disjoint from its conjugate by a Schwartz--Bruhat function, and $\fv$ is
rapid; integrability then holds once the matrices at which $\dL$ does not vanish give left-invertible transposed
coordinate matrices at $\mathsf{p}$ (Input~3). Where $\dL\ne0$ the matrix $\sy(\fA_1\mid\fA_2)$ of the standard
datum has a unit determinant, and the coordinate matrix of $\mathsf{p}$ on the frame of $\fS$ is the adelic frame
times a block triangular matrix with diagonal blocks $\sq$ and an invertible one; this gives the left inverse.

(2) Move the datum through the two factors of $\mathsf{p}$. The radical element contributes the character
$\mathrm{e}_{\A}(\frac12\tr(\mathsf{w}\,\mathsf{v}^{*}\mathrm{J}_{\sV}\mathsf{v}))$ and the Heisenberg operator of $\mathsf{v}'\mathsf{c}^{*}\operatorname{diag}(\sW_1)$; the Levi element
contributes $\omega'_{\ba}(\det\sq)^{-1}|\det\sq|^{|\ba|/2}$ and replaces $\mathsf{v}'$ by $\mathsf{v}'\sq$ (Input~3, with
$\omega''=\omega\upsilon^{|\bb|}$). What remains is the source term of the local datum on the adapted frame: the
vacuum factor times, when the finite part is integral, the theta functional of the vector named by the slice of
$\dL$, applied to the real slice of $\fv$ (Input~3). The slice of the standard datum $\dL$ at $\bar{\sy}$
is the function named in the statement.

(3) The coefficient is the sum over the rational $\mathsf{v}$ with $\mathsf{v}^{*}\mathrm{J}_{\sV}\mathsf{v}=T$ of integrals of theta terms. By (2)
each term is a multiple of the functional of Definition~\ref{d:threal} at the real slice
$\mathrm{sl}_{\fS}(\sk^{-1}Y;\fv)$. If this slice were non-zero, its Gram matrix would be represented by $\mathsf{t}'$;
but that Gram matrix is, up to the invertible factor $\sq$ and the unitary part of the point, the image of
$\mathsf{v}^{*}\mathrm{J}_{\sV}\mathsf{v}=T$, which is not represented by $\mathsf{t}'$. So every term vanishes.
\end{proof}
\leanline{\lean{thetaTermstdDatumspec}}

\begin{lemma}[products of data]\label{l:box}
Let $E$ be a CM field with a CM type and Weil data, $\sW_1$ and $\sV'$ hermitian spaces with index sets $\bn$ and
$\bb$. Let $\omega$, $\omega'$ be characters of $\U(1)(\A)$; $F'$ and $F$ finite data for $(\sW_1,\sV')$ and
$(\sW_1,\sW_1)$; $\sx$ an adelic $\bb\times\bn$ matrix; $\sg$, $\mathsf{u}$, $\mathsf{l}$ points of
$\U(\sW_1)(\A)$, $\U(\sV')(\A)$, $\U(\sW_1)(\A)$; and $\fv'$, $\fv$ real data for $(\sW_1,\sV')$ and
$(\sW_1,\sW_1)$. Then
\begin{multline*}
\thetaOnRealData{\sW_1}{\sV'\oplus\sW_1}{\omega}{\omega\omega'}{F'\boxtimes F}{\tbinom{\sx}{0}}{(\sg,\mathsf{u}\oplus\mathsf{l})}(\fv'\boxtimes\fv)\\
=\thetaOnRealData{\sW_1}{\sV'}{\omega}{\omega'}{F'}{\sx}{(\sg,\mathsf{u})}(\fv')\cdot
\theta_{\sW_1,\sW_1}(\omega,\omega;\fv\otimes F)(\sg,\mathsf{l}).
\end{multline*}
\end{lemma}
\begin{proof}
The model of the pair $(\sW_1,\sV'\oplus\sW_1)$ is the tensor product of the models of $(\sW_1,\sV')$ and
$(\sW_1,\sW_1)$; the Weil representation of a block-diagonal point, the Heisenberg operator of a vector in the
first summand and the theta functional are compatible with this product. Unfolding the definitions gives the
identity.
\end{proof}
\leanline{\lean{thetaOnRealDataboxProdeqmul}}

\begin{proposition}[the equal-rank theta integral is constant]\label{p:const}
Let $E$ be a CM field with a CM type and Weil data, $\sX$ and $\sW$ hermitian spaces with index sets $\bm$ and
$\ba$, $\fS_1$ a base-adapted splitting of $\sX$ over $\sW$ along one line, and $\fc$ a split tame character.
Suppose that $|\bm|$ is even, and let $\fv$ be a rapid real datum for $(\sX,\sX)$. Then for every point $\sg$ of
$\U(\sX)(\A)$
\[\thL{\sX}{\sX}{\fc}{\fv}{\stc{\chi_{\fc}}{1}}(\sg)=\thL{\sX}{\sX}{\fc}{\fv}{\stc{\chi_{\fc}}{1}}(1).\]
\end{proposition}
\begin{proof}
Apply Proposition~\ref{p:lift} to $(\sX,\sX)$ and the zero space, with constant $0$: the partners of the empty
frame are $\stc{\chi_{\fc}}{1}$ on both sides, so the function in question is $\cD(\stc{\chi_{\fc}}{1})$ for an
equivariant lift. It is therefore automorphic, has level supported at $\Sigma(\fc)$ and is invariant under $\mathcal{K}_f$
off $\Sigma(\fc)$. By the tower property (Input~5) its constant term along the isotropic line of
$\fS_1$ is a theta integral of the smaller pair, so that on the Levi subgroup it depends only on the component in
$\U(\sW)$. The constancy criterion from strong approximation (Input~5) applies.
\end{proof}
\leanline{\lean{thetaLargeeqapplyoneofevenrank}}

\begin{definition}[the equal-rank ratio; the standard exchange identity]\label{d:exch}
Let $E$ be a CM field with a CM type and Weil data.

\textup{(a)} For a hermitian space $\sV$ with index set $\ba$, a split tame character $\fc$ and a real datum $\fv$
for $(\sV,\sV)$,
\[\equalRankThetaRatio{\sV}{\fc}{\fv}=\frac{\bigl|\thL{\sV}{\sV}{\fc}{\fv}{\stc{\chi_{\fc}}{1}}(1)\bigr|^2}{\LP{\fc}{\lfloor|\ba|/2\rfloor}} .\]

\textup{(b)} Let $\sX$, $\sV$, $\sV'$ be hermitian spaces with index sets $\bm$, $\bn$, $\bb$, let $\mu$ be
weights, $\se\colon\bb\sqcup\bn\simeq\bm$ an identification, $\fv$ and $\mathfrak{y}$ real data for $(\sX,\sV)$ and
$(\sX,\sV')$, and $\cN$ a real function on the split tame characters of weight $\mu$. The statement
\[\StdExchange{\sX}{\sV}{\sV'}{\mu}{\se}{\fv}{\mathfrak{y}}{\cN}\]
holds if $\fv$ and $\mathfrak{y}$ are rapid and there is a real number $\rC$ such that every split tame character
$\fc$ of weight $\mu$ has a complex number $\rc$ with
$\rC\,|\rc|^2=\LP{\fc}{|\bn|-\lfloor|\bm|/2\rfloor}\cdot\cN(\fc)$ and
\[\thL{\sX}{\sV}{\fc}{\fv}{\str{\se}{\chi_{\fc}}{\sz}}=\rc\cdot\thetaSmall{\sX}{\sV'}{\fc}{\mathfrak{y}}{\stc{\chi_{\fc}}{\sz}}
\qquad\text{for every integral }\sz\in\Mat{\bb}{\bm}(\sF_{\fc}).\]
\end{definition}
\leanline{\lean{equalRankThetaRatio}\sep \lean{StdExchange}}

\sect{The real pair, and the proofs of Theorems 3 and 1}{s:real}
\begin{proof}[Proof of Theorem 3]
Both data are rapid by assumption. \emph{The constant.} The inclusion of $\U(\sV')\times\U(\sW_1)$ in
$\U(\sV'\oplus\sW_1)$ is a closed embedding. Unfolding in stages (Input~5), with the weight that is
$\Omega$ at the infinite places on the points whose finite part lies in $\mathcal{K}_f$ up to a block
element, and $0$ on the other points, gives a number $\mathrm{R}\ne0$. Let $\mathsf{a}=\det(\mathrm{sc}_E\,\mathrm{x}_{\sV'})$,
an idele, and
$\mathrm{k}=\operatorname{vac}_{\fS}(\sV'\oplus\sW_1)\,\operatorname{vac}_{\fS}(\sV')^{-1}\,|\mathsf{a}|^{|\bn|/2}$.
Take $\rC=|\mathrm{R}\,\mathrm{k}|^{-2}$. Fix a split tame character $\fc$ of weight $\mu$, and let $\fA$ be the
frame of $\fS$ adapted to the places of $\fc$. By Lemma~\ref{l:rowstd} there is $\rc_1$ with $|\rc_1|=1$ and
$\str{\se'}{\chi_{\fc}}{\fA_3'}=\mathrm{st}_{\fA}(\sy\mapsto\rc_1\inu{\chi_{\fc}}(\det\sy))$. Let $(\cM,\cD)$ be the
equivariant lift of Proposition~\ref{p:lift} for $\sX$, $\sV'\oplus\sW_1$, $\sV'$, the data $\fv$, $\fv'$ and the
constant
\[\rc=\mathrm{R}\,\mathrm{k}\,\rc_1\,\omega^{\mathrm{bc}}_{\fc}(\mathsf{a})\,\mathrm{l}(\fc).\]
Then $\rC|\rc|^2=|\mathrm{l}(\fc)|^2$, and this is $\LP{\fc}{|\bb|+|\bn|-|\bm|/2}$ times
$|\mathrm{l}(\fc)|^2/\LP{\fc}{|\bn|/2}$ because $|\bm|=|\bn|+2|\bb|$ and the products of $L$-values are positive.
For an integral $\sz$ the row and column data of $\sz$ are partners, so the asserted identity is
$\cD(\stc{\chi_{\fc}}{\sz})=0$. By Theorem~2 it suffices that $\dS_0=\stc{\chi_{\fc}}{\fA_3'}$ lies in $\cM$, which
holds because it is the partner of a row datum, and that the centre coefficients of
$\cD(\dS_0)$, the integral on the large side of the partner of $\dS_0$ minus $\rc$ times the integral on the small side of $\dS_0$, vanish in
the two cases of Theorem~2.

\emph{Case \textup{(i)}: the index $\mathrm{J}_{\sV'}$.} By Proposition~\ref{p:terms} the theta terms of both
pairs are integrable and explicit. The coefficient of the large integral is the integral over the quotient by the
stabilizer of the rational frame $\binom{1}{0}$ of $\sV'$ in $\sV'\oplus\sW_1$, which is $\U(\sW_1)$; that of the
small integral is the integral of the theta term at $1$ over $\U(\sV')(\A)$, the stabilizer being trivial
(Input~3). By unfolding in stages the first is $\mathrm{R}$ times the number $\rc'$ for which the inner integral
over $[\U(\sW_1)]$ of the large theta term at a point $\sh$, shifted by a point of $\U(\sV')\times\U(\sW_1)$ that
may depend on $\sh$, is $\Omega(\sh)\,\rc'$ times the small theta term; the shift of Input~5 is the point
$\mathsf{p}$ of the period condition. If the finite part of the point lies in $\mathcal{K}_f$ up to a block element, insert the explicit terms of
Proposition~\ref{p:terms}\,(2): for the large pair with the constant $\rc_1$ and the matrix $1$, for the small one
with the frame $\binom10$; the standard row datum is fixed by the integral point, the named vector of the large
pair is the product of the vector named by the indicator of the integral matrices and of $F_{\fc}$, and the
period condition, used with $B=\sy^{-1}Y$, evaluates the integral over $[\U(\sW_1)]$. The characters are collected
as follows: $\omega(\det B)$ splits into its infinite and finite parts; the quotient of the values of
$\inu{\chi_{\fc}}$ at $\det\sx$ and at its conjugate is the value of $\omega_{\fc}$ at the finite part (unit 38);
$|\det\sy|=1$ for unitary $\sy$. The result is $\rc'=\mathrm{k}\,\rc_1\,\omega^{\mathrm{bc}}_{\fc}(\mathsf{a})\,\mathrm{l}(\fc)$.
If the finite part does not lie in $\mathcal{K}_f$ up to a block element, the large theta term vanishes: were the
finite matrix integral, Witt's theorem over $\hat{\mathcal{O}}_E$ (Input~2) would produce such a block element. So
the coefficient of $\cD(\dS_0)$ is $\mathrm{R}\rc'-\rc=0$ times the small coefficient.

\emph{Case \textup{(ii)}: $T$ invertible and not represented by the signs of $\sV'$.} The small coefficient
vanishes by Proposition~\ref{p:terms}\,(3), the slices of the pair $(\sX,\sV')$ having Gram matrices represented
by the signs of $\sV'$. The large one vanishes by Proposition~\ref{p:terms}\,(3) under the first alternative of
the support or sign condition, and by Proposition~\ref{p:nonneg} under the second: $\sV'$ is negative definite,
$\sW_1$ has a positive sign, and the bound on the weights follows from $2\negc(\sW_1)\le|\bn|$ and the positivity
of the weights of $\fc$.
\end{proof}

\begin{lemma}[a datum with prescribed slices]\label{l:slice}
Let $E$ be a CM field with a CM type and Weil data; let $\sX$, $\sW_1$, $\sV'$ be hermitian spaces with index sets
$\bm$, $\bn$, $\ba$, and $\fS$ a base-adapted splitting of $\sX$ over $\sW_1$ along $\bb$. Let $\Omega_0$ be a real
function on the space of real coordinates of the slices that is smooth, has compact support, is non-negative and
is non-zero at a point $\sy$, and let $\fv$ be a rapid real datum for $(\sW_1,\sV')$. Then there are a complex
number $\mathrm{k}$ and a rapid real datum $\fx'$ for $(\sX,\sV')$ such that the slice of $\fx'$ at $0$ for the
splitting opposite to $\fS$ is $\fv$, and the slice of $\fx'$ at every $\sx$, in the mixed model of $\fS$, is
$\Omega_0(\sx)\,\mathrm{k}\cdot\omega_\infty(1)\fv$.
\end{lemma}
\begin{proof}
The integral of $\Omega_0$ is positive; let $\mathrm{k}$ be its inverse. By the existence of sections in the mixed
model (Input~3: a smooth compactly supported coefficient times a fixed rapid vector) there is a rapid $\fx'$ with
the prescribed slices. The slice at $0$ for the opposite splitting is the integral over $\sx$ of the slices
(Input~3), which is $\bigl(\int\Omega_0\bigr)\mathrm{k}\,\fv=\fv$; two real data are equal when the integrated data
built from their zero terms agree.
\end{proof}
\leanline{\lean{existsrapidwithrealSlice}}

\begin{lemma}[Witt's theorem, smooth local form]\label{l:witt}
Let $\bb$, $\bm$ be finite index sets, $\mathrm{J}$ a complex $\bm\times\bm$ matrix with $\mathrm{J}^{*}=\mathrm{J}$,
and $\sx$ a complex $\bm\times\bb$ matrix with $\det(\sx^{*}\mathrm{J}\sx)\ne0$. There is a map $\sk$ from the
$\bm\times\bb$ matrices to the unitary group of $\mathrm{J}$ such that $\sy\mapsto\sk(\sy)$ is smooth as a map of
real vector spaces and, for every $\sy$ in a neighbourhood of $\sx$, $\sk(\sy)\,\sx\,\mathsf{c}=\sy$ for some
$\bb\times\bb$ matrix $\mathsf{c}$.
\end{lemma}
\begin{proof}
Put $\mathsf{w}=(\sx^{*}\mathrm{J}\sx)^{-1}\sx^{*}$ and, for an $\bm\times\bm$ matrix $\sg$,
$\mathsf{S}_0(\sg)=(\sg\sx\mathsf{w}-\sx\mathsf{w}\sg^{*})\mathrm{J}$, which is skew for $\mathrm{J}$, so that
$\exp\mathsf{S}_0(\sg)$ is unitary. The map $\sg\mapsto\exp(\mathsf{S}_0(\sg))(\sg-\mathsf{S}_0(\sg))$ has derivative the identity at
$\sg=1$, so it has a smooth local inverse $\mathsf{A}$ near $1$ (and a smooth function defined everywhere agrees
with it near $1$, unit 112). With $\mathsf{p}(\sy)=1-(\sx-\sy)\mathsf{w}\mathrm{J}$ one has
$\mathsf{p}(\sy)\sx=\sy$ and $\mathsf{p}(\sx)=1$; put $\sk(\sy)=\exp\mathsf{S}_0(\mathsf{A}(\mathsf{p}(\sy)))$. Near $\sx$,
$\sk(\sy)\bigl(\mathsf{A}\mathsf{p}(\sy)-\mathsf{S}_0(\mathsf{A}\mathsf{p}(\sy))\bigr)\sx=\mathsf{p}(\sy)\sx=\sy$, and
$\bigl(\sg-\mathsf{S}_0(\sg)\bigr)\sx$ lies in the column space of $\sx$, since $\sg\sx\mathsf{w}\mathrm{J}\sx=\sg\sx$.
\end{proof}
\leanline{\lean{existscontDiffwitt}}

\begin{proof}[Proof of Theorem 1]
Let $\mathrm{J}$ be the diagonal form of $\sV'\oplus\sW_1$ at a real place and $\sx_0=\binom10$ the frame of
$\sV'$, so that $\sx_0^{*}\mathrm{J}\sx_0$ is the diagonal form of $\sV'$. Lemma~\ref{l:witt} gives a smooth
$\sk$ with $\sk(\sy)\sx_0\mathsf{c}=\sy$ near $\sx_0$. Choose a smooth bump $\Omega_1\ge0$ at $\sx_0$ (in the
coordinates at every infinite place) with compact support in the neighbourhood where this holds and where the
Gram matrix $\sx^{*}\mathrm{J}\sx$ is invertible, and put
$\mathsf{f}_0(\sx)=\sx\,(\sx^{*}\mathrm{J}\sx)^{-1}(\sx^{*}\mathrm{J}\sx_0)$, so $\mathsf{f}_0(\sx_0)=\sx_0$. Where
$\Omega_1(\mathsf{f}_0(\sx))\ne0$ one has $\sx=\sk(\mathsf{f}_0(\sx))\,\sx_0\,\mathsf{c}$ with
$\mathsf{c}^{*}\operatorname{diag}(\sV')\mathsf{c}=\sx^{*}\mathrm{J}\sx$.

\emph{The small datum.} Lemma~\ref{l:slice}, with the weight
$\sx\mapsto\Omega_1(\sx_0\sx)\,\bigl|\prod_\tau\det(\sx_\tau)^{|\bn|/2+\mu(\tau)}\bigr|^2$ (smooth, non-negative,
compactly supported, non-zero at $1$), gives $\mathrm{k}$ and a rapid $\fx'$ whose slice at $0$ for the opposite
splitting is $\fv$.

\emph{The large datum.} Let $\fx$ be a rapid real datum for $(\sX,\sV'\oplus\sW_1)$ whose slice at $\sx$ is,
with $\sy=\mathsf{f}_0(\sx)$,
\[\Omega_1(\sy)\,\Omega_1\bigl(\sk(\sy)^{-1}\sx\bigr)\,\prod_\tau\overline{\det\bigl((\sk(\sy)^{-1}\sx)_{\tau,1}\bigr)}^{\,2\mu(\tau)}\,\mathrm{k}\cdot
\omega_\infty\bigl(1,\sk(\sy)\bigr)(\fv\boxtimes\fv_1),\]
where $(\,\cdot\,)_{\tau,1}$ is the first block at $\tau$ and $\omega_\infty$ the real Weil representation of weights
$|\bn|/2-\mu$. It exists by the existence of sections (Input~3): the coefficient is smooth because $\sk$ is, it
has compact support because its support lies in the image of a product of two compact sets, and the weighted
norms of the vector field are bounded by those of the vacuum, of $\fv_1$ and of $\fv$.

\emph{The period condition.} Apply Theorem~3 to $\fx$, $\fx'$, the weight
$\Omega(\sh)=\Omega_1(\mathsf{f}_0(\sh^{-1}\sx_0))$, which is non-negative and non-zero near $1$, and
$\mathrm{l}(\fc)=\thL{\sW_1}{\sW_1}{\fc}{\fv_1}{\stc{\chi_{\fc}}{1}}(1)$. The support condition holds in its first
form by the relation $\mathsf{c}^{*}\operatorname{diag}(\sV')\mathsf{c}=\sx^{*}\mathrm{J}\sx$. For the period
condition let $\sy=\mathsf{f}_0(\sh^{-1}\sx_0)$. If $\Omega_1(\sy)=0$ both sides vanish. Otherwise there are
$\mathsf{c}_\tau$ preserving the form of $\sV'$ with $\sk(\sy_\tau)\sx_0\mathsf{c}_\tau=(\sh^{-1}\sx_0)_\tau$, hence
elements $\mathsf{u}_\tau$ of the real unitary group of $\sV'$; then $\sh_\tau\sk(\sy_\tau)(\mathsf{u}_\tau\oplus1)$
fixes $\sx_0$, so it is $1\oplus\mathsf{q}_\tau$ with $\mathsf{q}_\tau$ in the real unitary group of $\sW_1$. Use
the point $\mathsf{p}=(\mathsf{u},\mathsf{q}^{-1})$. At the slice argument $\sx=\sh^{-1}\sx_0\mathsf{u}^{-1}B$ one
checks $\mathsf{f}_0(\sx)=\sy$, so the slices of $\fx$ and $\fx'$ are those written above. The Weil representation
of the point on $\fv\boxtimes\fv_1$ factors by Lemma~\ref{l:box}; the theta kernel of $(\sW_1,\sW_1)$ with
$\fv_1$, integrated over $[\U(\sW_1)]$, is the theta integral on the large side, which is constant
(Proposition~\ref{p:const}, by $\fS_1$ and $|\bn|$ even) and so equals $\mathrm{l}(\fc)$; the character of $\fc$
at the infinite places is read from its weights (unit 96), which produces $\omega^{\mathrm{bc}}_{\fc}(\det_\infty B)$ and the
power of $|\det_\infty B|$. Clearing denominators gives the period condition. So Theorem~3 yields the standard
exchange identity with the function $\fc\mapsto|\mathrm{l}(\fc)|^2/\LP{\fc}{|\bn|/2}$, which is
$\equalRankThetaRatio{\sW_1}{\fc}{\fv_1}$ by Definition~\ref{d:exch}.
\end{proof}

\sect{Determinant vectors and the cone integral}{s:det}
The tower property (Input~5) produces a \emph{cone integral} of the real datum. One floor down in the tower, the
second note needs to know that the real data of a special shape are reproduced by it.

\begin{definition}[determinant vectors]\label{d:detvec}
Let $E$ be a CM field with a CM type and Weil data, $\sW$ and $\sV$ hermitian spaces with index sets $\bn$ and
$\br$ and signs $\operatorname{sgn}_{\sW}$ and $\operatorname{sgn}_{\sV}$. Let $F$ be an adelic $\bn\times\br$
matrix and $\mathsf{u}$ a point of the real coordinate space of the Fock model of the pair $(\sW,\sV)$
(coordinates indexed by an infinite place and a pair of indices).
\begin{enumerate}[label=\textup{(\alph*)},leftmargin=2.3em,itemsep=1pt]
\item $\posCoordMatrix{\operatorname{sgn}_{\sW}}{\operatorname{sgn}_{\sV}}{F}{\mathsf{u}}$ is the family, over the
infinite places, of the square matrices $\mathsf{u}^{+}F$, where $\mathsf{u}^{+}$ is the matrix of the coordinates
of the projection of $\mathsf{u}$ to the coordinates that join basis vectors of equal sign.
\item For a natural number $e$, natural numbers $\mu(\tau)$ at the infinite places and $\rc\in\CC$,
\[\detVectorFn{\operatorname{sgn}_{\sW}}{\operatorname{sgn}_{\sV}}{e}{\mu}{\rc}{F}{\mathsf{u}}
=\rc\cdot\operatorname{vac}(\mathsf{u})\cdot\overline{\prod_\tau\det\nolimits_{+}(\mathsf{M}_\tau)^{\mu(\tau)}\,\det\nolimits_{-}(\mathsf{M}_\tau)^{\mu(\tau)+2e}},\]
where $\mathsf{M}=\posCoordMatrix{\operatorname{sgn}_{\sW}}{\operatorname{sgn}_{\sV}}{F}{\mathsf{u}}$, $\det_{\pm}$ are the
principal minors on the indices of $\br$ of sign $\pm1$, and $\operatorname{vac}$ is the Gaussian, the vacuum
vector.
\item $\detVectors{e}{\mu}$ is the set of the vectors of the Fock model of $(\sW,\sV)$ that are rapid and equal to
$\mathsf{u}\mapsto\detVectorFn{\operatorname{sgn}_{\sW}}{\operatorname{sgn}_{\sV}}{e}{\mu}{\rc}{F}{\mathsf{u}}$ for
some $\rc$ and $F$. The pair $(\sW,\sV)$ is an implicit argument; we name it when needed.
\end{enumerate}
\end{definition}
\leanline{\lean{posCoordMatrix}\sep \lean{detVectorFn}\sep \lean{detVectors}}

\begin{lemma}[the real frame of an adapted splitting]\label{l:frame}
\textup{(a)} Let $E$ be a CM field with a CM type, $\sW$ and $\sW'$ hermitian spaces with index sets $\bn$ and
$\bn'$, and $\fS$ a splitting of $\sW$ over $\sW'$ along $\br$ that is adapted at the point $1$ with signs
$\mathsf{t}$. For every infinite place $\tau$ the real frame matrix $\mathsf{P}$ of $\fS$ at $\tau$ is unitary, and
$\operatorname{diag}(\sW)\,\mathsf{P}=\mathsf{P}\,\mathsf{B}$, where $\mathsf{B}$ has the block
$\bigl(\begin{smallmatrix}0&1\\1&0\end{smallmatrix}\bigr)$ on $\br\sqcup\br$ and the diagonal form of $\sW'$ on
$\bn'$.

\textup{(b)} For a finite index set $\bn$ with signs, the projection onto the coordinates of positive sign is
$\frac12(1+\operatorname{diag})$ of the signs.
\end{lemma}
\begin{proof}
(a) Adaptedness gives $\mathsf{P}^{*}\mathsf{P}=1$, and $\mathsf{P}$ is square, so $\mathsf{P}\mathsf{P}^{*}=1$.
The Gram matrix of the frame of a splitting is $\mathsf{B}$ (unit 38), that is,
$\mathsf{P}^{*}\operatorname{diag}(\sW)\mathsf{P}=\mathsf{B}$; multiply by $\mathsf{P}$. (b) By definition.
\end{proof}
\leanline{\lean{frameAtrealunitary}\sep \lean{posProjeq}}

\begin{lemma}[slices of polynomial vectors]\label{l:posslice}
Let $E$ be a CM field with a CM type and Weil data, $\sW$ and $\sW'$ hermitian spaces with index sets $\bn$ and
$\bn'$, and $\fS$ a splitting of $\sW$ over $\sW'$ along $\br$ that is adapted at $1$ with signs $\mathsf{t}$; let
$\mathsf{L}$, $\mathsf{I}$, $\mathsf{K}$ be the blocks of its mixed model and $\mathrm{P}_{+}$ the projections to
the positive coordinates. For all coordinate vectors $\mathsf{a}$ and $\mathsf{u}$ there is $\mathsf{U}$ with
$\mathrm{P}_{+}\mathsf{U}=\mathsf{L}(-\mathsf{a})+\mathsf{I}\,\mathrm{P}_{+}\mathsf{u}+\mathsf{K}(-\mathsf{a})$
such that, for every polynomial in the positive coordinates, with Fock vector $\fv$: the slice of $\fv$ at
$\mathsf{a}$, evaluated at $\mathsf{u}$, is $\fv(\mathsf{U})$.
\end{lemma}
\begin{proof}
Put $\mathsf{J}_0(\mathsf{a},\mathsf{z},\mathsf{b})=\mathsf{L}\mathsf{a}+\mathsf{I}\mathsf{z}+\mathsf{K}\mathsf{b}$ and
$\mathsf{U}=\mathsf{J}_0(-\mathsf{a},\mathsf{u},-\mathsf{a})$. Lemma~\ref{l:frame} gives
$\mathrm{P}_{+}\mathsf{J}_0(\mathsf{a},\mathsf{z},\mathsf{b})=\mathsf{J}_0(\tfrac12(\mathsf{a}+\mathsf{b}),\mathrm{P}_{+}\mathsf{z},\tfrac12(\mathsf{a}+\mathsf{b}))$
and $\langle\mathsf{J}_0(\mathsf{a},\mathsf{u},\mathsf{a}),\mathsf{J}_0(\mathsf{c},\mathsf{z},\mathsf{c})\rangle=2\langle\mathsf{a},\mathsf{c}\rangle+\langle\mathsf{u},\mathsf{z}\rangle$,
the real frame matrix being orthonormal. Both sides of the asserted identity are continuous and linear in the
polynomial vector, so it suffices to compare them on the Heisenberg translates of the vacuum (unit 37), where the
slice is explicit and both sides are the same Gaussian exponential by the two formulas.
\end{proof}
\leanline{\lean{existssliceeq}}

\begin{proposition}[the cone integral is onto the determinant vectors]\label{p:cone}
Let $E$ be a CM field with a CM type and Weil data; let $\sW$, $\sW'$, $\sV$, $\sV'$ be hermitian spaces with
index sets $\bn$, $\bn'$, $\br$, $\br'$; let $\fA$ be a base-adapted splitting of $\sW'$ over $\sW$ along one line
and $\fS$ a base-adapted splitting of $\sV'$ over $\sV$ along one line. Let $e$ be a natural number, $\mu(\tau)$
natural numbers at the infinite places, $\mathrm{b}$ a parameter of the cone integral and $\mathrm{z}$ a real
number with $1\le\mathrm{z}$, and suppose $|\bn'|=2(e+\negc(\sW'))$. Then the cone integral with weights $e+\mu$,
\[\fv'\longmapsto\iint\operatorname{cone}_{\fA,\fS}(e+\mu,\mathrm{b},\mathrm{z};\fv',\mathsf{p})\,d\mathsf{p}\]
(the integral of the cone integrand of the tower property for the product of two Haar measures) has this
property: every element of the set $\detVectors{e}{\mu}$ of the pair $(\sW,\sV)$ is the cone integral of an element
of the set $\detVectors{e}{\mu}$ of the pair $(\sW',\sV')$. (Nothing is said about the cone integrals of the other
determinant vectors of the larger pair.)
\end{proposition}
\begin{proof}
Let $\fv=\detVectorFn{\operatorname{sgn}_{\sW}}{\operatorname{sgn}_{\sV}}{e}{\mu}{\rc}{F}{\,\cdot\,}$ be a
determinant vector. The line-term expansion of the cone integral (unit 43) provides a splitting, adapted at $1$,
whose mixed terms are the line terms, and its frame in terms of the blocks of $\fA$ and $\fS$. Let $G$ be the
adelic $\bn'\times\br'$ matrix that is $F$ on the smaller spaces plus a multiple of the hyperbolic part, and for
$\rc'\in\CC$ let $\mathrm{R}(\rc')$ be the Fock vector of the polynomial
$\rc'\prod_\tau \Delta_{\tau,+}^{\mu(\tau)}\Delta_{\tau,-}^{\mu(\tau)+2e}$, where $\Delta_{\tau,\pm}$ are determinants of matrices of
linear forms with coefficients from $G$. The real Weil representation of a block-diagonal point acts on a
polynomial vector by powers of principal minors and a linear substitution (unit 43); with the exponents fixed by
$|\bn'|=2(e+\negc(\sW'))$ the minors cancel, and
$\mathrm{R}(\rc')=\detVectorFn{\operatorname{sgn}_{\sW'}}{\operatorname{sgn}_{\sV'}}{e}{\mu}{\rc'}{G}{\,\cdot\,}$;
it is rapid, being a polynomial vector. By Lemma~\ref{l:posslice} the slice of a polynomial vector is its value at
an explicit point, and by Lemma~\ref{l:frame} the matrix
$\posCoordMatrix{\operatorname{sgn}_{\sW'}}{\operatorname{sgn}_{\sV'}}{G}{\,\cdot\,}$ at that point is conjugate,
by the unitary frame matrix, to the block sum of the Levi coordinate $\mathsf{q}_\tau$ on the hyperbolic plane and of
$\posCoordMatrix{\operatorname{sgn}_{\sW}}{\operatorname{sgn}_{\sV}}{F}{\mathsf{u}}$; so each of its principal
minors is $\mathsf{q}_\tau$ times the corresponding minor of the smaller matrix. Therefore the cone integrand of
$\mathrm{R}(\rc')$ at a point $\mathsf{p}$ is $\exp(\mathrm{E}_0(\mathsf{p}))\,\rc'\rc^{-1}\fv$ with the real function
$\mathrm{E}_0(\mathsf{p})=-\sum_\tau|\mathsf{q}_\tau|^2+\log\bigl(|\mathsf{q}|^{\mathrm{z}}\prod_\tau|\mathsf{q}_\tau|^{2(e+\mu(\tau))}\bigr)$. For
$\fv\ne0$ the function $\exp\mathrm{E}_0$ is integrable with positive integral, and
$\rc'=\rc\,(\int\exp\mathrm{E}_0)^{-1}$ gives a determinant vector with cone integral $\fv$; for $\fv=0$ take $0$.
\end{proof}
\leanline{\lean{surjOnconeintegraldetVectors}}

\sect{Two norm formulas}{s:norms}
\begin{setting}[auxiliary pairs and $L$-value data]\label{set:aux}
The following objects are those of units 134, 138, 140 and 141 and of the second note of this unit; we recall
what the statements need.

\textup{(a)} $E$ is a CM field with a CM type and Weil data, $\br$ and $\bn$ are finite index sets. An
\kw{auxiliary pair} $\fa$ for $E$, $\br$, $\bn$ consists of two families of finite places (the chosen places and
the twist places), a hermitian space $\sW(\fa)$ with index set $\br$, an embedding of $\br$ into $\bn$ with
$|\br|<|\bn|$, a weight $\operatorname{wt}(\fa)$ of forms on $\sI_{\bn}$ and a character. It determines: the characters and
the theta kernel of the pair $(\sW(\fa),\sI_{\bn})$; a finite datum $\Upsilon(\fa)$ and, for the highest weight
vector of $\operatorname{wt}(\fa)$, a real datum $\fv(\fa)$; the space $\mathcal{A}(\fa)$ of algebraic forms on $\sI_{\bn}$ of
the level, character and weight of $\fa$, with its inner product $\langle\ ,\ \rangle$; the class-sum lift
$\theta_{\fa}$ from $\mathcal{A}(\fa)$ to forms on $\sW(\fa)$; the residue ring $\sR(\fa)$ of the product of all
places of $\fa$; the level indicator $\mathbb{1}_{\fa}$, a function on $\Mat{\bn}{\br}(\sR(\fa))$; the number
$\mathrm{t}(\fa)$, the product over the twist places of the orders of $\GL_{\br}$ of the residue fields; the level
group of $\fa$ and its index $\mathrm{i}(\fa)$ in the stabilizer of the integral lattice; and the action of the
product of the groups $\GL_{\br}$ of the residue fields of the chosen places on $\mathcal{A}(\fa)$. A parameter
on $\sW(\fa)$ may \kw{fit} $\fa$. A form $F\in\mathcal{A}(\fa)$ \kw{lifts} a family $\mathsf{p}$ of complex
numbers $\mathsf{p}(\fp,n)$, $n$ a natural number, if at every finite place $\fp\ne\bar{\fp}$ off the places of $\fa$ and their conjugates
it is spherical with the Hecke polynomial that Hecke transfer (Input~7) attaches to the polynomial of
$\mathsf{p}$ at $\fp$.

\textup{(b)} $P$ are parameters, $D$ admissible data, $\vartheta$ a tame character of $D$ and $i$ a lift index,
$1\le i\le m$. Then $\sW_i(D)$ is the $i$-th space of the tower of the rows, $\fa_D(\vartheta,i)$ the $i$-th
auxiliary pair, the \kw{forms of principal level} on a space are those of the parameter of the rows fixed by the
principal level at the chosen places, and the \kw{adjoint lift} $\theta^{*}_{D,\vartheta,i}$ is the adjoint of
$\theta_{\fa_D(\vartheta,i)}$, a linear map from the forms of principal level on $\sW_i(D)$ to
$\mathcal{A}(\fa_D(\vartheta,i))$. The \kw{eigenpart} of a space for a family $\mathsf{p}$ consists of the forms
of principal level on which each Hecke operator acts by the value of $\mathsf{p}$. $\Theta^{\circ}_{D,\vartheta}$
is the chain of the $m$ lifts as a map on all forms of principal level on the first space $\sY_0(D)$ of the
chain.

\textup{(c)} \kw{$L$-value data} $\fl$ for $E$ and $\bn$ consist of finitely many chosen places and twist places
(finite places of $E$; $\Sigma(\fl)$ is the set of all of them, and $\overline{\Sigma(\fl)}$ the set of their
conjugates), jet degrees at the infinite places, and a character $\omega_{\fl}$ of the corresponding weight,
unramified off the places, with $|\bn|>1$. They determine an auxiliary pair $\fa(\fl)$ for the one-point index
set and $\bn$, with $\sW(\fa(\fl))=\sI_1$, and the lift $\Xi(\fl)\in\mathcal{A}(\fa(\fl))$ of the trivial character,
a class sum of kernel forms. $\operatorname{hook}(\fl)$ is the hook number of the weight of $\fa(\fl)$ in rank
$|\bn|$, $\operatorname{weyl}(\fl)$ the number of elements of its Weyl index set, and $\operatorname{sg}(\fp)\in\{\pm1\}$ the sign
that is $-1$ for one place of each pair of conjugate split places and $1$ otherwise.
\end{setting}

\begin{lemma}[the adjoint of the class-sum lift is a theta lift]\label{l:adjoint}
In Setting~\ref{set:aux}\,(a), let $\fa$ be an auxiliary pair for $E$, $\br$, $\bn$ with $\br$ non-empty, $l$ a
parameter on $\sW(\fa)$ that fits $\fa$, $y$ a form on $\sW(\fa)$ of the type of $l$, and $F\in\mathcal{A}(\fa)$
such that $\langle F,\Xi\rangle=\pet{y}{\theta_{\fa}(\Xi)}$ for every $\Xi\in\mathcal{A}(\fa)$. Then $F$ is the
algebraic form of the theta lift of $y$ to $\sI_{\bn}$ with the datum $\fv(\fa)\otimes\Upsilon(\fa)$.
\end{lemma}
\begin{proof}
Since $l$ fits $\fa$, the real datum is an edge, so with $\Upsilon(\fa)$ it gives an arrow from the type of $l$
to the type on $\sI_{\bn}$ and, for the swapped pair, an arrow back (unit 134, Input~7). The algebraic form of the
lift of $y$ lies in $\mathcal{A}(\fa)$. For every $\Xi$, $\theta_{\fa}(\Xi)$ is a fixed multiple of the conjugate
of the descent of $\bar\Xi$, and the adjoint of the lift along an arrow is the conjugate of the descent; so the
algebraic form of the lift has the same inner products with every $\Xi$ as $F$.
\end{proof}
\leanline{\lean{adjointthetaeqformLift}}

\begin{lemma}[the datum of the level indicator]\label{l:moddatum}
In Setting~\ref{set:aux}\,(a), let $\fa$ be an auxiliary pair, $\fv$ a real datum, and $y$ a function on
$\U(\sW(\fa))(\A)$ that is automorphic and rapidly decreasing and is fixed by the principal level at the chosen
places. Then the theta lift of $y$ with the datum $\fv\otimes\Upsilon(\fa)$ equals its theta lift with the datum
whose real part is $\fv$ and whose finite part is the function $\mathrm{t}(\fa)\cdot\mathbb{1}_{\fa}$ modulo the
places of $\fa$.
\end{lemma}
\begin{proof}
$\Upsilon(\fa)$ is the sum, over the elements $\mathsf{A}$ of the product of the groups $\GL_{\br}$ of the
residue fields at the twist places, of the translates of the datum of $\mathbb{1}_{\fa}$ by lifts of
$\mathsf{A}^{-1}$ that reduce to $1$ at the chosen places (unit 134). Each such lift lies in the principal level
at the chosen places and fixes $y$; by the equivariance of the theta lift each summand gives the same lift.
\end{proof}
\leanline{\lean{liftauxDatumeqliftmodDatum}}

\begin{lemma}[the level indicator]\label{l:indicator}
Let $E$ be a CM field with a CM type, $\br$ and $\bn$ finite index sets and $\fa$ an auxiliary pair for them. Put
$G=\mathrm{t}(\fa)\cdot\mathbb{1}_{\fa}$. Then the rows of every $\sy$ with $G(\sy)\ne0$ that are indexed by the
image of $\br$ form the identity matrix, and
\[\sum_{\sy}|G(\sy)|^2=\mathrm{t}(\fa)^2\cdot\Bigl(\prod_{\mathfrak{q}}\Np\mathfrak{q}\Bigr)^{(|\bn|-|\br|)\,|\br|},\]
the product over the twist places $\mathfrak{q}$ of $\fa$.
\end{lemma}
\begin{proof}
A matrix over $\sR(\fa)$ is the family of its reductions at the places of $\fa$, and this is a bijection
(unit 104). In these terms $\mathbb{1}_{\fa}$ is the indicator of a product of sets, which are counted place by
place.
\end{proof}
\leanline{\lean{levelIndicatorsupportandnormSq}}

\begin{definition}[the Gram factor and the terms of the inner product]\label{d:gram}
\textup{(a)} In Setting~\ref{set:aux}\,(a), for an auxiliary pair $\fa$,
\[\gramFactor{\fa}=\frac{\mathrm{i}(\fa)\cdot\mathrm{t}(\fa)^2\cdot\bigl(\prod_{\mathfrak{q}}\Np\mathfrak{q}\bigr)^{(|\bn|-|\br|)\,|\br|}}{\#\GL_{\br}(\sR(\fa))},\]
the product over the twist places $\mathfrak{q}$ of $\fa$.

\textup{(b)} In Setting~\ref{set:aux}\,(b), let $\sV$ be a hermitian space with the index set of the rows, $y'$
and $y$ forms of principal level on $\sV$, and $\bar{\sg}$ a coset of the finite adelic unitary group of the
definite space of the auxiliary pairs modulo its elements that are integral off the places of the modulus
$\mathfrak{m}_D(\vartheta,1)$ of $\fa_D(\vartheta,1)$ and their conjugates; here $\mathfrak{m}_D(\vartheta,e)$ is, as
in the other notes, the product of the chosen places to the power $e$ times the places of $\vartheta$. Then
\[\rallisTerm{D}{\vartheta}{i}{\sV}{y'}{y}{\bar{\sg}}=\frac{\mathrm{ct}(\bar{\sg})}{\#\Mat{\bn}{\br}\bigl(\mathcal{O}_E/\mathfrak{m}_D(\vartheta,1)\bigr)},\]
where $\mathrm{ct}(\bar{\sg})$ is the coset term of the inner product formula (Input~6) for the character of
$\fa_D(\vartheta,i)$ on the definite side: the vacuum coefficient at a representative $\sg$ times
$\pet{y'}{R(\sg)y}$; here $\bn$ and $\br$ are the index sets of the definite space and of the rows.
\end{definition}
\leanline{\lean{gramFactor}\sep \lean{rallisTerm}}

\begin{lemma}[the chain of lifts is spherical]\label{l:spherical}
Let $P$ be parameters and $D$ admissible data, and let $\mathsf{p}$ be a family of complex numbers
$\mathsf{p}(\fp,n)$. There is a family $G$ of functions $G(\fp)\colon\CC\to\CC$, one for each finite place, such
that for every form $y$ in the eigenpart of $\sY_0(D)$ for $\mathsf{p}$ and every finite place $\fp\ne\bar{\fp}$ such
that neither $\fp$ nor $\bar{\fp}$ is a chosen place, the form $\Theta^{\circ}_{D,1}(y)$ is spherical at $\fp$ with
Hecke polynomial $G(\fp)$.
\end{lemma}
\begin{proof}
A form is spherical with a given Hecke polynomial exactly when its algebraic form is (unit 102). Induct on the
length of the chain. A form in the eigenpart is spherical at such $\fp$ with the Hecke polynomial of $\mathsf{p}$
(unit 138). Each further lift is, up to a scalar, the theta lift by a datum that is an arrow (unit 141), and Hecke
transfer for theta lifts (Input~7) applies: the characters are unramified at $\fp$ because their tame parts are
trivial off the places of the auxiliary pairs, and the trivial twist has no twist places.
\end{proof}
\leanline{\lean{existsisSphericalWithiota}}

\begin{lemma}[adjointness]\label{l:inner}
In Setting~\ref{set:aux}\,(b), for parameters $P$, data $D$, a tame character $\vartheta$, a form $y$ of
principal level on $\sW_i(D)$ and $\Xi\in\mathcal{A}(\fa_D(\vartheta,i))$:
$\langle\theta^{*}_{D,\vartheta,i}\,y,\Xi\rangle=\pet{y}{\theta_{\fa_D(\vartheta,i)}(\Xi)}$.
\end{lemma}
\begin{proof}
This is the definition of the adjoint lift (unit 138).
\end{proof}
\leanline{\lean{innerauxLift}}

\begin{proof}[Proof of Theorem 4]
The real datum of the auxiliary pair $\fa_D(1,i)$ of the trivial twist is an edge (second note; unit 134), and
$2\operatorname{rk}\sW_i(D)<\operatorname{rk}\sI_{\bn}$. The Euler form of the inner product formula (Input~6)
gives $A_0\ne0$, valid for all characters of the weights of the edge, hence for all $\vartheta$. Put $A=A_0$
times the number of elements of the Weyl index set. Fix $\vartheta$ and write $\fa=\fa_D(\vartheta,i)$,
$G=\mathrm{t}(\fa)\mathbb{1}_{\fa}$. For a form $y$ of principal level, Lemma~\ref{l:inner} and
Lemma~\ref{l:adjoint} show that $\theta^{*}_{D,\vartheta,i}\,y$ is the algebraic form of the lift of $y$ by
$\fv(\fa)\otimes\Upsilon(\fa)$, which by Lemma~\ref{l:moddatum} is the lift by the datum of $G$. This datum is
sharp for the principal level at the places of $\fa$ and their conjugates (Input~6): the modulus has places
(unit 140), the tame part of the character is trivial off them, $G$ has the support property of
Lemma~\ref{l:indicator} and $\sum|G|^2\ne0$. So the lift is injective, and so is $\theta^{*}_{D,\vartheta,i}$.
The inner product of two algebraic forms is $\mathrm{i}(\fa)$ times the number of elements of the Weyl index set
times the Petersson product of the corresponding forms (second note). Multiplying the series of Input~6 by this
factor and inserting the constant of sharpness, with the value of $\sum|G|^2$ from Lemma~\ref{l:indicator}, gives
the terms $A\cdot\gramFactor{\fa}\cdot\rallisTerm{D}{\vartheta}{i}{\sW_i(D)}{y'}{y}{\bar{\sg}}$.
\end{proof}

\begin{lemma}[the adjoint lift transfers eigensystems]\label{l:lifts}
In Setting~\ref{set:aux}\,(b), let $P$ be parameters, $D$ data, $\vartheta$ a tame character, $\mathsf{p}$ a
family of complex numbers and $y$ a form in the eigenpart of $\sW_i(D)$ for $\mathsf{p}$. Then
$\theta^{*}_{D,\vartheta,i}\,y$ lifts $\mathsf{p}$.
\end{lemma}
\begin{proof}
By Lemmas \ref{l:inner} and~\ref{l:adjoint} it is the algebraic form of a theta lift along an arrow; $y$ is
spherical at the places in question with the Hecke polynomial of $\mathsf{p}$ (unit 138), and Hecke transfer
(Input~7) gives the polynomial required, the characters being unramified there.
\end{proof}
\leanline{\lean{liftsauxLift}}

\begin{corollary}[injectivity]\label{c:inj}
In Setting~\ref{set:aux}\,(b), for parameters $P$, data $D$ and a tame character $\vartheta$, the adjoint lift
$\theta^{*}_{D,\vartheta,i}$ is injective.
\end{corollary}
\begin{proof}
It is the first assertion of Theorem~4.
\end{proof}
\leanline{\lean{auxLiftinjective}}

\begin{lemma}[the adjoint lift is equivariant at the chosen places]\label{l:rep}
In Setting~\ref{set:aux}\,(b), let $P$ be parameters, $D$ data, $\vartheta$ a tame character, $\ssx$ an element
of the product of the groups $\GL$ of the residue fields of the chosen places, in the rank of the rows, and $y$ a
form of principal level on $\sW_i(D)$. Then $\theta^{*}_{D,\vartheta,i}(\ssx\cdot y)=\ssx\cdot\theta^{*}_{D,\vartheta,i}(y)$,
for the action of $\ssx$ on forms of principal level by translation and its action on
$\mathcal{A}(\fa_D(\vartheta,i))$.
\end{lemma}
\begin{proof}
Two elements of $\mathcal{A}(\fa)$ are equal when their inner products with all $\ssx\cdot\Xi$ agree, the action
being bijective. By Lemma~\ref{l:inner}, $\langle\theta^{*}(\ssx y),\ssx\Xi\rangle=\pet{\ssx y}{\theta_{\fa}(\ssx\Xi)}$;
$\theta_{\fa}(\ssx\Xi)$ is the translate of $\theta_{\fa}(\Xi)$ by a lift of $\ssx$ (unit 138), and the Petersson
product and the inner product of $\mathcal{A}(\fa)$ are invariant (units 102, 138).
\end{proof}
\leanline{\lean{auxLiftprincipalRep}}

\begin{lemma}[the lift of the trivial character is spherical]\label{l:trivlifts}
In Setting~\ref{set:aux}\,(c), let $E$ be a CM field with a CM type and Weil data, $\bn$ a finite index set and
$\fl$ $L$-value data. Then $\Xi(\fl)$ lifts the constant family $1$ for the auxiliary pair $\fa(\fl)$.
\end{lemma}
\begin{proof}
Hecke transfer for an equivariant map (Input~7) applies to the map from finite data to forms given by the class
sum over $\U(\sI_1)$ of the kernel values, and to the datum $\Upsilon(\fa(\fl))$, whose image is $\Xi(\fl)$. On
$\U(\sI_1)$ the representation is trivial, so the Hecke polynomial on that side is the one of the constant
family~$1$.
\end{proof}
\leanline{\lean{auxPairliftstrivialLift}}

\begin{proof}[Proof of Theorem 5]
Write $\fa=\fa(\fl)$ and $G=\mathrm{t}(\fa)\mathbb{1}_{\fa}$. \emph{Reduction to a doubling integral.} The inner
product of $\Xi(\fl)$ with itself is an index times a Petersson product (second note), and $\Xi(\fl)$ is, up to
the inverse of an index, the algebraic form of the theta lift of the constant function $1$ on $\U(\sI_1)$ with
the datum $\fv(\fa)\otimes\Upsilon(\fa)$, hence with the datum of $G$ (Lemma~\ref{l:moddatum}; the constant
function is automorphic and, $\sI_1$ being definite, rapidly decreasing). The doubling form of the inner product
formula (Input~6; the rank condition is $2<|\bn|$) expresses the Petersson product as
$\operatorname{vol}[\U(\sI_{\bn})]$ times the integral over $\U(\sI_1)(\A)$ of the matrix coefficient of the datum.

\emph{The places off $\fl$.} The datum is the vacuum off the places of $\fl$ and their conjugates. Integrating
out one more place of $E^{+}$ off these, that is, the set $\{\fp,\bar{\fp}\}$ (one place of $E$ if
$\fp=\bar{\fp}$, two if not), multiplies the partial integral by the local doubling sum
\[\prod_{\mathfrak{u}\in\{\fp,\bar{\fp}\}}\frac{1-\operatorname{sg}(\mathfrak{u})\,\Np\mathfrak{u}^{-e}}{1-\omega^{\mathrm{bc}}_{\fl}(\varpi_{\mathfrak{u}})^{-1}\,\Np\mathfrak{u}^{-e}}\]
(Input~6; here $e>1$ and $\omega_{\fl}$ is unramified at these places, and $\omega^{\mathrm{bc}}_{\fl}$ is the base
change of $\omega_{\fl}$); for $\fp=\bar{\fp}$ it is the single factor
$(1-\Np\fp^{-e})/(1-\omega^{\mathrm{bc}}_{\fl}(\varpi_{\fp})^{-1}\Np\fp^{-e})$, the sign being $1$. By induction on
finite sets of places and passage to the limit, the partial integrals converging and the Euler products
converging for $e>1$ (Input~7), the integral is
\[L^{\Sigma(\fl)\cup\overline{\Sigma(\fl)}}\bigl(e,(\omega_{\fl}^{-1})^{\mathrm{bc}}\bigr)\cdot\prod_{\mathfrak{u}\notin\Sigma(\fl)\cup\overline{\Sigma(\fl)}}\bigl(1-\operatorname{sg}(\mathfrak{u})\,\Np\mathfrak{u}^{-e}\bigr)\]
times the partial integral at the places of $\fl$, at their conjugates and at infinity.

\emph{The places of $\fl$ and infinity.} The datum of $G$ is sharp (Input~6, by Lemma~\ref{l:indicator}), so the
remaining partial integral is its constant of sharpness times the archimedean integral, which over the compact
group is $\operatorname{vol}[\U(\sI_1)]$ times the Fock pairing of the highest weight vector with itself; the
polynomial of the highest vector is a renaming of the harmonic polynomial of $\fa$ (second note), the pairing is
invariant under renaming, and its value is the hook number.

\emph{Collecting the factors.} The local factors that enter are the following. At infinity:
$\operatorname{vol}[\U(\sI_1)]\cdot\operatorname{hook}(\fl)$; and the passage from the Petersson product to the inner
product of algebraic forms contributes $\operatorname{weyl}(\fl)$ and an index (second note). At a place $\fp$ of
$\fl$, with residue field of $\Np\fp$ elements: the stabilizer of a line in $\GL_{\bn}$ of the residue field has
index $\sum_{n<|\bn|}\Np\fp^{\,n}$, and the stabilizer of a non-zero vector has index $\Np\fp^{|\bn|}-1$ (units 44,
103); the number of all vectors is $\Np\fp^{|\bn|}$, and the number of units of the residue field is $\Np\fp-1$.
At the twist places, moreover, $\sum|G|^2=\mathrm{t}(\fa)^2\,(\prod_{\mathfrak{q}}\Np\mathfrak{q})^{|\bn|-1}$
(Lemma~\ref{l:indicator} with $|\br|=1$), where $\mathrm{t}(\fa)=\prod_{\mathfrak{q}}(\Np\mathfrak{q}-1)$. These numbers
combine as follows. A place $\fp$ of $\fl$ is not its own conjugate, the sign takes both values on the pair
$\{\fp,\bar{\fp}\}$, and the product of the two sign factors of the pair is
\[\bigl(1-\operatorname{sg}(\fp)\Np\fp^{-e}\bigr)\bigl(1-\operatorname{sg}(\bar{\fp})\Np\fp^{-e}\bigr)=1-\Np\fp^{-|\bn|}
=\frac{\Np\fp^{|\bn|}-1}{\Np\fp^{|\bn|}},\]
the index of the stabilizer of a vector divided by the number of vectors; with these factors the product of the
sign factors off $\fl$ becomes the product over all finite places, the inverse of the Euler product in the
denominator of $\rC$. The quotient of the index of the stabilizer of a vector by that of a line is $\Np\fp-1$,
which leaves $\prod(\Np\fp-1)$ over the chosen places; at the twist places only the power
$(\prod\Np\mathfrak{q})^{|\bn|-1}$ of Lemma~\ref{l:indicator} remains in the result. The formal proof makes this collection in one long step of
simplification, which the readers of the dossier mark as coarse; it is reported here by its local factors and its
result and was not recomputed line by line. Finally $\rC\ne0$ because the volumes and the Euler product are
non-zero.
\end{proof}

\sect{Notes on the formal text}{s:notes}
\begin{enumerate}[label=\textup{\arabic*.},leftmargin=1.6em,itemsep=5pt]
\item \emph{Half ranks.} Where the formal text divides a number of elements by $2$ it uses the division of
natural numbers; in Theorems 1 and~3 and in Propositions \ref{p:lift} and~\ref{p:const} the number is assumed
even, and in Definition~\ref{d:exch} the floor is written. The argument of $\LP{\fc}{\cdot}$ in
Definition~\ref{d:exch}\,(b) is an integer, possibly negative.
\item \emph{Local sections depend on choices.} The functional $\cL_0$ and the representatives of
Lemma~\ref{l:reps} are fixed by the axiom of choice; Definition~\ref{d:locsec} and Definition~\ref{d:match}
depend on them, and the results of §\ref{s:local} hold for every choice.
\item \emph{Hypotheses that are not used, or only formally.} In Lemma~\ref{l:clopen} the identification $\se$ is an
argument of the formal statement although the conclusion does not mention it. In Proposition~\ref{p:nonneg} the
hypothesis on the index $i$ is the existence of a positive sign of $\sY$. Lemma~\ref{l:indicator} is stated in the
formal text without Weil data, the surrounding statements with them.
\item \emph{Statements that begin with a local definition.} The formal statements of Proposition~\ref{p:unimod},
Proposition~\ref{p:terms}, Lemma~\ref{l:indicator}, Lemma~\ref{l:block}, Lemma~\ref{l:posslice} and Theorem~5
begin with a local definition (of $S_{\sz}$, of the datum and the point, of $G$, of the blocks, of the Fock vector,
of $\rC$). The reader's tool takes such a definition for the beginning of the proof; the dossier restores the
statements, and they were read here in the Lean text.
\item \emph{Arguments that only Lean can settle.} For some declarations of this note the dossier gives the list
of arguments by Lean's rules, applied by hand, because the text alone does not determine it; no compiler was run.
They are: Lemma~\ref{l:reps}, Lemma~\ref{l:clopen}, Lemma~\ref{l:block}, Proposition~\ref{p:unimod},
Proposition~\ref{p:terms}, Lemma~\ref{l:indicator}, Theorems 3, 4 and~5, Lemma~\ref{l:trivlifts},
Definition~\ref{d:exch}\,(b) and Definition~\ref{d:detvec}\,(c). The doubt concerns implicit and instance
arguments only (which index sets, signs and instances a declaration takes); the explicit arguments and the
hypotheses are those written in the Lean statements. For instance arguments of definitions the tool marks some
as taken only if the defining term uses them.
\item \emph{A reading in the dossier that the types contradict.} For Lemma~\ref{l:rowstd} the dossier gives the
identification as $\bb\sqcup(\bb\sqcup\bn)\simeq\bm$. The frame of the formal statement has its rows indexed by
$\bn$, so the identification that the types allow is $\bb\sqcup(\bb\sqcup\bm)\simeq\bn$, as stated here; in the
application (Theorem~3) the two index sets appear with exchanged names. This was inferred from the types in the
Lean text, not checked by a compiler.
\item \emph{The second reading.} A second reader compared the statements of the unit with the Lean text. Three
of her points concern this note and are followed here: the constant $\rC$ of Theorem~5 depends only on the field
and the rank, and its $L$-function omits the Euler factors at the places of the datum and at their conjugates
(the dossier allowed a dependence on the chosen places and named the places only); in the period condition of
Theorem~3 the character in the denominator is the base change of $\omega_{\fc}$, evaluated at an idele;
Proposition~\ref{p:cone} asserts preimages and no more; and in the proof of Proposition~\ref{p:factor} the pieces
are cut out by the primitivity of $\sx^{\mathsf T}\mathsf{r}^{-1}$, as written here from the Lean text.
\item \emph{What is taken for granted.} Everything listed as Inputs. In particular the two forms of the inner
product formula, Hecke transfer, the matrix Gauss sum, Freitag's lemma, Landherr's theorem, the tower property and
the constancy criterion are results of the background library. The Inputs give, for each, the relation to the
literature that the second reader of the unit found by comparing the formal statement with the sources. For the
two forms of the inner product formula, for the matrix Gauss sum and for Hecke transfer that relation is
\emph{unresolved}: the sources named there contain an analogue or a part, and not the formal statement. None of
them is called a special case of a published theorem here.
\item \emph{Weil data.} In the formal text, Weil data include three conditions that an earlier edition of this
note left out: $E$ contains a primitive root of unity of order $5040$; $(1-\jmath^{-1})/2$ is integral; and
$\upsilon$ is a Hecke character of $E$ with the value $-1$ at the uniformisers of the places $\fp=\bar{\fp}$. The
first is used in this note: the diagonalisation of unimodular hermitian matrices over $\hat{\mathcal{O}}_E$
(Input~2), in the proof of Theorem~2, is proved in the formal text under it.
\item \emph{Two local characters.} The partners in all statements are partners for $\chi_{\fc}$, the local
character of rank $0$. In the proof of Proposition~\ref{p:lift} the large side is invariant for the local
character $\chi^{(|\bn|)}_{\fc}$ of rank $|\bn|$, and the formal proof factors through the matched pairs of that
character; the two agree on units, which is all that partners over the integers see.
\item \emph{Where the results go.} Theorem~1, Proposition~\ref{p:lift}, Lemma~\ref{l:stable}, Lemma~\ref{l:block},
Lemma~\ref{l:slice}, Lemma~\ref{l:box}, Proposition~\ref{p:const}, Proposition~\ref{p:cone} and Theorem~5 are used
in the second note of the unit; Theorem~4 and Lemmas \ref{l:lifts}, \ref{l:rep} and~\ref{l:spherical} are used by
other units of the argument.
\end{enumerate}

\subhead{How this note was written.} From the dossier of the unit, which three readers made in three parts; this
note follows the first part. The program check of the whole dossier covered \UnitSteps{} leaf steps of
\UnitProofs{} proofs and found every statement at its lines. I read all statements of this part in the dossier,
all structured proofs in their summaries, and step by step those of Theorems 2 and~3, of Theorem~1, of
Proposition~\ref{p:nonneg}, Proposition~\ref{p:factor}, Lemma~\ref{l:stable}, Proposition~\ref{p:unimod},
Proposition~\ref{p:lift}, Theorems 4 and~5 and Proposition~\ref{p:cone}. In the Lean text I read the statements of
all declarations of the note, the proof of Proposition~\ref{p:factor}, and the definitions of the background
library that the statements rest on (partners, standard data, integral frames, tame characters, split tame
characters, local data, theta integrals on the large side, equivariant lifts, the products of $L$-values,
representedness). The proofs here are shorter than the formal ones: bookkeeping with characters, measures and
matrix identities is merged, and the proofs of Lemma~\ref{l:matched} and Lemma~\ref{l:stable} are written as the
computations that the formal steps amount to. The readers of the dossier mark the long automation steps of
Theorems 2, 3, 5 and of Propositions \ref{p:nonneg} and~\ref{p:cone} as coarse; these proofs are given here in
outline with every input named, and they have not been checked independently of the dossier. Nothing was compiled.

\subhead{Revision.} This edition answers a referee's report on the three notes of the unit. Two proofs were wrong
and are corrected from the Lean text, which I read for both: in the proof of Proposition~\ref{p:unimod} the Gauss
sum is formed with the shears $S_{\varpi^{-1}\sz}$ and the left inverse is the block $\fA_3'$; in the proof of
Proposition~\ref{p:lift} the invariance of the large side holds for the local character of rank $|\bn|$. The
definition of Weil data was completed, the relations of the Inputs to the literature follow the second reader's
verdicts, the quoted criteria of Input~5 are stated with their hypotheses, and letters were changed to agree with
the handout (the partner on the large side is $\dL$, the lift index is $i$).

\head{Appendix A. Concordance with the Lean artifact}
{\fontsize{8.6}{11.6}\selectfont\renewcommand{\leanfont}{\ttfamily\fontsize{7.7}{10}\selectfont\color{leancol}}
\setlength{\LTleft}{0pt}\setlength{\LTright}{\fill}\setlength{\tabcolsep}{0pt}
\begin{longtable}{@{}p{30mm}p{56mm}p{43mm}r@{}}
\textbf{In this note} & \textbf{Lean declaration} & \textbf{File} & \textbf{Lines}\\\midrule\endhead
\ConcordanceRows
\end{longtable}}

\subhead{Three statements as they are written in the artifact.}
\quoted{existsstdExchangetop}
\quoted{StdExchange}
\quoted{auxLiftinjectiveandhasSuminner}
Legend: in these texts {\leanfont X}, {\leanfont W₁}, {\leanfont W}, {\leanfont V'} are the spaces $\sX$, $\sW_1$,
$\sW$, $\sV'$; {\leanfont D}, {\leanfont D₁} the splittings $\fS$, $\fS_1$; {\leanfont σ'} is $\se'$; {\leanfont v},
{\leanfont v₁}, {\leanfont x}, {\leanfont x'} are the real data $\fv$, $\fv_1$, $\fx$, $\fx'$; {\leanfont 𝔠} is
$\fc$ and {\leanfont 𝔠.χ} is $\chi_{\fc}$; {\leanfont ℓ} is $\cN$; {\leanfont thetaLarge}, {\leanfont stdRowDatum},
{\leanfont stdColDatum}, {\leanfont lValueProdInt} are the theta integral on the large side, the standard data and
$\LP{\fc}{e}$. In the third, {\leanfont D} is the data $D$, {\leanfont auxLift} the adjoint lift and {\leanfont
auxPair} the auxiliary pair.

\head{Notation}
\begingroup\fontsize{8.6}{10.9}\selectfont\raggedright
\setlength{\columnsep}{9mm}\setlength{\parskip}{0pt}
\begin{multicols}{2}
\nt{$E$, $E^{+}$, $\A_E$, $\mathcal{O}_E$, $\hat{\mathcal{O}}_E$}{a CM field, its maximal totally real subfield, its adeles, its integers, its integral finite adeles}
\nt{$\tau$; $\fp$, $\bar{\fp}$, $\mathfrak{q}$, $\mathfrak{u}$; $\varpi_{\fp}$; $\Np\fp$}{an infinite place of $E$; finite places, the conjugate place; a uniformiser at $\fp$; the norm of a place}
\nt{$\jmath$, $\upsilon$, $\mathrm{sc}_E$, $\mathrm{e}_{\A}$}{Weil data: the purely imaginary element, the splitting character (a Hecke character of $E$, and its restriction to $\U(1)(\A)$), the scale idele, the additive character of the adeles}
\nt{$P$, $D$, $m$, $\vartheta$, $i$}{parameters, admissible data, the number of lifts, a tame character of $D$, a lift index ($1\le i\le m$) in §\ref{s:norms}}
\nt{$\omega^{\mathrm{bc}}$, $\omega^{\mathrm{bc}}_{\fc}$, $\omega^{\mathrm{bc}}_{\fl}$; $\xi$}{the base change $\mathsf{x}\mapsto\omega(\mathsf{x}/\bar{\mathsf{x}})$ of a character of $\U(1)(\A)$, a Hecke character of $E$; a Hecke character}
\nt{$L^{\Sigma\cup\bar\Sigma}(s,\omega^{\mathrm{bc}})$; $\bar\Sigma$; $s$}{the Hecke $L$-function of the base change without the Euler factors at $\Sigma$ and at the conjugate places $\bar\Sigma$; the set of conjugates; a real number $\ge1$}
\nt{$\ba,\bb,\bm,\bn,\br,\bn',\br'$; $|\bn|$}{finite index sets; the number of elements}
\nt{$\sF$, $\sO$, $\varpi$; $\sR$, $\bar{\sR}$}{a semilocal field, its compact open subring, its topologically nilpotent unit; a commutative ring; the residue ring of a tame character}
\nt{$\sF_{\fc}$, $\sF_\Sigma$}{the product of the completions of $E$ at the places of $\fc$, of $\Sigma$}
\nt{$\Mat{\ba}{\bn}(\sR)$, $\GL_{\bn}(\sR)$; $\sx^{\mathsf T}$, $\sx^{*}$, $\tr$}{matrices with rows $\ba$ and columns $\bn$, invertible matrices; transpose, conjugate transpose, trace}
\nt{$\sg,\sh,\sk,\sx,\sy,\sz,\ssx,\sq$; $\mathsf{a},\mathsf{b},\mathsf{c},\mathsf{d},\mathsf{f},\mathsf{j},\mathsf{l},\mathsf{n},\mathsf{p},\mathsf{r},\mathsf{t},\mathsf{u},\mathsf{v},\mathsf{w},\mathsf{z}$}{matrices, vectors and points of groups: $\sk$ in a general linear or unitary group, $\sz$ an integral matrix or an element of $\sF$, $\ssx$ a left inverse, $\sq$ a Levi coordinate, $\mathsf{v}$ a rational frame, $\mathsf{c},\mathsf{w}$ the coordinates of a radical element, $\mathsf{t}$ signs, $\mathsf{p}$ a point or a family of eigenvalues, $\mathsf{r}$ a coset representative}
\nt{$\mathsf{A},\mathsf{B},\mathsf{I},\mathsf{K},\mathsf{L},\mathsf{M},\mathsf{P},\mathsf{U}$, $\mathsf{J}_0$, $\mathsf{S}_0$, $\mathrm{P}_{+}$}{matrices of §\ref{s:real} and §\ref{s:det}: a local inverse, the form of a split frame, blocks of a mixed model, a matrix of positive coordinates, a frame matrix, a point, a combination of blocks, a skew matrix, the projection to positive coordinates}
\nt{$\se$, $\se'$}{identifications of index sets, $\bb\sqcup\ba\simeq\bn$}
\nt{$\colk{\se}{\sk}$, $\rowk{\se}{\sk}$}{the columns of $\sk$ indexed by $\bb$; the rows of $\sk^{-1}$ indexed by $\ba$}
\nt{$\chi$, $\chi'$, $\inu{\chi}$, $\bar\chi$; $\mathrm{e}_{\sF}$, $\bar{\mathrm{e}}$}{multiplicative characters, the inverse on integral units, the primitive character of the residue ring; an additive character and its reduction}
\nt{$\chi_{\fc}$, $\chi^{(e)}_{\fc}$, $\chi^{\circ}$}{the local character of a split tame character in rank $0$, in rank $e$, and in rank $|\bn|$ (proof of Proposition~\ref{p:lift})}
\nt{$\dL$, $\dS$, $\dS_0$, $\dS'$}{functions on the large side ($\ba\times\bn$ matrices) and on the small side ($\bn\times\bb$ matrices); partners (the formal text writes $\Phi$, $\Psi$)}
\nt{$\stc{\chi}{\ssx}$, $\str{\se}{\chi}{\ssx}$}{the standard column datum and the standard row datum}
\nt{$\cS(\mathcal{U})$, $\cS^{\mathrm{inj}}$, $\cS^{\mathrm{sur}}$}{Schwartz--Bruhat functions supported in a set $\mathcal{U}$, on matrices with a left inverse, on matrices whose transpose has one}
\nt{$\cL_0$, $\mathcal{R}$, $\mathcal{R}_0$}{the functional and the coset representatives of Lemma~\ref{l:reps}; a finite set of them}
\nt{$\localSection{\chi}{\sy}$; $\matchingPairs{\chi}{\se}$}{the local section (Definition~\ref{d:locsec}); the matched pairs (Definition~\ref{d:match})}
\nt{$\cL$, $\cL'$, $\cL_1$; $\mathcal{X}$}{linear maps and functionals; a set}
\nt{$\fA$, $(\fA_1\mid\fA_2\mid\fA_3)$, $\fA_1'$, $\fA_2'$, $\fA_3'$; $\fA^{\mathrm{op}}$; $\mathrm{st}_{\fA}(G)$}{an integral frame, its column blocks, the row blocks of its inverse; the adapted frame of the opposite splitting; the function $G$ read in the frame}
\nt{$\cP$, $\cM$; $S_{\sz}$, $\Gamma_0$, $\mathcal{Z}$}{a stable space of functions, the domain of a lift; the shears of a standard datum, their Gauss sum, a set of integral representatives}
\nt{$G$, $G'$, $G_0$; $\mathrm{n}$, $\mathrm{K}_0$; $i$, $n$}{functions; a counting function, a positive constant; an index or an exponent (in §§\ref{s:local}--\ref{s:coeff}; in §\ref{s:norms} the lift index), a natural number}
\nt{$\sW,\sV,\sV',\sX,\sW_1,\sY,\sZ$; $\sI_{\bn}$}{hermitian spaces; the standard positive definite space}
\nt{$\mathrm{J}_{\sW}$, $\mathrm{x}_{\sW}$, $\negc(\sW)$, $\operatorname{diag}(\sW)$}{the Gram matrix, the adelic frame, the number of negative signs, the diagonal form of the signs}
\nt{$\U(\sW)$, $[\U(\sW)]$, $\operatorname{vol}[\U(\sW)]$, $\mathcal{K}_f$, $d\sk$}{the unitary group, its adelic quotient and the mass of the quotient (normalised in §\ref{s:setting}), the stabilizer of the integral lattice, integration for a Haar measure}
\nt{$\omega$, $\omega'$, $\omega''$, $\delta$; $\chi'$; $\mu$}{characters of $\U(1)(\A)$; a character of $\sF_{\fc}$; weights, integers at the infinite places}
\nt{$\fc$, $\omega_{\fc}$, $\Sigma(\fc)$; $\Sigma$}{a split tame character, its character, its set of places; a set of finite places}
\nt{$\fv,\fv',\fv_1,\fx,\fx',\mathfrak{y}$; $F$, $F'$, $F_{\fc}$, $F_{\sy}$; $f$}{real data; finite data (also an algebraic form in §\ref{s:norms}); a datum}
\nt{$\theta$, $\thI{\sW}{\sV}{\omega}{\omega'}{f}$}{the theta functional; the theta integral over $[\U(\sV)]$}
\nt{$\thL{\sW}{\sV}{\fc}{\fv}{\dL}$, $\thetaSmall{\sW}{\sV}{\fc}{\fv}{\dS}$}{the theta integrals on the large and on the small side}
\nt{$\thetaOnRealData{\sW_1}{\sV}{\omega}{\omega'}{F}{\sx}{\mathsf{p}}$}{the theta functional on real data (Definition~\ref{d:threal})}
\nt{$\LP{\fc}{e}$, $\Pi(\Sigma,\omega,e)$; $e$, $e_0$}{products of partial Hecke $L$-values; a natural number or an integer}
\nt{$\fS$, $\fS_1$; $\cf{\fS}{T}{f}$; $T$, $T'$, $Y$}{splittings; the centre coefficient; matrix indices ($T'=\jmath T$ in the proof of Theorem~2), an adelic matrix}
\nt{$\operatorname{rad}_{\fS}(\mathsf{c},\mathsf{w})$, $\operatorname{lev}_{\fS}(\sq,\sg)$; $\mathrm{sl}_{\fS}(\sx;\fv)$; $\operatorname{vac}_{\fS}(\sV)$}{the radical and Levi elements of the parabolic subgroup of a splitting; the real slice; the vacuum factor of the mixed model}
\nt{$\operatorname{heis}(\sx)$, $\omega_\infty$, $\det_\infty B$}{the Heisenberg operator of a vector; the real Weil representation; the idele that is $\det B$ at the infinite places and $1$ at the finite ones}
\nt{$(\cM,\cD)$}{an equivariant lift}
\nt{$\equalRankThetaRatio{\sV}{\fc}{\fv}$}{the equal-rank ratio (Definition~\ref{d:exch})}
\nt{$\StdExchange{\sX}{\sV}{\sV'}{\mu}{\se}{\fv}{\mathfrak{y}}{\cN}$}{the standard exchange identity (Definition~\ref{d:exch})}
\nt{$\cN$; $\rc$, $\rC$, $\mathrm{R}$, $\mathrm{k}$, $\mathrm{b}$, $\mathrm{z}$; $\mathrm{l}$}{a real function on split tame characters; constants ($\mathrm{R}$ is the constant of unfolding in stages; it is not the right translation $R(\sg)$, the commutative ring $\sR$, or the remainder $R_{d,M}$ of the handout); a complex function on split tame characters}
\nt{$\Omega$, $\Omega_0$, $\Omega_1$; $B$}{weight and bump functions; an invertible adelic $\bb\times\bb$ matrix}
\nt{$\varkappa$, $\mathcal{C}$}{a point of the unit circle; a compact group (proof of Proposition~\ref{p:nonneg})}
\nt{$\posCoordMatrix{\operatorname{sgn}_{\sW}}{\operatorname{sgn}_{\sV}}{F}{\mathsf{u}}$}{the matrix of positive coordinates (Definition~\ref{d:detvec})}
\nt{$\detVectorFn{\operatorname{sgn}_{\sW}}{\operatorname{sgn}_{\sV}}{e}{\mu}{\rc}{F}{\mathsf{u}}$, $\detVectors{e}{\mu}$}{a determinant vector, the set of determinant vectors}
\nt{$\Delta_{\tau,\pm}$}{determinants of matrices of linear forms (proof of Proposition~\ref{p:cone})}
\nt{$\fa$, $\sW(\fa)$, $\mathcal{A}(\fa)$, $\theta_{\fa}$, $\Upsilon(\fa)$, $\fv(\fa)$, $\mathbb{1}_{\fa}$, $\sR(\fa)$, $\mathrm{t}(\fa)$, $\mathrm{i}(\fa)$}{an auxiliary pair and its objects (Setting~\ref{set:aux}); the upright $\mathrm{t}(\fa)$ is an order of a finite group, not the step $t(D)$ of the handout}
\nt{$\gramFactor{\fa}$, $\rallisTerm{D}{\vartheta}{i}{\sV}{y'}{y}{\bar{\sg}}$}{the Gram factor and the terms of the inner product (Definition~\ref{d:gram})}
\nt{$\fl$, $\omega_{\fl}$, $\Sigma(\fl)$, $\fa(\fl)$, $\Xi(\fl)$}{$L$-value data, their character, the set of their places, their auxiliary pair, the lift of the trivial character}
\nt{$\operatorname{hook}(\fl)$, $\operatorname{weyl}(\fl)$, $\operatorname{sg}(\fp)$}{the hook number, the number of elements of the Weyl index set, the sign of a finite place (Setting~\ref{set:aux})}
\nt{$A$, $A_0$; $\rc_F$, $\mathrm{ct}(\bar{\sg})$}{the constants of the inner product formula; the constant of a sharp datum, the coset term (Input~6)}
\nt{$y$, $y'$, $\Xi$; $\pet{y}{y'}$; $R(\sg)$; $l$}{forms; the Petersson product; right translation; a parameter of a discrete series type}
\nt{$\sW_i(D)$, $\sY_0(D)$, $\fa_D(\vartheta,i)$, $\theta^{*}_{D,\vartheta,i}$, $\Theta^{\circ}_{D,\vartheta}$, $\mathfrak{m}_D(\vartheta,e)$}{the objects attached to the data (Setting~\ref{set:aux}); the modulus of the chosen places to the power $e$ times the places of $\vartheta$}

\end{multicols}
\endgroup

\end{document}
